\PassOptionsToPackage{table}{xcolor}
\documentclass{article}
\usepackage{preprint}
\usepackage{newtxtext,newtxmath}
\usepackage{microtype}
\usepackage{graphicx}
\usepackage{booktabs}
\usepackage{multirow}   % header labels that span both rows of a two-row header, so they centre vertically
\usepackage{float}
\usepackage{placeins}
\usepackage{longtable}
\usepackage{enumitem}
\usepackage{amsmath}
\let\Bbbk\relax
\usepackage{amssymb}
\usepackage{xcolor}
\usepackage{url}
\usepackage{hyperref}
\hypersetup{hidelinks}
\input{math_commands}
\input{tables/cf_colors}
\input{tables/cf_numbers}
\setlength{\textfloatsep}{6pt plus 2pt minus 2pt}
\setlength{\abovecaptionskip}{5pt}
\setlength{\belowcaptionskip}{0pt}

\title{Readability Does Not Confer Control over Clinical Concepts}

\authorline{%
Hanqi Jiang\aff{1,2,*} \quad
Weihang You\aff{1,*} \quad
Qingchan Zhu\aff{1} \quad
Junhao Chen\aff{1} \quad
Yi Pan\aff{1}\\
Zhengliang Liu\aff{1} \quad
Wei Liu\aff{3} \quad
Lei Xing\aff{2,\dag}}

\affiliations{%
\aff{1}School of Computing, University of Georgia\\
\aff{2}School of Medicine, Stanford University\\
\aff{3}Department of Radiation Oncology, Mayo Clinic, Phoenix, AZ, USA}

\titlenotes{%
\aff{*}Equal contribution.\\
\aff{\dag}Corresponding author: \texttt{lei@stanford.edu}.}

\begin{document}
\maketitle

\begin{abstract}
A linear probe reads a clinical finding from the visual representation that a vision-language model's language model consumes. Activation steering suggests the converse: writing the probe direction back into that representation should move the model's answer about the finding, making the direction a handle on the concept. Evaluations score whether models encode findings and answer questions about them, but not whether such a handle exists. We test it with \emph{ownership}: at a fixed dose on the same patients, does a concept's own direction move its own answer more than the directions fitted for every competing finding? Across \cfCheckpoints{} checkpoints on NIH ChestX-ray14, CheXpert Plus and COCO, clinical concepts are readable in \cfChestReadablePct{}\% of chest checkpoint--concept cells and answerable in \cfChestAnswerablePct{}\%, but owned in \cfChestOwnedPct{}\%. Two failures combine: most clinical writes do not move their own answer beyond random directions of the same dose, and the same directions move the other findings' answers. The procedure owns \cfCocoOwnedPct{}\% of natural-image object cells, and the gap persists among cells the model answers well (\cfStratChestTopPct{}\% against \cfStratNatTopPct{}\% at a clean-answer AUROC of at least \cfStratTopEdge{}). Non-clinical attributes of the same radiographs are owned no more often than findings, and checkpoints that share a vision tower, and therefore write identical vectors, own different concepts. The deficit follows the chest-radiograph representation and the model that reads it rather than the concept's clinical status. A direction fitted to the model's own answers provides the handle that the label direction lacks. A probe's readability does not confer control through its direction.
\end{abstract}

\section{Introduction}
\label{sec:introduction}

Linear probes are the standard tool for testing what a vision-language model (VLM) represents. They read out a concept by fitting a classifier to activations at a chosen layer \citep{alain2017understanding,belinkov2022probing}. In medical imaging, probes often succeed at linearly decoding clinical findings such as pleural effusion or cardiomegaly from the visual representation of general-purpose and medical VLMs \citep{zhu2026lost,pedersen2026probes}. Activation steering promises the converse: writing the probe direction back into the representation to change the model's behaviour \citep{turner2023steering,panickssery2024caa}. Together, these methods promise to make the probe direction a handle on the concept: read it to detect, write it to control. We ask whether that handle exists.

A steering write that changes the answer does not settle the question. In Qwen2.5-VL-7B on NIH ChestX-ray14, a probe reads effusion from the consumed visual tokens, and writing its direction back raises the model's Effusion answer by $\cfSeedWEff$, more than the 95th percentile of 119 random directions and a coordinate-permutation sham written at the same dose. By the usual steering criteria, the write works. Writing the Nodule direction into the same patients at the same dose raises the Effusion answer further, by $\cfSeedWNodule$ (Figure~\ref{fig:framework}c), although Nodule is the competing direction least similar to Effusion (Appendix~\ref{app:cf-robustness}). Random and sham directions carry no label, so they cannot expose a direction whose effect is shared with other findings \citep{pahde2025navigating}. Exposing it requires writing every competing clinical direction at the same locus and dose in the same patients and asking whether the concept's own direction moves the concept's own answer more than any other. We call the outcome \emph{ownership}. A probe direction that is readable but not owned reads its concept without writing it: reading is not writing.

We turn this comparison into an evaluation (Section~\ref{sec:framework}, Figure~\ref{fig:framework}). At the final visual block whose output the language model consumes, each concept receives three grades. It is \emph{readable} if its probe beats random-label probes of equal capacity, \emph{answerable} if the model's clean yes/no answer beats chance, and \emph{owned} if a fixed-dose write along its direction clears the random and sham references and moves its own answer more than every competing clinical direction under simultaneous bounds. We apply the evaluation to \cfCheckpoints{} checkpoints from \cfFamiliesWord{} families (3B--72B; general-purpose and medical) on two chest-radiograph datasets with report-derived labels, NIH ChestX-ray14 and CheXpert Plus, and on COCO objects as a natural-image control (Section~\ref{sec:setup}).

The three grades separate on chest radiographs (Section~\ref{sec:results}, Table~\ref{tab:cf-main}, Figure~\ref{fig:overview}). Clinical concepts are readable in \cfChestReadablePct\% of chest cells and answerable in \cfChestAnswerablePct\%, but owned in \cfChestOwnedPct\%, and the seed cell reproduces on new patients. The deficit has two parts. Most clinical writes do not move their own answer beyond the random and sham references, and the effect the writes do have spreads across findings: the written concept's own question receives \cfOwnShareNih\% of the positive answer change on NIH, against \cfOwnShareCoco\% on COCO. The write itself is intact. The same procedure owns \cfCocoOwnedPct\% of COCO cells, six small or fine-grained object categories are owned at the rate of the easy ones, and among cells with a clean-answer AUROC of at least \cfStratTopEdge{}, \cfStratChestTopPct\% of chest cells are owned against \cfStratNatTopPct\% of natural-image cells.

Three further results locate the deficit. Checkpoints that share a vision tower write identical vectors into identical representations, yet own different concepts, so the model that reads the representation decides ownership, not the representation alone. Non-clinical attributes of the same radiographs (projection, sex and age) are owned no more often than findings, so the deficit follows the chest-radiograph representation and its reader rather than the concept's clinical status. The reader responds to a different direction: a direction fitted to the model's own clean answer margin is owned in \cfAnsChestOwnedA{} of \cfAnsChestN{} chest cells where the label direction is owned in \cfAnsChestOwnedLabel{}, and on chest radiographs it is nearly orthogonal to the probe direction.

\begin{figure}[t]
    \centering
    \resizebox{\textwidth}{!}{\includegraphics{figures/fig1_method.pdf}}
    \caption{\textbf{Reading and writing a concept at the consumed block.} (a) The probe reads the concept from the consumed visual tokens; its lifted normal is written back to them at $\alpha=+0.25$. (b) The write matrix $W_{q,d}$ and the three comparisons that decide ownership. (c) The seed cell, Effusion in Qwen2.5-VL-7B on NIH ChestX-ray14, whose verdict is stronger competitor.}
    \label{fig:framework}
\end{figure}

\paragraph{Contributions.} (i) We separate reading a concept from writing it and make the difference measurable. Ownership scores a probe direction against directions fitted the same way for the competing concepts, written at the same locus and dose in the same patients, alongside readability and answerability at the consumed visual block (Section~\ref{sec:framework}). (ii) We measure the three grades across \cfCheckpoints{} checkpoints, \cfFamiliesWord{} families, three datasets and \cfCells{} concept cells. Clinical concepts are read and answered far more often than they are owned, while the chest-to-natural-image ownership gap persists where both models answer well (Sections~\ref{sec:headline}--\ref{sec:coco}). (iii) We locate the deficit in how the model reads the chest-radiograph representation: identical writes are owned differently by different language models, non-clinical attributes of the radiograph share the deficit, and the model's own answer direction provides the handle that the label direction lacks (Sections~\ref{sec:reader}--\ref{sec:ansdir}). We release the evaluation, scorecards and write matrices; a new model needs one adapter naming its consumed block.
\section{Background and Related Work}
\label{sec:related}

\paragraph{Probing and its controls.}
Linear classifiers have long been used to track information across network depth \citep{alain2017understanding}. A probe score reflects both the representation and the learner fitted to it. This motivates matched-capacity probes and control tasks \citep{hewitt2019designing}. The probing literature distinguishes information a readout can recover from information the model uses \citep{belinkov2022probing,ravichander2021probing}. Causal erasure pursues the complementary test \citep{ravfogel2020null,elazar2021amnesic}, and a closed-form projection removes a concept from every linear reader at once \citep{belrose2023leace}. Erasure asks what breaks when a concept is deleted; the question here is whether the concept's own direction is the one that moves its answer. Our readability grade adopts the control-task idea: it measures selectivity against random-label probes of equal capacity fitted within recurring image types. Depth-wise probing of medical multimodal models has mapped classification information across towers, connectors, and language layers \citep{zhu2026lost}. Layer-wise probes of MedCLIP expose acquisition and device shortcuts \citep{pedersen2026probes}. Both map what can be read; the write test here intervenes on the decoded clinical direction at the locus where it is read.

\paragraph{Concept directions and what they inherit.}
A concept activation vector is the normal of a classifier separating concept examples from a negative set, and it grades a class by the directional derivative of the logit along that normal \citep{kim2018tcav}. We keep the direction and change the measurement: a bounded forward write at a fixed fraction of each token's norm, scored against directions built the same way for the competing concepts, rather than a gradient sensitivity read at the output. Such a normal shifts with the layer it is fitted at, entangles with correlated concepts, depends on the negative set it is fitted against, and carries where in the image the concept was read \citep{nicolson2025explaining}. Section~\ref{sec:robust-geometry} tests each: alternative constructions, a second locus, and writes concentrated on the tokens the probe reads.

\paragraph{Steering with probe normals.}
Representation engineering makes the step from reading a concept to writing it back \citep{zou2023representation}. Activation engineering changes intermediate representations at inference time and measures the resulting output changes \citep{turner2023steering,panickssery2024caa}. A probe normal is an appealing intervention because it supplies a label and a signed vector from the same representation. Separability-trained concept vectors can absorb distractor signals and diverge from the concept they are meant to model \citep{pahde2025navigating}. A classifier's weight vector is a filter, not the pattern it detects \citep{haufe2014interpretation}. Mass-mean directions steer truth judgements in language models more effectively than logistic normals \citep{marks2024geometry}. The features that identify a concept in the input are largely not the features that control the output, and selecting for output effect improves steering \citep{arad2025saes}. Benchmarks accordingly score concept detection and steering as separate axes of the same method \citep{wu2025axbench}. Random-direction and sham controls test generic sensitivity and coordinate dependence. They do not test whether an equally constructed direction with a different label would move the same answer more. Ownership adds this comparison and evaluates it in the forward path at the consumed locus, using each token's norm as the dose scale. Alignment search certifies a subspace as causal by showing that it implements an interpretable causal variable \citep{geiger2024alignments}; ownership asks a narrower question of a labelled direction: whether the label names the direction that moves its answer. Prior work replaces the separability-trained vector with a better estimator of the same concept and evaluates the replacement against the concept's own label \citep{pahde2025navigating,haufe2014interpretation,marks2024geometry}. Ownership instead fixes the estimator and varies the label: the reference family contains a direction fitted the same way for every other concept in the set, so a construction that improves estimation is scored by whether it wins that comparison. Section~\ref{sec:robust-geometry} runs five constructions and both lifts through the same test.

\paragraph{Writing into a vision-language model.}
Image features are linearly decodable in a vision-language model's hidden layers, and targeted edits to them change downstream answers \citep{rajaram2025lineofsight}. Steering one sparse-autoencoder feature of the vision tower on an empty image and asking the language model what it sees turns the language model into a read-out of what that feature encodes \citep{ferrando2026explain}. That work steers a feature to recover its meaning; this paper starts from the meaning, fits a direction to a clinical label, and asks whether the write keeps that label's selectivity when the other findings' directions are written at the same locus, dose, and patients. Causal probing of multimodal models intervenes on the language model's residual stream with difference-in-means vectors over paired images, and reads out where entity and abstract concepts sit across depth \citep{deng2026causal}. It intervenes inside the reader and scores each concept against its own effect; the write here acts on the representation the reader consumes and scores each concept against the competing concepts' directions, which is the comparison that decides whether the label names the handle.
\section{Reading, Answering, and Writing at a Consumed Locus}
\label{sec:framework}

We ask three questions about a concept $c$ in a checkpoint $M$ on a dataset with binary labels: whether $c$ is readable from the visual representation the language model consumes, whether $M$ answers the question about $c$, and whether writing $c$'s own direction moves $c$'s answer more than any alternative. A registered protocol fixes an estimator and a decision rule for each question (Appendix~\ref{app:cf-labels}).

\subsection{Consumed locus and matched read--write positions}
\label{sec:locus}

Every VLM in the grid passes an image through a vision tower and hands a specific tensor to a connector or merger that produces the language model's inputs. The \emph{primary locus} is the output of the last visual block whose output that consumer reads, restricted to the token positions the consumer actually uses (CLS or padding rows are excluded where the consumer drops them); the \emph{connector locus} is the consumer's output. Appendix~\ref{app:repro-loci} identifies the block for every family. For image $i$ and consumed token $t$ let $\hvec_{it}\in\R^D$ be the block output; mean pooling over the consumed tokens gives the read vector $\hvec_i$. The write acts on the same tokens,
\begin{equation}
\hvec'_{it}=\hvec_{it}+\alpha\,\lVert\hvec_{it}\rVert_2\,\hat\wvec_c,
\label{eq:intervention}
\end{equation}
so that the dose $\alpha$ is a fraction of each token's own norm and is comparable across families whose activation scales differ by orders of magnitude. Reading and writing therefore use the same representation and the same positions. The connector and the language model, everything downstream of the primary locus, form the \emph{reader} of that representation.

\subsection{Directions, controls, and the reference family}
\label{sec:directions}

A fixed seed-0 Gaussian matrix $R\in\R^{D\times512}$ projects every checkpoint to a common readout dimension; featurewise scale $\boldsymbol s$ is fitted on training rows only; logistic regression ($C=1$, 2{,}000 iterations) is fitted per concept with three seeds. The write direction lifts the positive-class coefficient $\boldsymbol\beta_c$ back to the consumed space,
\begin{equation}
\hat\wvec_c=\frac{R\left(\boldsymbol\beta_c\oslash\max(\boldsymbol s,10^{-8})\right)}{\left\lVert R\left(\boldsymbol\beta_c\oslash\max(\boldsymbol s,10^{-8})\right)\right\rVert_2}.
\label{eq:lift}
\end{equation}
Two families of controls bracket the concept direction. For reading, 20 control probes are fitted to random labels drawn within recurring image types (view, sex, and age decade on chest sets; aspect and area buckets on COCO) so that they have the same capacity and the same nuisance structure as the concept probe; selectivity is
\begin{equation}
S=\operatorname{AUROC}(y,r)-\tfrac{1}{20}\textstyle\sum_{k}\operatorname{AUROC}(y^{\mathrm{ctrl}}_k,r^{\mathrm{ctrl}}_k).
\label{eq:selectivity}
\end{equation}
For writing, the \emph{reference family} contains 119 seed-0 random unit directions, a coordinate-permutation sham of $\hat\wvec_c$ (same coordinates, scrambled arrangement), and the five directions of the other clinical concepts, all written at the same dose to the same images. Cosine and projection depend on the inner product chosen for the representation \citep{park2024linear}, so each cosine names its metric: cosines between directions are reported in the raw model space, and the two comparisons that carry the argument, the answer direction against the probe normal and the expert-label refit against the report-label direction, are also reported under the training covariance.

\subsection{Ownership}
\label{sec:ownership}

For question $q$ and direction $d$, $W_{q,d}$ is the paired mean change in $P(\mathrm{yes})$ of $q$ under a write along $d$, where $P(\mathrm{yes})=\sigma(\ell_{\mathrm{yes}}-\ell_{\mathrm{no}})$ at the first answer position and each log-probability is the maximum over case and leading-space variants that tokenise to one token. The $6\times6$ matrix $W$ over concept questions and concept directions is the object of interest (Figure~\ref{fig:write-matrices}). Ownership is the diagonal margin
\begin{equation}
O_q=W_{q,q}-\max_{d\neq q}W_{q,d},
\label{eq:margin}
\end{equation}
the advantage of the concept's own direction over its strongest clinical competitor on the concept's own question. Ownership is a comparison at one fixed dose, on the same patients, against directions built the same way for the other concepts. $W_{q,q}$ measures how far the concept's own direction moves its answer; $O_q$ measures whether it moves that answer further than every such competitor. A large $W_{q,q}$ shows that $d=q$ can write the answer; only $O_q>0$ shows that the label names the effective direction. A concept can therefore fail ownership in two ways: its own write stays within the random and sham references, or the write clears them without holding its advantage over every competitor. Section~\ref{sec:headline} measures both. The seed cell illustrates the latter: the Effusion direction raises the Effusion answer by $\cfSeedWEff$, above the random 95th percentile and the sham, and the Nodule direction raises the same answer by $\cfSeedWNodule$, so $O_{\mathrm{Effusion}}=\cfSeedEffO$ (Section~\ref{sec:seed}). Uncertainty over patients is quantified by 2{,}000 bootstrap draws of test rows with the maximum recomputed inside every draw, and the 30 contrasts $W_{q,q}-W_{q,d}$ receive simultaneous 95\% lower bounds from a centred studentised max-$T$ bootstrap (Appendix~\ref{app:simultaneous-bounds}).

\subsection{Graded outcomes}
\label{sec:grades}

The readable and answerable grades are computed on the 400 calibration rows, which carry no write. The write matrix and every ownership quantity are computed on the 600 test rows, which share no patient with the calibration rows.
\begin{itemize}[leftmargin=*,itemsep=1pt,topsep=2pt]
\item \textbf{Readable:} the one-sided 95\% lower bound of $S$ on the calibration rows is above zero, with at least ten known positives and ten known negatives among those 400 rows.
\item \textbf{Answerable:} the one-sided 95\% lower bound of the clean-answer AUROC on the yes/no question, on the same calibration rows, is above $0.5$, under the same ten-and-ten support rule.
\item \textbf{Owned:} on the test rows, the write meets the \emph{steering reference} ($W_{q,q}>0$, above the 95th percentile of the 119 random writes, and above the absolute sham) and every simultaneous lower bound of $W_{q,q}-W_{q,d}$ is positive, the \emph{advantage} verdict. A cell with any simultaneous upper bound below zero has a \emph{stronger competitor}; the rest are \emph{unresolved}.
\end{itemize}
A block contributes to the answerable and owned counts only when its write matrix and its calibration pass both completed. A probe-graded block whose write matrix is ineligible or incomplete contributes readability alone.
The grades are defined for one direction construction, one relative dose per module with a five-point sweep, six templates, pooled consumed-token features, and one primary and one connector locus; Section~\ref{sec:robust-geometry} scores four alternative direction constructions by the same decision rule, against the write matrix's own random family and sham. Table~\ref{tab:cf-ledger} lists the rows, reference, decision rule and sham of every graded analysis.
Every block additionally reports the signed dose range of $O_q$ over $\alpha\in\{-0.5,\ldots,+0.5\}$, the standard deviation of $O_q$ over the three fit seeds (seeds 1 and 2 resample the training units with replacement and refit the scaler and the six probes), ownership and the absolute write effects at the connector locus, the label gap (under the concept's own write, the mean shift of the answer margin $\ell_{\mathrm{yes}}-\ell_{\mathrm{no}}$ on label-positive rows minus that on label-negative rows, in logits, median over concepts), and the wording (\emph{is} versus \emph{show}) and mapping (A-present versus B-present) contrasts of $O_q$ across six question templates.

\subsection{Acceptance gates}
\label{sec:gates}

Seven checks run on 16 held-out preflight rows at both loci before block scoring (Appendix~\ref{app:cf-gates}): bitwise determinism; identity at $\alpha=0$; reach; batched-versus-single agreement of the answer logits; image-free semantic mapping of every template; throughput; and fp32 recomputation of the answer logits. To pass reach, the write must change the consumer's input and answer logits without changing anything outside the consumed tokens. We declare batched-versus-single deviations above $0.25$ and neutralise them by scoring every condition of a question in one fixed batch composition. We mark templates that fail semantic mapping as ineligible and close them as such. If a checkpoint's yes/no template fails, we grade it using its eligible B-positive template (IB for Llama~3.2 Vision). The checkpoint's scorecard records the template used. The block's scorecard records every declared deviation.
\section{Models, Datasets, and Scale}
\label{sec:setup}

\paragraph{Checkpoints.}
We evaluate \cfCheckpoints{} instruction-tuned checkpoints across \cfFamiliesWord{} families (Table~\ref{tab:cf-models}): Qwen2.5-VL 3B/7B/32B/72B \citep{bai2025qwen25vl}, Qwen3-VL 4B/8B/32B \citep{qwen2025qwen3vl}, InternVL3.5 8B/14B/38B \citep{wang2025internvl35}, Gemma~3 4B/12B/27B \citep{gemma2025gemma3}, MedGemma 4B/27B \citep{sellergren2025medgemma}, Lingshu 7B/32B \citep{lasa2025lingshu}, LLaVA-1.5 7B/13B \citep{liu2024improved}, the released LLaVA-Med~1.5 7B checkpoint \citep{li2023llavamed}, Llama~3.2~Vision 11B \citep{meta2024llama32vision}, CheXagent-2-3B and CheXagent-8B \citep{chen2024chexagent}, LLaVA-Rad-7B \citep{zambranochaves2025llavarad}, and HuatuoGPT-Vision-7B \citep{chen2024huatuogptvision}. Of these, \cfMedicalFamiliesWord{} are medical fine-tunes (MedGemma, Lingshu, LLaVA-Med, CheXagent, LLaVA-Rad, HuatuoGPT-Vision). The rest are general-purpose. The checkpoints span \cfTowerDesignsWord{} vision-tower designs and three consumer designs: token insertion after a merger, token insertion after a projector, and cross-attention. The towers include CLIP \citep{radford2021clip} in the LLaVA families, SigLIP \citep{zhai2023siglip} in Gemma~3, and a SigLIP~2 design \citep{tschannen2025siglip2} in Qwen3-VL; MedGemma's MedSigLIP tower is a medically tuned SigLIP released with that model \citep{sellergren2025medgemma}. Three towers are trained on chest radiographs: the XraySigLIP tower of CheXagent-2-3B, the EVA-CLIP tower of CheXagent-8B, and the BiomedCLIP-CXR tower of LLaVA-Rad-7B. HuatuoGPT-Vision-7B tunes on medical data over the same stock CLIP tower as LLaVA-1.5. Consumed-token counts range from \cfTokensMin{} to \cfTokensMax{} per image. We run each checkpoint in bf16 at a pinned revision with its native processor settings and a fixed image resolution per family.

\paragraph{Datasets and concepts.}
NIH ChestX-ray14 \citep{wang2017chestxray} and CheXpert Plus \citep{chambon2024chexpertplus} provide chest radiographs with report-derived labels; the labeler that produces the CheXpert labels, its uncertainty convention, and its radiologist-annotated validation studies come from CheXpert \citep{irvin2019chexpert}. We evaluate six findings per dataset (NIH: Effusion, Atelectasis, Pneumothorax, Cardiomegaly, Mass, Nodule; CheXpert: Effusion, Atelectasis, Pneumothorax, Cardiomegaly, Consolidation, Edema). COCO \citep{lin2014coco} provides the natural-image control, with binary presence labels for six object categories (person, dog, car, chair, bottle, bicycle). Seeded hashes of patient (or image) identity select one fixed cohort per dataset: 20{,}000 training rows (\cfCohortNihTrainRows{} on NIH, every row of that dataset's training split), 400 calibration rows, 600 test rows, and 16 preflight rows. We fit the probes and the answer direction on training rows. We assign readable and answerable grades on calibration rows. We apply every write to test rows. All checkpoints use the same cohorts. Every chest image is a frontal study. On CheXpert every role holds one study per patient. On NIH the calibration, test, and preflight roles hold one study per patient and the training role holds every study of each training patient: \cfCohortNihTrainRows{} studies from \cfCohortNihTrainUnits{} patients, \cfCohortNihTrainMulti{} of whom contribute more than one and one of whom contributes \cfCohortNihTrainMax{}. Each COCO role holds one image. No patient and no image appears in two roles of a dataset (Appendix~\ref{app:repro-cohorts}).

\paragraph{Questions and templates.}
We ask a yes/no question about each concept with six templates that cross two wordings (\emph{is there} versus \emph{does the image show}) with three answer interfaces (yes/no, A-present, B-present). A template is named by its wording and its interface: I or W for \emph{is} or \emph{show}, then Y, A, or B for the yes/no or letter mapping. We use the \emph{is}/yes-no template for the primary grade. The other five provide wording and mapping contrasts. The main write dose is $\alpha=+0.25$. The dose module adds $\alpha\in\{-0.5,-0.25,-0.1,+0.1,+0.5\}$.

\paragraph{Scale.}
Each block pairs one checkpoint with one dataset. It scores the $6\times6$ write matrix, the sham, and 119 random directions on 600 test rows (457{,}200 outcomes). The campaign targets \cfCoveragePlannedWord{} modules per block: the write matrix, the calibration pass, the five additional templates, the dose sweep, the two probe refits, the connector locus, and the connector-locus calibration. \cfCoverageBlocksAllSeven{} of the \cfBlocks{} blocks carry all \cfCoveragePlannedWord{} (Appendix~\ref{app:cf-coverage}). The campaign comprises \cfBlocks{} blocks, \cfCells{} scored concept cells, about 30 million scored answers, and roughly 3{,}000 A100-hours. Every checkpoint has core write and calibration results on all three datasets. We grade Llama~3.2~Vision 11B on its eligible B-positive template (Appendix~\ref{app:campaign}; label handling, registration, and the released artifact are in Appendix~\ref{app:cf-labels}).
\section{Results}
\label{sec:results}

\subsection{Reading is not writing}
\label{sec:headline}

Both chest datasets follow the same pattern (Table~\ref{tab:cf-main}, Figure~\ref{fig:overview}). Table~\ref{tab:cf-main} orders checkpoints by their clinical owned share and shades each cell by magnitude. Clinical concepts are readable in \cfChestReadable{} of \cfChestReadableN{} cells (\cfChestReadablePct\%): median probe selectivity over the matched random-label controls is $\cfSelChestMed$, with quartiles $\cfSelChestQOne$ and $\cfSelChestQThree$. The models also answer: \cfChestAnswerable{} of \cfChestAnswerableN{} cells (\cfChestAnswerablePct\%) have a clean-answer AUROC lower bound above chance. Of these answerable cells, \cfReadAns{} are also readable. Yet only \cfChestOwned{} cells (\cfChestOwnedPct\%) are owned. Readability and answerability do not suffice for writing: only \cfReadAnsOwned{} of the \cfReadAns{} readable and answerable cells are owned (\cfReadAnsOwnedPct\%), compared with \cfRestOwned{} of the remaining \cfRestN{} (\cfRestOwnedPct\%).

Ownership on chest radiographs breaks at the write itself. \cfChestRefBelow{} of \cfChestCellsN{} probe directions fail the random and sham reference for their own question (Table~\ref{tab:cf-contingency}). Among those below-reference cells, \cfChestBelowStrong{} have a stronger clinical competitor; they account for \cfChestBelowStrongPct\% of all \cfChestCompetitor{} stronger-competitor verdicts. Once an own write clears the reference, \cfChestOwned{} of \cfChestRefMet{} cells are owned. Readable and answerable cells follow the same pattern, while on COCO \cfCocoRefMet{} of \cfCocoOwnedN{} writes clear the reference and \cfCocoOwned{} are owned. The chest deficit therefore starts before fine-grained attribution: most probe directions scarcely move the answers they name, while other finding directions often move those same answers more.

\begin{table}[t]
\centering
\scriptsize
\setlength{\tabcolsep}{2pt}
\renewcommand{\arraystretch}{1.08}
\input{tables/cf_colors}
\caption{\textbf{Reading, answering, and writing per checkpoint.} Per dataset: mean probe selectivity (\textsc{Read}), clean-answer AUROC (\textsc{Ans}), and ownership (\textsc{Own}), superscripted with the graded concepts of six, then the owned share of clinical and COCO cells and their ratio. Shading darkens with magnitude: sage for \textsc{Read} and \textsc{Ans}, terracotta for positive and blue for negative \textsc{Own}, terracotta for the owned share. Grey marks an ineligible template (inel.), an incomplete write matrix (inc.); ``--'' is a block not run.}
\label{tab:cf-main}
\begin{tabular}{l@{\hspace{3pt}}ccc@{\hspace{4pt}}ccc@{\hspace{4pt}}ccc@{\hspace{4pt}}ccc}
\toprule
\multirow{2}{*}{Checkpoint} & \multicolumn{3}{c}{NIH ChestX-ray14} & \multicolumn{3}{c}{CheXpert Plus} & \multicolumn{3}{c}{COCO (control)} & \multicolumn{3}{c}{Owned share}\\
\cmidrule(lr){2-4}\cmidrule(lr){5-7}\cmidrule(lr){8-10}\cmidrule(lr){11-13}
& \multicolumn{1}{c}{\textsc{Read}} & \multicolumn{1}{c}{\textsc{Ans}} & \multicolumn{1}{c}{\textsc{Own}} & \multicolumn{1}{c}{\textsc{Read}} & \multicolumn{1}{c}{\textsc{Ans}} & \multicolumn{1}{c}{\textsc{Own}} & \multicolumn{1}{c}{\textsc{Read}} & \multicolumn{1}{c}{\textsc{Ans}} & \multicolumn{1}{c}{\textsc{Own}} & \multicolumn{1}{c}{clin.} & \multicolumn{1}{c}{COCO} & \multicolumn{1}{c}{ratio}\\
\midrule
Lingshu 32B & \cellcolor{cfSage2}0.08$^{3}$ & \cellcolor{cfSage3}0.78$^{6}$ & \cellcolor{cfSlate2}-0.02$^{2}$ & \cellcolor{cfSage4}0.20$^{4}$ & \cellcolor{cfSage3}0.87$^{5}$ & \cellcolor{cfTerra2}+0.04$^{5}$ & \cellcolor{cfSage1}0.03$^{1}$ & \cellcolor{cfSage4}0.95$^{5}$ & \cellcolor{cfTerra4}+0.28$^{6}$ & \cellcolor{cfTerra4}58\% & \cellcolor{cfTerra4}100\% & 58\% \\
CheXagent-2-3B & \cellcolor{cfSage3}0.13$^{5}$ & \cellcolor{cfSage3}0.86$^{6}$ & \cellcolor{cfTerra1}+0.00$^{2}$ & \cellcolor{cfSage4}0.21$^{5}$ & \cellcolor{cfSage4}0.95$^{5}$ & \cellcolor{cfTerra1}+0.01$^{4}$ & \cellcolor{cfSage3}0.13$^{4}$ & \cellcolor{cfSage3}0.81$^{5}$ & \cellcolor{cfTerra2}+0.03$^{2}$ & \cellcolor{cfTerra4}50\% & \cellcolor{cfTerra2}33\% & 150\% \\
Qwen2.5-VL-72B & \cellcolor{cfSage2}0.07$^{2}$ & \cellcolor{cfSage2}0.62$^{4}$ & \cellcolor{cfSlate2}-0.03$^{3}$ & \cellcolor{cfSage3}0.20$^{4}$ & \cellcolor{cfSage2}0.71$^{4}$ & \cellcolor{cfSlate2}-0.09$^{1}$ & \cellcolor{cfSage1}0.01$^{1}$ & \cellcolor{cfSage4}0.99$^{5}$ & \cellcolor{cfTerra4}+0.66$^{6}$ & \cellcolor{cfTerra3}33\% & \cellcolor{cfTerra4}100\% & 33\% \\
Lingshu 7B & \cellcolor{cfSage2}0.08$^{3}$ & \cellcolor{cfSage3}0.76$^{6}$ & \cellcolor{cfTerra2}+0.02$^{1}$ & \cellcolor{cfSage4}0.21$^{4}$ & \cellcolor{cfSage3}0.86$^{5}$ & \cellcolor{cfTerra2}+0.04$^{3}$ & \cellcolor{cfSage1}0.02$^{1}$ & \cellcolor{cfSage4}0.97$^{5}$ & \cellcolor{cfTerra4}+0.68$^{6}$ & \cellcolor{cfTerra3}33\% & \cellcolor{cfTerra4}100\% & 33\% \\
Qwen3-VL-4B & \cellcolor{cfSage2}0.08$^{3}$ & \cellcolor{cfSage2}0.71$^{6}$ & \cellcolor{cfSlate3}-0.11$^{0}$ & \cellcolor{cfSage4}0.20$^{4}$ & \cellcolor{cfSage2}0.74$^{4}$ & \cellcolor{cfSlate1}-0.01$^{2}$ & \cellcolor{cfSage1}0.04$^{4}$ & \cellcolor{cfSage4}0.99$^{5}$ & \cellcolor{cfTerra4}+0.34$^{6}$ & \cellcolor{cfTerra3}17\% & \cellcolor{cfTerra4}100\% & 17\% \\
Qwen3-VL-32B & \cellcolor{cfSage2}0.10$^{2}$ & \cellcolor{cfSage2}0.70$^{6}$ & \cellcolor{cfSlate2}-0.02$^{1}$ & \cellcolor{cfSage3}0.17$^{5}$ & \cellcolor{cfSage3}0.82$^{5}$ & \cellcolor{cfSlate2}-0.02$^{1}$ & \cellcolor{cfSage2}0.05$^{4}$ & \cellcolor{cfSage4}0.99$^{5}$ & \cellcolor{cfTerra3}+0.15$^{6}$ & \cellcolor{cfTerra3}17\% & \cellcolor{cfTerra4}100\% & 17\% \\
InternVL3.5-38B & \cellcolor{cfSage2}0.10$^{5}$ & \cellcolor{cfSage3}0.78$^{6}$ & \cellcolor{cfSlate1}-0.00$^{0}$ & \cellcolor{cfSage3}0.17$^{4}$ & \cellcolor{cfSage3}0.86$^{5}$ & \cellcolor{cfSlate1}-0.00$^{2}$ & \cellcolor{cfSage2}0.08$^{4}$ & \cellcolor{cfSage4}0.99$^{5}$ & \cellcolor{cfTerra1}+0.01$^{6}$ & \cellcolor{cfTerra3}17\% & \cellcolor{cfTerra4}100\% & 17\% \\
Gemma 3 12B & \cellcolor{cfSage2}0.07$^{2}$ & \cellcolor{cfSage2}0.62$^{5}$ & \cellcolor{cfSlate4}-0.37$^{1}$ & \cellcolor{cfSage3}0.18$^{4}$ & \cellcolor{cfSage1}0.52$^{0}$ & \cellcolor{cfSlate4}-0.38$^{1}$ & \cellcolor{cfSage2}0.07$^{5}$ & \cellcolor{cfSage4}0.95$^{5}$ & \cellcolor{cfTerra4}+0.44$^{6}$ & \cellcolor{cfTerra3}17\% & \cellcolor{cfTerra4}100\% & 17\% \\
MedGemma 4B & \cellcolor{cfSage3}0.12$^{3}$ & \cellcolor{cfSage3}0.80$^{6}$ & \cellcolor{cfSlate3}-0.13$^{0}$ & \cellcolor{cfSage4}0.21$^{4}$ & \cellcolor{cfSage3}0.87$^{5}$ & \cellcolor{cfSlate3}-0.16$^{2}$ & \cellcolor{cfSage2}0.06$^{4}$ & \cellcolor{cfSage4}0.94$^{5}$ & \cellcolor{cfTerra3}+0.23$^{6}$ & \cellcolor{cfTerra3}17\% & \cellcolor{cfTerra4}100\% & 17\% \\
InternVL3.5-14B & \cellcolor{cfSage2}0.07$^{4}$ & \cellcolor{cfSage2}0.73$^{6}$ & \cellcolor{cfSlate1}-0.00$^{1}$ & \cellcolor{cfSage3}0.16$^{4}$ & \cellcolor{cfSage3}0.87$^{5}$ & \cellcolor{cfTerra1}+0.00$^{1}$ & \cellcolor{cfSage3}0.12$^{4}$ & \cellcolor{cfSage4}0.98$^{5}$ & \cellcolor{cfTerra1}+0.02$^{5}$ & \cellcolor{cfTerra3}17\% & \cellcolor{cfTerra3}83\% & 20\% \\
Qwen2.5-VL-32B & \cellcolor{cfSage2}0.09$^{3}$ & \cellcolor{cfSage2}0.63$^{4}$ & \cellcolor{cfSlate2}-0.06$^{1}$ & \cellcolor{cfSage3}0.16$^{4}$ & \cellcolor{cfSage2}0.70$^{4}$ & \cellcolor{cfSlate3}-0.11$^{0}$ & \cellcolor{cfSage1}0.02$^{1}$ & \cellcolor{cfSage4}0.99$^{5}$ & \cellcolor{cfTerra4}+0.38$^{6}$ & \cellcolor{cfTerra2}8\% & \cellcolor{cfTerra4}100\% & 8\% \\
InternVL3.5-8B & \cellcolor{cfSage2}0.06$^{2}$ & \cellcolor{cfSage2}0.73$^{6}$ & \cellcolor{cfSlate1}-0.01$^{0}$ & \cellcolor{cfSage3}0.19$^{4}$ & \cellcolor{cfSage3}0.86$^{5}$ & \cellcolor{cfSlate1}-0.00$^{1}$ & \cellcolor{cfSage2}0.12$^{5}$ & \cellcolor{cfSage4}0.98$^{5}$ & \cellcolor{cfTerra2}+0.03$^{6}$ & \cellcolor{cfTerra2}8\% & \cellcolor{cfTerra4}100\% & 8\% \\
HuatuoGPT-Vision-7B & \cellcolor{cfSage2}0.08$^{3}$ & \cellcolor{cfSage2}0.71$^{6}$ & \cellcolor{cfSlate2}-0.03$^{1}$ & \cellcolor{cfSage3}0.19$^{4}$ & \cellcolor{cfSage3}0.78$^{4}$ & \cellcolor{cfSlate1}-0.01$^{0}$ & \cellcolor{cfSage2}0.07$^{4}$ & \cellcolor{cfSage4}0.99$^{5}$ & \cellcolor{cfTerra1}+0.01$^{4}$ & \cellcolor{cfTerra2}8\% & \cellcolor{cfTerra3}67\% & 12\% \\
LLaVA-Med 1.5 7B & \cellcolor{cfSage2}0.08$^{3}$ & \cellcolor{cfSage1}0.53$^{0}$ & \cellcolor{cfSlate1}-0.01$^{0}$ & \cellcolor{cfSage3}0.19$^{4}$ & \cellcolor{cfSage1}0.59$^{2}$ & \cellcolor{cfSlate1}-0.01$^{1}$ & \cellcolor{cfSage2}0.07$^{4}$ & \cellcolor{cfSage3}0.76$^{5}$ & \cellcolor{cfTerra1}+0.00$^{1}$ & \cellcolor{cfTerra2}8\% & \cellcolor{cfTerra2}17\% & 50\% \\
CheXagent-8B & \cellcolor{cfSage3}0.16$^{5}$ & \cellcolor{cfSage2}0.73$^{6}$ & \cellcolor{cfSlate1}-0.00$^{0}$ & \cellcolor{cfSage4}0.23$^{5}$ & \cellcolor{cfSage3}0.88$^{5}$ & \cellcolor{cfSlate1}-0.02$^{1}$ & \cellcolor{cfSage2}0.05$^{3}$ & \cellcolor{cfSage1}0.52$^{1}$ & \cellcolor{cfSlate2}-0.03$^{0}$ & \cellcolor{cfTerra2}8\% & \cellcolor{cfTerra1}0\% & -- \\
Qwen2.5-VL-3B & \cellcolor{cfSage2}0.07$^{2}$ & \cellcolor{cfSage1}0.58$^{2}$ & \cellcolor{cfSlate2}-0.08$^{0}$ & \cellcolor{cfSage3}0.17$^{4}$ & \cellcolor{cfSage1}0.57$^{2}$ & \cellcolor{cfSlate2}-0.07$^{0}$ & \cellcolor{cfSage1}0.05$^{4}$ & \cellcolor{cfSage4}0.97$^{5}$ & \cellcolor{cfTerra4}+0.52$^{6}$ & \cellcolor{cfTerra1}0\% & \cellcolor{cfTerra4}100\% & 0\% \\
Qwen2.5-VL-7B & \cellcolor{cfSage2}0.08$^{3}$ & \cellcolor{cfSage2}0.62$^{3}$ & \cellcolor{cfSlate3}-0.14$^{0}$ & \cellcolor{cfSage3}0.20$^{4}$ & \cellcolor{cfSage2}0.65$^{2}$ & \cellcolor{cfSlate2}-0.08$^{0}$ & \cellcolor{cfSage1}0.02$^{1}$ & \cellcolor{cfSage4}0.98$^{5}$ & \cellcolor{cfTerra4}+0.67$^{6}$ & \cellcolor{cfTerra1}0\% & \cellcolor{cfTerra4}100\% & 0\% \\
Qwen3-VL-8B & \cellcolor{cfSage2}0.10$^{4}$ & \cellcolor{cfSage2}0.69$^{6}$ & \cellcolor{cfSlate1}-0.01$^{0}$ & \cellcolor{cfSage3}0.17$^{4}$ & \cellcolor{cfSage2}0.75$^{4}$ & \cellcolor{cfSlate2}-0.04$^{0}$ & \cellcolor{cfSage1}0.04$^{3}$ & \cellcolor{cfSage4}0.99$^{5}$ & \cellcolor{cfTerra4}+0.58$^{6}$ & \cellcolor{cfTerra1}0\% & \cellcolor{cfTerra4}100\% & 0\% \\
Gemma 3 27B & \cellcolor{cfSage2}0.07$^{2}$ & \cellcolor{cfSage2}0.60$^{3}$ & \cellcolor{cfSlate2}-0.03$^{0}$ & \cellcolor{cfSage3}0.18$^{4}$ & \cellcolor{cfSage1}0.55$^{2}$ & \cellcolor{cfSlate1}-0.01$^{0}$ & \cellcolor{cfSage2}0.07$^{5}$ & \cellcolor{cfSage4}0.96$^{5}$ & \cellcolor{cfTerra4}+0.31$^{6}$ & \cellcolor{cfTerra1}0\% & \cellcolor{cfTerra4}100\% & 0\% \\
Gemma 3 4B & \cellcolor{cfSage2}0.07$^{2}$ & \cellcolor{cfSage1}0.56$^{2}$ & \cellcolor{cfSlate3}-0.11$^{0}$ & \cellcolor{cfSage3}0.18$^{4}$ & \cellcolor{cfSage1}0.53$^{1}$ & \cellcolor{cfSlate4}-0.27$^{0}$ & \cellcolor{cfSage2}0.07$^{5}$ & \cellcolor{cfSage4}0.94$^{5}$ & \cellcolor{cfTerra3}+0.24$^{5}$ & \cellcolor{cfTerra1}0\% & \cellcolor{cfTerra3}83\% & 0\% \\
LLaVA-1.5-7B & \cellcolor{cfSage2}0.08$^{2}$ & \cellcolor{cfSage1}0.51$^{0}$ & \cellcolor{cfSlate2}-0.04$^{0}$ & \cellcolor{cfSage3}0.17$^{4}$ & \cellcolor{cfSage1}0.52$^{0}$ & \cellcolor{cfSlate2}-0.03$^{0}$ & \cellcolor{cfSage4}0.28$^{5}$ & \cellcolor{cfSage4}0.97$^{5}$ & \cellcolor{cfTerra2}+0.03$^{5}$ & \cellcolor{cfTerra1}0\% & \cellcolor{cfTerra3}83\% & 0\% \\
MedGemma 27B & \cellcolor{cfSage3}0.12$^{3}$ & \cellcolor{cfSage3}0.80$^{6}$ & \cellcolor{cfSlate3}-0.14$^{0}$ & \cellcolor{cfSage4}0.21$^{4}$ & \cellcolor{cfSage3}0.87$^{5}$ & \cellcolor{cfSlate2}-0.06$^{0}$ & \cellcolor{cfSage2}0.06$^{4}$ & \cellcolor{cfSage4}0.97$^{5}$ & \cellcolor{cfTerra3}+0.12$^{4}$ & \cellcolor{cfTerra1}0\% & \cellcolor{cfTerra3}67\% & 0\% \\
LLaVA-1.5-13B & \cellcolor{cfSage2}0.08$^{2}$ & \cellcolor{cfSage1}0.53$^{0}$ & \cellcolor{cfSlate2}-0.02$^{0}$ & \cellcolor{cfSage3}0.17$^{4}$ & \cellcolor{cfSage1}0.51$^{1}$ & \cellcolor{cfSlate1}-0.01$^{0}$ & \cellcolor{cfSage4}0.28$^{5}$ & \cellcolor{cfSage4}0.96$^{5}$ & \cellcolor{cfTerra2}+0.02$^{2}$ & \cellcolor{cfTerra1}0\% & \cellcolor{cfTerra2}33\% & 0\% \\
Llama 3.2 Vision 11B & \cellcolor{cfSage2}0.08$^{2}$ & \cellcolor{cfSage2}0.65$^{4}$ & \cellcolor{cfSlate2}-0.03$^{0}$ & \cellcolor{cfSage3}0.16$^{4}$ & \cellcolor{cfSage2}0.73$^{4}$ & \cellcolor{cfSlate2}-0.04$^{0}$ & \cellcolor{cfSage3}0.12$^{5}$ & \cellcolor{cfSage4}0.95$^{5}$ & \cellcolor{cfSlate1}-0.01$^{1}$ & \cellcolor{cfTerra1}0\% & \cellcolor{cfTerra2}17\% & 0\% \\
LLaVA-Rad-7B & \cellcolor{cfSage3}0.15$^{4}$ & \cellcolor{cfSage2}0.71$^{5}$ & \cellcolor{cfSlate1}-0.01$^{0}$ & \cellcolor{cfSage4}0.27$^{5}$ & \cellcolor{cfSage3}0.87$^{5}$ & \cellcolor{cfSlate1}-0.02$^{0}$ & \cellcolor{cfSage3}0.12$^{4}$ & \cellcolor{cfSage1}0.50$^{0}$ & \cellcolor{cfSlate2}-0.03$^{0}$ & \cellcolor{cfTerra1}0\% & \cellcolor{cfTerra1}0\% & -- \\
\midrule
\textbf{All cells} & 0.09 & 0.68 & -0.06 & 0.19 & 0.74 & -0.05 & 0.08 & 0.92 & +0.23 & 13\% & 75\% & 17\% \\
\quad count / cells & 74/150 & 110/150 & 13/150 & 104/150 & 89/150 & 25/150 & 90/150 & 116/150 & 113/150 & & & \\
\bottomrule
\end{tabular}
\end{table}

\begin{table}[t]
\centering
\small
\setlength{\tabcolsep}{4pt}
\caption{\textbf{Steering reference against verdict.} Primary write-matrix cells of the included blocks, cross-classified by the steering reference and by the simultaneous-bound verdict against the five other clinical directions.}
\label{tab:cf-contingency}
\begin{tabular}{llrrrr}
\toprule
Cells & Steering reference & \multicolumn{1}{c}{advantage} & \multicolumn{1}{c}{stronger competitor} & \multicolumn{1}{c}{unresolved} & \multicolumn{1}{c}{total}\\
\midrule
Chest, all cells & met & 38 & 9 & 9 & 56 \\
 & not met & 24 & 187 & 33 & 244 \\
Chest, readable and answerable & met & 28 & 5 & 5 & 38 \\
 & not met & 13 & 74 & 16 & 103 \\
COCO, all cells & met & 113 & 0 & 4 & 117 \\
 & not met & 10 & 17 & 6 & 33 \\
\bottomrule
\end{tabular}
\end{table}

\begin{figure}[t]
    \centering
    \resizebox{\textwidth}{!}{\includegraphics{figures/fig2_overview.pdf}}
    \caption{\textbf{Reading, answering, and writing across the grid.} (a) Readable, answerable, and owned cells (of 150) per dataset, with descriptive 95\% bootstrap intervals over cells. (b) Mean readability against mean ownership per checkpoint. (c) Median ownership and its interquartile band over 25 blocks against the relative dose; the dotted line marks $\alpha=+0.25$, the dose at which every grade is scored, and the COCO band is clipped at $0.2$. (d) Cumulative share of cells by the rank of the concept write among the 120 compared effects; markers at rank 1 give the share ranked first.}
    \label{fig:overview}
\end{figure}

\subsection{Effusion exposes the ownership gap on new patients}
\label{sec:seed}

The seed finding reproduces on 600 new patients (Table~\ref{tab:cf-seed}). The Effusion write raises the Effusion answer by $\cfSeedWEff$, above the random 95th percentile ($\cfSeedEffRandP$) and the sham ($\cfSeedEffSham$), and ranks second among the 120 compared effects. Yet the Nodule direction raises the Effusion answer by $\cfSeedWNodule$. Effusion ownership is $O_{\mathrm{Effusion}}=\cfSeedEffO$ (95\% CI $[\cfSeedEffOLo,\cfSeedEffOHi]$), matching the $\cfSeedPriorO$ $[\cfSeedPriorOLo,\cfSeedPriorOHi]$ of the two-model NIH study this campaign builds on (Appendix~\ref{app:seed}). The prompt dependence of Mass also reproduces (Appendix~\ref{app:seed-results}). Effusion and Cardiomegaly are the most readable of the four concepts both chest sets share, with mean selectivity $\cfSelEffusion$ and $\cfSelCardiomegaly$ over \cfChestProbeBlocks{} chest blocks. Effusion is owned in \cfEffusionOwnedBlocks{} of \cfChestBlocks{} chest blocks. In \cfEffusionOwnedSmall{} of these blocks, $O_{\mathrm{Effusion}}\leq0.03$.

\subsection{The attribution deficit across the grid}
\label{sec:deficit}

On the chest sets, ownership is the rarest and least stable of the three grades. An effect floor of $0.05$ on both $W_{q,q}$ and $O_q$ retains \cfChestFloored{} of the \cfChestOwned{} owned chest cells (\cfChestFlooredPct\% of \cfChestCellsN{}), compared with \cfCocoFloored{} of the \cfCocoOwned{} owned COCO cells (\cfCocoFlooredPct\% of \cfCocoOwnedN{}). Ownership survives both probe refits in \cfRefitGradeStableNih{} of the \cfRefitGradeOwnedNih{} owned NIH cells, compared with \cfRefitGradeStableCoco{} of \cfRefitGradeOwnedCoco{} on COCO (Table~\ref{tab:cf-refit}). Clinical writes are not selective on the chest sets. The written concept's own question receives \cfOwnShareNih\% of the positive answer change on NIH and \cfOwnShareChex\% on CheXpert, compared with \cfOwnShareCoco\% on COCO (Figure~\ref{fig:write-flow}). This pattern is the attribution deficit. The model makes the concept readable to a probe. The probe's direction is not a handle on the concept. Writing it back moves other findings' answers as much as the concept's own.

\begin{figure}[t]
    \centering
    \resizebox{\textwidth}{!}{\includegraphics{figures/fig7_write_flow.pdf}}
    \caption{\textbf{Where the write goes.} Bands join each written direction $d$ to the questions $q$ it moves, with width proportional to $P_{q,d}$, the mean over checkpoints of $\max(W_{q,d},0)$. The percentage is the own-question share $\sum_d P_{d,d}/\sum_{q,d} P_{q,d}$.}
    \label{fig:write-flow}
\end{figure}

The pattern is visible in individual images (Figure~\ref{fig:example}). On a chest radiograph, the Effusion write raises the Effusion answer from $\cfExChestClean$ to $\cfExChestOwn$ and the Nodule write raises it further, to $\cfExChestComp$. On a natural image, the dog write moves the dog answer from $\cfExCocoClean$ to $\cfExCocoOwn$ and moves nothing else. The write matrix reveals two components of the deficit (Figure~\ref{fig:write-matrices}, Table~\ref{tab:cf-controls}). First, concept writes in the chest datasets have little effect on their own answers. The median $W_{q,q}$ is $\cfMedianWqqNih$ on NIH and $\cfMedianWqqChex$ on CheXpert, compared with $\cfMedianWqqCoco$ on COCO. Second, competing directions share the effect whenever a chest write moves an answer (Table~\ref{tab:cf-contingency}). Across the grid, the Atelectasis, Effusion, and Consolidation directions raise each other's answers by similar amounts. On COCO, the diagonal entry of $W$ dominates its row. On NIH and CheXpert, the row is flat or peaks off the diagonal. At every Qwen2.5-VL size, a competing direction raises the Effusion answer more than the Effusion direction does: Cardiomegaly at 3B, Nodule at 7B and 32B, and Atelectasis at 72B.

\subsection{The write mechanism is intact on natural objects}
\label{sec:coco}

On COCO, the same procedure yields the opposite result (Table~\ref{tab:cf-main}, right block). Object ownership holds in \cfCocoOwned{} of \cfCocoOwnedN{} cells (\cfCocoOwnedPct\%). Ownership margins span $\cfCocoMarginMin$--$\cfCocoMarginMax$ in the Qwen, Lingshu, and Gemma families. \cfCocoRefMet{} cells meet the steering reference. Only \cfCocoCompetitor{} cells have a stronger competitor. \cfCocoBlocksFourPlus{} of the \cfCocoBlocks{} checkpoints with a COCO block own at least four of the six objects. Dose response tracks the write. The dose module ran on \cfDoseBlocksNih{} NIH and \cfDoseBlocksCoco{} COCO blocks. On COCO, median ownership across questions and blocks rises from $\cfDoseCocoLow$ at $\alpha=-0.1$ to $\cfDoseCocoPeak$ at $\alpha=+0.25$. Median ownership remains positive at $\alpha=+0.5$ ($\cfDoseCocoHalf$). On NIH, median ownership stays between $\cfDoseNihMin$ and $\cfDoseNihMax$ at every dose (Figure~\ref{fig:overview}c). The write, direction construction, reference family, and decision rules thus work as intended.

The natural-image result extends to smaller, fine-grained objects. Six small or fine-grained COCO categories (handbag, book, backpack, traffic light, knife, and cell phone) are owned in \cfFgFineOwned{} of \cfFgFineCells{} cells, a rate of \cfFgFineRate{}, against \cfFgEasySameOwned{} of \cfFgEasySameCells{} (\cfFgEasySameRate{}) for the easy six in the same \cfFgSameBlocks{} checkpoints. Within a checkpoint the two counts of six differ by $\cfFgSameDiffMin$ to $\cfFgSameDiffMax$ cells, and an exact two-sided sign test over the checkpoints gives $p=\cfFgSameSignP$. Median clean-answer AUROC is \cfFgFineAuroc{} against \cfFgEasySameAuroc{}, and median probe selectivity \cfFgFineSel{} against \cfFgEasySameSel{} (Table~\ref{tab:cf-fgobj}). Most of these fine-grained cells remain owned. On the chest sets, ownership stays rare even where the model already answers the question well. At a clean-answer AUROC of at least \cfStratTopEdge{}, \cfStratChestTopOwned{} of \cfStratChestTopN{} chest cells are owned (\cfStratChestTopPct\%) against \cfStratNatTopOwned{} of \cfStratNatTopN{} natural-image cells (\cfStratNatTopPct\%), a difference of $\cfStratCvnDiff$ $[\cfStratCvnLo,\cfStratCvnHi]$. Table~\ref{tab:cf-fgobj} gives the remaining bins, the per-dataset counts, and the matched estimator, and Appendix~\ref{app:cf-crossed} states the category rule, the bin edges, the matching windows, and the bootstrap.

\begin{figure}[t]
    \centering
    \resizebox{\textwidth}{!}{\includegraphics{figures/fig3_example.pdf}}
    \caption{\textbf{Two images, six questions, three writes.} $P(\mathrm{yes})$ per question for one NIH radiograph (top) and one COCO image (bottom) under the clean input, the concept write, and the strongest competing write, at $\alpha=+0.25$ in Qwen2.5-VL-7B. The shaded row is the written concept, absent from the image; values at the bar tips give its two write answers, and the dashed line marks $P(\mathrm{yes})=0.5$.}
    \label{fig:example}
\end{figure}

\begin{figure}[t]
    \centering
    \resizebox{\textwidth}{!}{\includegraphics{figures/fig4_write_matrices.pdf}}
    \caption{\textbf{Write matrices.} $W_{q,d}$, the mean change in $P(\mathrm{yes})$ of question $q$ (rows) under the write of direction $d$ (columns), for three checkpoints on NIH (top) and three on COCO (bottom). A corner mark flags owned concepts.}
    \label{fig:write-matrices}
\end{figure}

\subsection{Checkpoints that share a vision tower own different concepts}
\label{sec:reader}

Checkpoints that share a vision tower differ in what they own. Gemma~3 shares one SigLIP tower across its 4B, 12B, and 27B checkpoints, MedGemma one MedSigLIP tower across 4B and 27B, and LLaVA-1.5 one CLIP tower across 7B and 13B. Intervening at the tower's final block changes the reader while holding the representation and the write fixed: within each family, pooled features, fitted probes, and lifted write vectors are identical across sizes, readability is identical by construction, and only the connector and language model that read them differ. Table~\ref{tab:cf-pairs} records the bitwise relation of every pair. The connector and language model produce every difference in Figure~\ref{fig:readers} by reading the same written representation differently. Steering a feature of the vision tower and letting the language model say what it sees treats that language model as a stable read-out of the tower \citep{ferrando2026explain}. Here the read-out is what varies: the same written direction, read by a larger language model of the same family, keeps or loses its advantage over the competing clinical directions. The median paired difference in ownership is \cfPairsRatioMin{}--\cfPairsRatioMax{} times the refit standard deviation, up to \cfPairsRatioSingleMax{} times for single concepts, and it exceeds that standard deviation in \cfPairsRatioAboveOne{} of \cfPairsRatioN{} pairs. In \cfPairsFiveOrSix{} of the \cfPairsN{} shared-tower pair--dataset combinations, the paired patient-bootstrap interval for $\Delta O_q$ excludes zero for five or six of the six concepts. A crossed swap of the Gemma~3 4B and MedGemma 4B towers and readers points the same way: on NIH, changing the reader at a fixed tower changes $O_q$ by a mean absolute \cfSwapReaderNihOMin{}, comparable to the \cfSwapTowerNihOMin{} from changing the tower, and no combination owns a chest finding (Appendix~\ref{app:cf-crossed}, Table~\ref{tab:cf-towerswap}).

The differences reach individual verdicts. On NIH Effusion, the identical write on the same 600 patients yields an $O_q$ difference of $\cfPairDeltaEffFourTwelve$ $[\cfPairDeltaEffFourTwelveLo,\cfPairDeltaEffFourTwelveHi]$ between Gemma~3 4B and 12B, and $\cfPairDeltaEffTwelveTwentySeven$ $[\cfPairDeltaEffTwelveTwentySevenLo,\cfPairDeltaEffTwelveTwentySevenHi]$ between 12B and 27B. In 12B, competing directions own the chest answers outright ($O_{\mathrm{Effusion}}=\cfOGemmaTwelveEff$ on both chest sets, $O_{\mathrm{Mass}}=\cfOGemmaTwelveMass$ on NIH, $O_{\mathrm{Pneumothorax}}=\cfOGemmaTwelvePneu$ on CheXpert). In 27B, no direction moves the same three answers ($\cfOGemmaTwentySevenEff$, $\cfOGemmaTwentySevenMass$, $\cfOGemmaTwentySevenPneu$). In 4B, the clean answer is yes on \cfGemmaFourYesMin--\cfGemmaFourYesMax\% of chest rows; a cell is at ceiling when the clean answer is beyond $0.99$ or $0.01$ on more than 90\% of its rows. Every own write lowers its answer's logit margin by more than any competing write raises it ($O^m_q$ from $\cfPairsGemmaFourOmNear$ to $\cfPairsGemmaFourOmFar$, all intervals below zero). MedGemma 4B and 27B differ in the same way (Pneumothorax $\cfOMedFourPneu$ versus $\cfOMedTwentySevenPneu$ and Edema $\cfOMedFourEdema$ versus $\cfOMedTwentySevenEdema$ on CheXpert). LLaVA-1.5 7B and 13B own no concepts on either chest set, yet own \cfLlavaSevenCoco{} and \cfLlavaThirteenCoco{} COCO objects with the same write. The reader effect spans the whole matrix, including cells whose verdict stays the same. Ownership is therefore a property of the checkpoint that reads the representation, which is why a probe fitted to the tower alone cannot certify it.

\begin{figure}[t]
    \centering
    \resizebox{\textwidth}{!}{\includegraphics{figures/fig5_same_write.pdf}}
    \caption{\textbf{Same write, different reader.} Ownership $O_q$ per concept on NIH, CheXpert, and COCO for Gemma~3 4B/12B/27B and MedGemma 4B/27B, whose probes and write vectors are identical across sizes within each family. Corner marks flag owned cells and hatching marks ceiling cells.}
    \label{fig:readers}
\end{figure}

\subsection{Non-clinical attributes of the radiograph share the deficit}
\label{sec:attributes}

A non-clinical property of the same film is owned as rarely as a finding. We combine projection, sex, age, and the findings into one nine-direction family. The attribute probes read their labels at least as well as the clinical probes read the findings: median probe AUROC per attribute is \cfAttrAurocMin{}--\cfAttrAurocMax{}, compared with \cfAttrClinAuroc{} for findings in the same blocks, and all \cfAttrReadable{} of \cfAttrChestAttrN{} attribute cells are readable. Attributes own \cfAttrAllAttrOwned{} of \cfAttrAllAttrN{} cells; findings own \cfAttrAllClinOwned{} of \cfAttrAllClinN{}, the same rate. Among cells whose clean answer meets the answer-capable rule, attributes own \cfAttrCapAttrOwned{} of \cfAttrCapAttrN{} and findings \cfAttrCapClinOwned{} of \cfAttrCapClinN{}. Pairing each attribute cell with clinical cells from its own block within \cfAttrPairWindow{} clean-answer AUROC yields \cfAttrPairAttrN{} pairs, with attributes owned in \cfAttrPairAttrOwned{} and findings in \cfAttrPairClinOwned{}. At a clean-answer AUROC of at least \cfStratTopEdge{}, attributes are owned in \cfStratAttrTopOwned{} of \cfStratAttrTopN{} cells and findings in \cfStratFindTopOwned{} of \cfStratFindTopN{}. On NIH, ownership is \cfAttrNihAttrOwned{}/\cfAttrNihAttrN{} for attributes and \cfAttrNihClinOwned{}/\cfAttrNihClinN{} for findings (Table~\ref{tab:cf-attr}). The deficit therefore follows the chest-radiograph representation and its reader rather than the concept's clinical status.

\subsection{The direction the reader responds to}
\label{sec:ansdir}

For each question we fit the model's own clean answer margin from the same projected, train-scaled features by ridge regression, then lift and write the fitted direction exactly as we lift and write a probe normal (Appendix~\ref{app:cf-robustness}; Table~\ref{tab:cf-ansdir}). This answer direction $a_q$ is a handle where the label direction is not. The features that identify a concept differ from the features that control the output \citep{arad2025saes}; here that separation holds against matched clinical competitors across the checkpoints of the grid. Across the \cfAnsChestN{} chest cells of \cfAnsChestBlocks{} blocks, $a_q$ is owned in \cfAnsChestOwnedA{} and the label direction in \cfAnsChestOwnedLabel{}. The answer direction's write beats the strongest logistic competitor in \cfAnsChestBeats{}. Across the five held-out templates, $a_q$ remains owned in \cfAnsdirtChestKept{} of \cfAnsdirtChestKeptN{} chest and \cfAnsdirtCocoKept{} of \cfAnsdirtCocoKeptN{} COCO (cell, template) pairs, and over all pairs it is owned in \cfAnsdirtChestOwnedPairs{} of \cfAnsdirtChestPairsAll{} chest and \cfAnsdirtCocoOwnedPairs{} of \cfAnsdirtCocoPairsAll{} COCO pairs. In Qwen2.5-VL-7B, it is owned for \cfAnsQwenOwnedNih{} of the six NIH concepts (Effusion $O^a=\cfAnsQwenEffO$ with an own write of $\cfAnsQwenEffW$) and all \cfAnsQwenOwnedChex{} CheXpert concepts. In Lingshu 32B, it is owned for \cfAnsLingOwnedNih{} NIH and \cfAnsLingOwnedChex{} CheXpert concepts. In Gemma~3 12B, it is owned for \cfAnsGemOwnedNih{} NIH and \cfAnsGemOwnedChex{} CheXpert concepts.

The answer direction reads the model, not the label. Over the \cfValAurocAnsChestN{} chest cells, its median AUROC against the report label is $\cfValAurocAnsChest$, compared with $\cfValAurocProbeChest$ for the probe; against the model's own clean answer it is $\cfValAurocAnsCleanChest$, compared with $\cfValAurocProbeCleanChest$ (Table~\ref{tab:cf-validation}). In Qwen2.5-VL-7B on NIH, the label AUROC is $\cfValAurocAnsQwenNih$ for the answer direction and $\cfValAurocProbeQwenNih$ for the probe, and the clean-answer AUROC is $\cfValAurocAnsCleanQwenNih$ and $\cfValAurocProbeCleanQwenNih$. The two directions are far apart. The median raw model-space cosine between them spans $\cfAnsCosQwenMin$--$\cfAnsCosQwenMax$ across the two chest blocks of Qwen2.5-VL-7B and $\cfAnsCosLingMin$--$\cfAnsCosLingMax$ across those of Lingshu 32B, compared with $\cfAnsCosCocoMin$--$\cfAnsCosCocoMax$ across the COCO blocks. The whitened cosine is $\cos_\Sigma(a,w)=a^{\top}\Sigma w/\sqrt{a^{\top}\Sigma a\,w^{\top}\Sigma w}$, taken between the coefficient vectors in the projected, train-scaled space under the covariance $\Sigma$ of the training features; its median over each dataset's cells is $\cfValCosSigmaChestMin$--$\cfValCosSigmaChestMax$ on the two chest sets and $\cfValCosSigmaCoco$ on COCO. The metrics differ most on the CheXpert block of Lingshu 32B, where the median over its six concepts is $\cfValCosRawLingChex$ raw and $\cfValCosSigmaLingChex$ whitened. The column-selectivity index $S_d$ of Section~\ref{sec:robust-geometry} measures the share of a direction's effect that lands on its own question. Across the \cfAnsColSelChestN{} chest cells, its median is $\cfAnsColSelChest$ for answer directions and $\cfAnsColSelLabelChest$ for label directions in the same blocks; the corresponding medians are $\cfAnsColSelCoco$ and $\cfAnsColSelLabelCoco$ over the \cfAnsColSelCocoN{} COCO cells. On the chest sets, both answer and label directions move the other five answers along with their own, and the row comparison separates them. On natural images, the two directions coincide. On chest radiographs, the direction the model answers along is not the direction the probe fits.

The separation between the two families persists at four endpoints that fit neither direction: the negated form of the concept's yes/no question, a two-alternative forced choice against the concept's strongest competitor in each presentation order, and the finding word's log-probability in a report continuation (Table~\ref{tab:cf-semend}). Across the \cfSemCells{} endpoint cells of the \cfSemBlocks{} scored blocks, the label direction is owned in \cfSemLabelOwned{} and the answer direction in \cfSemAnswerOwned{}. The negated question moves opposite to the affirmative one in \cfSemNegOpposite{} of the \cfSemNegN{} concept--block pairs. The forced choice is owned in both presentation orders in \cfSemForcedOwned{} of \cfSemForcedN{} pairs, with a mean absolute gap of \cfSemGapMin{}--\cfSemGapMax{} between orders. The report continuation is owned in \cfSemReportOwned{} of \cfSemReportCells{} cells. Endpoints that the answer direction was never fitted to thus reproduce its advantage over the label direction.

\subsection{Size and family}
\label{sec:families}

Ownership on the chest sets is concentrated among a few readers and varies by concept and dataset (Figure~\ref{fig:ladders}; Table~\ref{tab:cf-models}). Qwen2.5-VL owns no concepts at 3B or 7B. At 32B, it owns Cardiomegaly. At 72B, it owns Atelectasis ($O=\cfOQwenSeventyTwoAtel$), Cardiomegaly, and Mass. Effusion is anti-owned at every size. Lingshu, the medical Qwen2.5-VL derivative, owns three CheXpert concepts at 7B. At 32B, it owns five, the highest clinical ownership in the grid. InternVL3.5 answers every NIH question at every size. On the chest sets, its writes remain within $\cfInternvlWqqMax$ of zero. LLaVA-1.5 reads the chest concepts. It answers none above chance. It owns none.

\subsection{Medical and chest-specialised checkpoints}
\label{sec:specialisation}

Three checkpoints put medical training on the vision side: CheXagent-2-3B, CheXagent-8B, and LLaVA-Rad-7B use vision towers trained on chest radiographs. By contrast, HuatuoGPT-Vision-7B is tuned on medical data over a stock CLIP tower. We score all four on every dataset, dividing the grid into three groups: chest-trained towers, medically tuned checkpoints on general towers, and general checkpoints.

The chest-trained towers read and answer the findings best in the grid. Their median probe AUROC across \cfSpecTowerProbeN{} NIH cells is \cfSpecTowerProbeAuroc{}, against \cfSpecStockProbeAuroc{} for HuatuoGPT-Vision's stock CLIP tower and \cfSpecRestProbeAuroc{} across the \cfSpecRestProbeN{} remaining NIH cells. Their median clean-answer AUROC over the chest sets is \cfSpecTowerAnsAuroc{}, against \cfSpecMedAnsAuroc{} for the medically tuned group and \cfSpecGenAnsAuroc{} for the general group. Across the same \cfSpecTowerN{} chest cells, they read \cfSpecTowerReadable{} and answer \cfSpecTowerAnswerable{}. They own \cfSpecTowerOwned{}.

Better chest representation raises ownership but does not confer it. Across the chest sets, chest-trained towers own \cfSpecTowerOwned{} of \cfSpecTowerN{} cells, medically tuned checkpoints own \cfSpecMedOwned{} of \cfSpecMedN{}, and the general group owns \cfSpecGenOwned{} of \cfSpecGenN{}. On NIH, the corresponding counts are \cfSpecNihTowerOwned{} of \cfSpecNihTowerN{}, \cfSpecNihMedOwned{} of \cfSpecNihMedN{}, and \cfSpecNihGenOwned{} of \cfSpecNihGenN{}. Medical training moves reading and answering further than it moves writing. The group that reads and answers the findings best still loses most of its chest cells to a competing clinical direction.

The natural-image control in Section~\ref{sec:coco} rests on checkpoints that answer natural-image questions. The general group answers \cfSpecGenCocoAnswerable{} of \cfSpecGenCocoN{} COCO cells at a median clean-answer AUROC of \cfSpecGenCocoAnsAuroc{} and owns \cfSpecGenCocoOwned{}; the medically tuned group answers \cfSpecMedCocoAnswerable{} of \cfSpecMedCocoN{} at \cfSpecMedCocoAnsAuroc{} and owns \cfSpecMedCocoOwned{}. The chest-trained towers answer \cfSpecTowerCocoAnswerable{} of \cfSpecTowerCocoN{} COCO cells at \cfSpecTowerCocoAnsAuroc{} and own \cfSpecTowerCocoOwned{}, so their COCO rows give the grade little to assess.

\subsection{Controls and alternatives}
\label{sec:controls}

The reference family and additional controls behave as designed across the grid (Figure~\ref{fig:overview}d; Table~\ref{tab:cf-controls}). Owned concept writes rank 1--6 among the 120 compared effects, and the concept write typically ranks in the top decile on the chest sets. Table~\ref{tab:cf-ledger} gives the rows, the reference, the rule and the sham of each of the \cfLedgerAnalyses{} graded analyses and their \cfLedgerCells{} cells.

\label{sec:robust-geometry}

\paragraph{Geometry, scale, saturation, refits, numerics.} The chest-to-COCO contrast persists across score scales, saturation, probe refits and numerical precision (Tables~\ref{tab:cf-geometry}--\ref{tab:cf-validation}; Appendix~\ref{app:cf-robustness}). The strongest competitor matches the most similar direction or the most co-occurring label at chance rate. The logit-scale grade agrees with the probability-scale grade in \cfScaleAgree{} of \cfScaleAgreeN{} cells. COCO blocks are more saturated than chest blocks, and restricting each cell to unsaturated rows preserves the full owned grade in \cfScaleUnsatGradeKept{} of \cfScaleUnsatGradeN{} owned chest cells. Refitting the probes on resampled training patients changes $O_q$ by a median of \cfScaleRefitMedian{}. The column-selectivity index $S_d=W_{d,d}/\sum_q|W_{q,d}|$ divides a direction's signed effect on its own question by its summed absolute effect on the six questions; across cells, its median is \cfValColSelNih{} on NIH and \cfValColSelChex{} on CheXpert, versus \cfValColSelCoco{} on COCO (Table~\ref{tab:cf-validation}). Rescoring in fp32 and at batch size one changes $W$ by at most \cfPrecisionMaxDW{} and alters \cfPrecisionGradeChanges{} of \cfPrecisionGradeCells{} grades (Table~\ref{tab:cf-scale}).

\paragraph{Label source.} The deficit holds under radiologist labels. On CheXpert, restricting evaluation to rows with a known label preserves the sign of $O_q$ in \cfValKnownSignAgree{} of \cfValKnownN{} cells; NIH and COCO have every test label known, so their known-row numbers reproduce the all-row numbers. Re-grading the same fitted directions on \cfValidRows{} CheXpert films labelled by radiologist consensus \citep{irvin2019chexpert} reproduces the owned grade in \cfValidOwnedAgree{} of \cfValidCells{} cells (Table~\ref{tab:cf-valid}). Refitting the six directions on the radiologist labels themselves leaves them a median raw model-space cosine of \cfValidfitCosRaw{} and a median whitened cosine of \cfValidfitCosWhitened{} from the report-label directions, across the \cfValidfitBlocks{} CheXpert blocks with the refit. The expert-label refit reads the radiologist labels with a median cross-fitted AUROC of \cfValidfitAurocExpert{}, versus \cfValidfitAurocReportDir{} for the report-label direction on the same films, and a report-label direction fitted on as many rows as the refit reaches \cfVfArmReportSubAuroc{}: the estimation sample moves the result by $\cfVfArmDeltaSize$ and the label source by $\cfVfArmDeltaLabel$. The expert-label directions are owned in \cfValidfitOwnedExpert{} of \cfValidfitCells{} cells, the report-label directions in \cfValidfitOwnedReport{}.

\paragraph{Alternative direction constructions.} The deficit persists beyond the logistic normal. We test four alternative directions with the same projected, train-scaled features and write protocol (Table~\ref{tab:cf-altdir}): the difference of class means, the Haufe pattern $\Sigma\boldsymbol\beta_c$ \citep{haufe2014interpretation}, the normal orthogonalised against the other five normals, and the normal refitted on features residualised against the other five labels. The last two address the entanglement and the negative-set dependence that a concept vector inherits from the way it is fitted \citep{nicolson2025explaining}. Across the \cfAltdirBlocks{} scored blocks, each construction and the logistic normal own \cfAltdirChestPctMin{}--\cfAltdirChestPctMax\% of a chest dataset's cells and \cfAltdirCocoPctMin{}--\cfAltdirCocoPctMax\% of COCO cells. The difference-of-means and pattern directions sit at a median absolute raw model-space cosine of $\cfAltdirCosMin$ to $\cfAltdirCosMax$ from the logistic normal. Each of these two directions owns more NIH cells and fewer COCO cells than the logistic normal. Lifted as displacements, they respectively own \cfAltdirdChestDom{} and \cfAltdirdChestPattern{} of \cfAltdirdChestN{} chest cells and \cfAltdirdCocoDom{} and \cfAltdirdCocoPattern{} of \cfAltdirdCocoN{} COCO cells (Table~\ref{tab:cf-altdir}). Every construction preserves the chest-to-COCO gap. At least one construction owns \cfAltdirNihAny{} of \cfAltdirNihN{} cells on NIH and \cfAltdirChexAny{} of \cfAltdirChexN{} on CheXpert, compared with \cfAltdirCocoAny{} of \cfAltdirCocoN{} on COCO. Extending the clinical family to every additional dataset label with enough support retains \cfExtcompRetained{} of the \cfExtcompOwned{} owned cells in the \cfExtcompBlocks{} scored blocks (Table~\ref{tab:cf-altdir}). A concept vector also carries where in the image the concept was read \citep{nicolson2025explaining}, so we weight writes along the logistic directions by the softmax of token probe scores, or restrict them to the top quarter of tokens at four times the weight. Across the \cfTokenwChestBlocks{} chest blocks with these writes, the softmax write owns \cfTokenwSoftChestOwned{} of \cfTokenwChestCells{} cells, the top-quarter write owns \cfTokenwTopqChestOwned{}, and the uniform write owns \cfTokenwUniformChestOwned{}. On COCO, the softmax and top-quarter writes own \cfTokenwSoftCocoOwned{} and \cfTokenwTopqCocoOwned{} of \cfTokenwCocoCells{} cells (Table~\ref{tab:cf-altdir}). Concentrating the write where the probe reads does not create a clinical handle.


\FloatBarrier
\section{Discussion}
\label{sec:implications}

\paragraph{Readability belongs to the probe, ownership to the reader.}
Readability describes a classifier fitted to the representation; ownership describes how the model reading that representation responds to a write. Checkpoints sharing a vision tower have identical probes and write vectors, yet own different concepts. The chest-trained towers read and answer findings best in the grid but own only \cfSpecTowerOwned{} of \cfSpecTowerN{} chest cells. Ownership belongs to the checkpoint's reader as well as its representation.

\paragraph{On chest radiographs, the label direction and the answer direction come apart.}
On natural images, the object-label direction and the model's answer direction nearly coincide (median raw model-space cosine $\cfAnsCosCocoMed$), and writing the label direction controls the answer. On chest radiographs the two are nearly orthogonal ($\cfAnsCosChestMed$). The label direction barely moves its own answer while moving other findings' answers; the answer direction is owned in about half of chest cells. The flat rows of chest write matrices show how an effective write spreads across findings. The useful unit of control is the direction--reader pair: a legible label need not mark the pathway along which the reader chooses its answer.

\paragraph{Concept control is a relation between a direction and a reader.}
A concept handle joins a label, a write direction, and a reader. Probe accuracy measures what the label reveals; random and sham writes show that the direction reaches the answer; ownership tests whether it beats equally constructed clinical alternatives. The write matrix $W$ exposes both the intended answer change and its spread to other findings. Natural-image objects show that the same procedure can produce selective control.

\section{Limitations}
\label{sec:limitations}

We write linear directions at the consumed visual block and connector and grade six concepts per dataset through discrete answers. MAIRA-2 does not pass the template-mapping gate for report generation (Appendix~\ref{app:cf-gates}). The answer direction tracks the model's answer more closely than the clinical label: its median label AUROC is $\cfValAurocAnsChest$, against $\cfValAurocProbeChest$ for the label probe. Answer control and clinical fidelity remain distinct objectives.
\newpage
\section{Conclusion}
\label{sec:conclusion}

Reading a clinical concept from a vision-language model's visual representation does not make its probe direction a handle on that concept. Across \cfCheckpoints{} checkpoints and two chest-radiograph datasets, concepts are readable in \cfChestReadablePct\% and answerable in \cfChestAnswerablePct\% of cells but owned in \cfChestOwnedPct\%. The concept's own direction barely moves its answer, while the directions fitted for the other findings move it as much or more. The same procedure owns \cfCocoOwnedPct\% of natural-object cells; even where clean-answer AUROC is at least \cfStratTopEdge{}, ownership is \cfStratChestTopPct\% on chest radiographs and \cfStratNatTopPct\% on natural images. Checkpoints that share a vision tower, and therefore write identical vectors, own different concepts. The model's own answer direction is a handle where the label direction is not: on chest radiographs it is nearly orthogonal to the probe direction in raw model space (median cosine $\cfAnsCosChestMed$, against $\cfAnsCosCocoMed$ on natural images). A concept is a handle for a reader when its direction moves the intended answer more than competing directions do.
\section*{Reproducibility statement}
The registered protocol file fixes the loci, projection, probe settings, dose grid, reference family, templates, modules, gates, and decision rules; its identifier and the evaluation-code version are recorded in every block's scorecard. The release contains the evaluation code and adapters, the protocol file, the cohort manifests with their seeds, the fitted directions, the per-image outcomes of every module, every block's scorecard and coverage file, the rendered prompts, and the pinned checkpoint revisions (Table~\ref{tab:cf-models}); every table and figure in the paper is regenerated from those files by the scripts in the release. Appendix~\ref{app:repro} states the consumed block of every family, the consumed-token mask, the projection and probe construction, the control sampling procedures, the cohort roles, and the protocol amendments with their dates.

\section*{Ethics statement}
The study uses three public datasets under their published terms: NIH ChestX-ray14 (NIH Clinical Center), CheXpert Plus (Stanford AIMI, accessed through its data-use agreement), and COCO (annotations CC BY 4.0). All radiographs are de-identified by their providers; no protected health information was accessed, and no image or report is redistributed. The write test evaluates model behaviour under activation edits and is not a diagnostic tool; its grades describe what a probe direction does inside a model and carry no claim of clinical benefit.
\bibliography{refs}
\bibliographystyle{plainnat}

\clearpage
\appendix
\section{Campaign details}
\label{app:campaign}

This appendix lists the checkpoints and their consumed loci, the acceptance decisions, the coverage of the campaign, and the per-concept results of every completed block.

\subsection{Labels, registration, and artifact}
\label{app:cf-labels}

\paragraph{Labels.} We use each dataset's supplied report-derived labels: NIH ChestX-ray14 finding labels \citep{wang2017chestxray}, the CheXpert labeler output \citep{irvin2019chexpert} distributed with CheXpert Plus \citep{chambon2024chexpertplus}, and COCO instance annotations \citep{lin2014coco}. Rows are positive for label 1, negative for label 0, and \emph{unknown} when the labeler output is uncertain or blank, which is the uncertainty convention of the CheXpert labeler \citep{irvin2019chexpert}. We exclude unknown rows from probe fitting, selectivity, and answer AUROC. We write and score them like all other rows. Fixed manifests drawn with recorded seeds define the cohorts, one manifest per dataset for every block. No patient and no image appears in two roles of a dataset. Every cohort is disjoint from every seed-study cohort. Table~\ref{tab:cf-prevalence} reports per-concept counts for the three scored cohorts: the training cohort for fitting the scaler, probes, and answer direction; the 400-row calibration cohort for the readable and answerable grades; and the 600-row test cohort for every write. A second CheXpert cohort contains \cfValidRows{} frontal films from the CheXpert validation set, one per patient, with the radiologists' consensus labels \citep{irvin2019chexpert}. This cohort has zero patient overlap with any campaign cohort (Table~\ref{tab:cf-valid}).

\paragraph{Registration.} We froze the protocol before scoring any grid block. The loci, projection and probe settings, dose, reference family, six templates, module set, seven gates, and three decision rules match the registered protocol file and its specified code version. We registered the seed study's crossover and confirmation analyses before running them (Appendix~\ref{app:seed}). One analysis pass over the packaged outcomes produces all reported results for each block. After completing the grid, we added the leaderboard, write-flow share, and per-image examples as descriptive summaries of these outputs.

\paragraph{Artifact.} The release includes evaluation code, adapters, the registered protocol file, cohort manifests with their seeds, and fitted directions. It provides per-image outcomes for every module (write, dose, refit, locus, and template rows), scorecards and coverage files for every block, and rendered prompts. Each new checkpoint requires one adapter that identifies its input block and passes the gates.

\subsection{Grid and consumed loci}
\label{app:cf-grid}

\begin{table}[h]
\centering
\scriptsize
\setlength{\tabcolsep}{3pt}
\input{tables/cf_colors}
\caption{\textbf{Checkpoints and families.} Per checkpoint: vision tower, consumed visual tokens per image at the primary locus, input image size, eligible templates, and the pinned Hugging Face revision. Bold rows give each family's mean selectivity and clean-answer AUROC over clinical cells and its owned clinical and COCO cells. Terracotta shading darkens with the owned share.}
\label{tab:cf-models}
\begin{tabular}{llllll@{\hspace{8pt}}rrrr}
\toprule
\multirow{2}{*}{Checkpoint} & \multirow{2}{*}{Vision tower} & \multirow{2}{*}{Tokens} & \multirow{2}{*}{Image} & \multirow{2}{*}{Templates} & \multirow{2}{*}{Revision} & \multirow{2}{*}{$\bar S$} & \multirow{2}{*}{$\overline{\mathrm{AUROC}}$} & \multicolumn{2}{c}{owned cells}\\
\cmidrule(lr){9-10}
&  &  &  &  &  &  &  & \multicolumn{1}{c}{clinical} & \multicolumn{1}{c}{COCO}\\
\midrule
\multicolumn{6}{l}{\textbf{Qwen2.5-VL}~\citep{bai2025qwen25vl}} & 0.132 & 0.635 & \cellcolor{cfTerra2}5/48 & \cellcolor{cfTerra4}24/24 \\
\quad Qwen2.5-VL-3B & native ViT 675M & 576$\to$144 & $336^2$ & IYWYIAWA & \texttt{6628554} & & & & \\
\quad Qwen2.5-VL-7B & native ViT 675M & 576$\to$144 & $336^2$ & all & \texttt{cc59489} & & & & \\
\quad Qwen2.5-VL-32B & native ViT 675M & 576$\to$144 & $336^2$ & all & \texttt{7cfb30d} & & & & \\
\quad Qwen2.5-VL-72B & native ViT 675M & 576$\to$144 & $336^2$ & all & \texttt{89c8620} & & & & \\
\midrule
\multicolumn{6}{l}{\textbf{Qwen3-VL}~\citep{qwen2025qwen3vl}} & 0.137 & 0.735 & \cellcolor{cfTerra2}4/36 & \cellcolor{cfTerra4}18/18 \\
\quad Qwen3-VL-4B & SigLIP2-based & $\leq$1024 & $\leq 512^2$ & all & \texttt{ebb281e} & & & & \\
\quad Qwen3-VL-8B & SigLIP2-based & $\leq$1024 & $\leq 512^2$ & all & \texttt{0c351dd} & & & & \\
\quad Qwen3-VL-32B & SigLIP2-based & $\leq$1024 & $\leq 512^2$ & all & \texttt{0cfaf48} & & & & \\
\midrule
\multicolumn{6}{l}{\textbf{InternVL3.5}~\citep{wang2025internvl35}} & 0.125 & 0.807 & \cellcolor{cfTerra2}5/36 & \cellcolor{cfTerra4}17/18 \\
\quad InternVL3.5-8B & InternViT-300M & 1024 & $448^2$ & all & \texttt{741a7d0} & & & & \\
\quad InternVL3.5-14B & InternViT-300M & 1024 & $448^2$ & all & \texttt{226b96d} & & & & \\
\quad InternVL3.5-38B & InternViT-6B & 1024 & $448^2$ & all & \texttt{7c830fc} & & & & \\
\midrule
\multicolumn{6}{l}{\textbf{Gemma 3}~\citep{gemma2025gemma3}} & 0.124 & 0.564 & \cellcolor{cfTerra2}2/36 & \cellcolor{cfTerra4}17/18 \\
\quad Gemma 3 4B & SigLIP-So400m & 4096 & $896^2$ & all & \texttt{093f9f3} & & & & \\
\quad Gemma 3 12B & SigLIP-So400m & 4096 & $896^2$ & all & \texttt{96b6f1e} & & & & \\
\quad Gemma 3 27B & SigLIP-So400m & 4096 & $896^2$ & all & \texttt{005ad34} & & & & \\
\midrule
\multicolumn{6}{l}{\textbf{MedGemma} (medical)~\citep{sellergren2025medgemma}} & 0.164 & 0.835 & \cellcolor{cfTerra2}2/24 & \cellcolor{cfTerra3}10/12 \\
\quad MedGemma 4B & MedSigLIP & 4096 & $896^2$ & all & \texttt{290cda5} & & & & \\
\quad MedGemma 27B & MedSigLIP & 4096 & $896^2$ & all & \texttt{2d3e00e} & & & & \\
\midrule
\multicolumn{6}{l}{\textbf{Lingshu} (medical)~\citep{lasa2025lingshu}} & 0.142 & 0.820 & \cellcolor{cfTerra4}11/24 & \cellcolor{cfTerra4}12/12 \\
\quad Lingshu 7B & Qwen2.5-VL ViT & 576$\to$144 & $336^2$ & all & \texttt{b98aecd} & & & & \\
\quad Lingshu 32B & Qwen2.5-VL ViT & 576$\to$144 & $336^2$ & all & \texttt{36b9827} & & & & \\
\midrule
\multicolumn{6}{l}{\textbf{LLaVA-1.5}~\citep{liu2024improved}} & 0.124 & 0.516 & \cellcolor{cfTerra1}0/24 & \cellcolor{cfTerra3}7/12 \\
\quad LLaVA-1.5-7B & CLIP ViT-L/14-336 & 576 & $336^2$ & IYWY & \texttt{b234b80} & & & & \\
\quad LLaVA-1.5-13B & CLIP ViT-L/14-336 & 576 & $336^2$ & IYWYIAWA & \texttt{5dda288} & & & & \\
\midrule
\multicolumn{6}{l}{\textbf{LLaVA-Med} (medical)~\citep{li2023llavamed}} & 0.136 & 0.560 & \cellcolor{cfTerra2}1/12 & \cellcolor{cfTerra2}1/6 \\
\quad LLaVA-Med 1.5 7B & CLIP ViT-L/14-336 & 576 & $336^2$ & IYWY & \texttt{91bb16c} & & & & \\
\midrule
\multicolumn{6}{l}{\textbf{Llama 3.2 Vision}~\citep{meta2024llama32vision}} & 0.119 & 0.688 & \cellcolor{cfTerra1}0/12 & \cellcolor{cfTerra2}1/6 \\
\quad Llama 3.2 Vision 11B & ViT-H/14 & 1601 & $560^2$ & IBWB & \texttt{9eb2daa} & & & & \\
\midrule
\multicolumn{6}{l}{\textbf{CheXagent} (medical)~\citep{chen2024chexagent}} & 0.180 & 0.857 & \cellcolor{cfTerra3}7/24 & \cellcolor{cfTerra2}2/12 \\
\quad CheXagent-2-3B & XraySigLIP ViT-L/16 & 1024 & $512^2$ & IYWYIAWA & \texttt{8f19b53} & & & & \\
\quad CheXagent-8B & EVA-CLIP ViT 1.4B & 1025$\to$128 & $448^2$ & IYWYIAWA & \texttt{4934e91} & & & & \\
\midrule
\multicolumn{6}{l}{\textbf{LLaVA-Rad} (medical)~\citep{zambranochaves2025llavarad}} & 0.209 & 0.793 & \cellcolor{cfTerra1}0/12 & \cellcolor{cfTerra1}0/6 \\
\quad LLaVA-Rad-7B & BiomedCLIP-CXR-518 & 1369 & $518^2$ & IYWY & \texttt{dcdbc6c} & & & & \\
\midrule
\multicolumn{6}{l}{\textbf{HuatuoGPT-Vision} (medical)~\citep{chen2024huatuogptvision}} & 0.136 & 0.746 & \cellcolor{cfTerra2}1/12 & \cellcolor{cfTerra3}4/6 \\
\quad HuatuoGPT-Vision-7B & CLIP ViT-L/14-336 & 576 & $336^2$ & all & \texttt{34dfcdb} & & & & \\
\bottomrule
\end{tabular}
\end{table}
\begin{figure}[h]
    \centering
    \resizebox{\textwidth}{!}{\includegraphics{figures/fig6_ladders.pdf}}
    \caption{\textbf{Size ladders.} Ownership $O_q$ per concept across sizes for Qwen2.5-VL on NIH, Qwen3-VL on NIH, and Lingshu on CheXpert. The concept with the largest ownership at the largest size is drawn as a thick solid line (Atelectasis, Nodule, Edema), the other concepts dashed; filled markers are owned cells. Each panel has its own vertical scale.}
    \label{fig:ladders}
\end{figure}

\begin{table}[h]
\centering
\scriptsize
\setlength{\tabcolsep}{4pt}
\caption{\textbf{Label counts per cohort and concept.} Positive / negative / unknown rows of the training, calibration, and test cohorts of the fixed manifests. Unknown means the CheXpert labeler was uncertain or blank.}
\label{tab:cf-prevalence}
\begin{tabular}{llrrrrrrrrr}
\toprule
Dataset & Concept & \multicolumn{3}{c}{train (pos / neg / unk)} & \multicolumn{3}{c}{calibration (pos / neg / unk)} & \multicolumn{3}{c}{test (pos / neg / unk)} \\
\cmidrule(lr){3-5}\cmidrule(lr){6-8}\cmidrule(lr){9-11}
NIH ChestX-ray14 & Effusion & 2423 & 15789 & 0 & 20 & 380 & 0 & 46 & 554 & 0 \\
 & Atelectasis & 2311 & 15901 & 0 & 33 & 367 & 0 & 57 & 543 & 0 \\
 & Pneumothorax & 1254 & 16958 & 0 & 13 & 387 & 0 & 17 & 583 & 0 \\
 & Cardiomegaly & 881 & 17331 & 0 & 15 & 385 & 0 & 23 & 577 & 0 \\
 & Mass & 1263 & 16949 & 0 & 15 & 385 & 0 & 32 & 568 & 0 \\
 & Nodule & 1549 & 16663 & 0 & 27 & 373 & 0 & 47 & 553 & 0 \\
\addlinespace
CheXpert Plus & Effusion & 5897 & 4277 & 9826 & 117 & 93 & 190 & 158 & 130 & 312 \\
 & Atelectasis & 2891 & 86 & 17023 & 67 & 4 & 329 & 84 & 2 & 514 \\
 & Pneumothorax & 1086 & 6456 & 12458 & 17 & 136 & 247 & 32 & 213 & 355 \\
 & Cardiomegaly & 2312 & 1903 & 15785 & 49 & 46 & 305 & 71 & 58 & 471 \\
 & Consolidation & 966 & 3652 & 15382 & 25 & 83 & 292 & 25 & 113 & 462 \\
 & Edema & 3814 & 2231 & 13955 & 70 & 54 & 276 & 133 & 59 & 408 \\
\addlinespace
COCO & person & 10877 & 9123 & 0 & 228 & 172 & 0 & 321 & 279 & 0 \\
 & dog & 693 & 19307 & 0 & 10 & 390 & 0 & 15 & 585 & 0 \\
 & car & 2051 & 17949 & 0 & 48 & 352 & 0 & 60 & 540 & 0 \\
 & chair & 2183 & 17817 & 0 & 59 & 341 & 0 & 63 & 537 & 0 \\
 & bottle & 1455 & 18545 & 0 & 34 & 366 & 0 & 47 & 553 & 0 \\
 & bicycle & 567 & 19433 & 0 & 9 & 391 & 0 & 14 & 586 & 0 \\
\bottomrule
\end{tabular}
\end{table}

The Llama~3.2 Vision projector also reads five intermediate local-layer outputs that bypass the final global block. LLaVA computes the final CLIP block but consumes the penultimate block. At every checkpoint, Gate (C) verifies that the write reaches the consumer and the answer logits and changes nothing outside the consumed tokens.

\subsection{Acceptance decisions}
\label{app:cf-gates}

All \cfWriteBlocks{} completed blocks pass checks (A), (B), (C), and (G): determinism error is zero, change at $\alpha=0$ is zero, and fp32-versus-model logit differences are below \cfGateFpMax{}. In check (D), the Qwen3-VL, Gemma~3, MedGemma, InternVL, and Llama families exceed the 0.25-logit tolerance (up to \cfGateBatchMax{} logits for Gemma~3 12B). We record these as declared deviations. Fixed batch composition neutralises them. Check (E) fails across all datasets for the A/B templates for LLaVA-1.5-7B and LLaVA-Med, the B-mapped templates for Qwen2.5-VL-3B and LLaVA-1.5-13B, and the yes/no and A-mapped templates for Llama~3.2 Vision 11B. We record affected cells as ineligible. Check (E) fails for every template of MAIRA-2 \citep{bannur2024maira} on NIH: the model answers ``No'' to every finding on all \cfSpecUngradedRows{} preflight images and ``A'' to every forced choice under both letter mappings, so no template carries a graded answer and the block produces no outcomes. Two Gemma~3 checkpoints answer ``yes'' to most chest questions; Gemma~3 4B does so on \cfGemmaFourYesMin--\cfGemmaFourYesMax\% of rows (clean $P(\mathrm{yes})>0.99$). The answerability grade records this behaviour. Three checkpoints (LLaVA-1.5 7B/13B, LLaVA-Med) answer at chance.

\subsection{Coverage and cost}
\label{app:cf-coverage}

Every checkpoint has a completed core write matrix and calibration pass on NIH ChestX-ray14, CheXpert Plus, and COCO (Table~\ref{tab:cf-coverage}). The campaign targets \cfCoveragePlannedWord{} modules per block: the write matrix, the calibration pass, the five additional templates, the dose sweep, the two probe refits, the connector locus, and the connector-locus calibration. \cfCoverageBlocksAllSeven{} of \cfBlocks{} blocks carry all seven modules. Qwen2.5-VL-72B on CheXpert Plus is the partial block: its core write matrix, calibration, dose sweep, refits, and connector-locus calibration are complete; its prompt-template and connector-locus modules are not. Llama~3.2 Vision 11B is graded on its eligible B-positive template IB. The campaign used about 3,000 A100-hours. A chest block for a 7B checkpoint takes about ten GPU-hours.

\begin{table}[htbp]
\centering
\scriptsize
\setlength{\tabcolsep}{4pt}
\caption{\textbf{Coverage.} Modules completed per block at packaging, from the campaign manifest. A module whose template is ineligible reads inel., a block not run reads ``--'', and the primary template is named when it is not the \emph{is}/yes-no template.}
\label{tab:cf-coverage}
\begin{tabular}{llll}
\toprule
Checkpoint & NIH ChestX-ray14 & CheXpert Plus & COCO\\
\midrule
Qwen2.5-VL-3B & C Cal P D R L Lc A Ad & C Cal P D R L Lc A Ad & C Cal P D R L Lc A Ad \\
Qwen2.5-VL-7B & C Cal P D R L Lc A Ad An E T Pr & C Cal P D R L Lc A Ad An E T & C Cal P D R L Lc A Ad An T Pr \\
Qwen2.5-VL-32B & C Cal P D R L Lc A Ad & C Cal P D R L Lc A Ad & C Cal P D R L Lc A Ad \\
Qwen2.5-VL-72B & C Cal P D R L Lc & C Cal D R Lc & C Cal P D R L Lc \\
Qwen3-VL-4B & C Cal P D R L Lc A Ad & C Cal P D R L Lc A Ad & C Cal P D R L Lc A Ad \\
Qwen3-VL-8B & C Cal P D R L Lc A Ad E T Pr & C Cal P D R L Lc A Ad & C Cal P D R L Lc A Ad T \\
Qwen3-VL-32B & C Cal P D R L Lc A Ad & C Cal P D R L Lc A Ad & C Cal P D R L Lc A Ad \\
InternVL3.5-8B & C Cal P D R L Lc A Ad E Pr & C Cal P D R L Lc A Ad & C Cal P D R L Lc A Ad \\
InternVL3.5-14B & C Cal P D R L Lc A Ad & C Cal P D R L Lc A Ad & C Cal P D R L Lc A Ad \\
InternVL3.5-38B & C Cal P D R L Lc A Ad & C Cal P D R L Lc A Ad & C Cal P D R L Lc A Ad \\
Gemma 3 4B & C Cal P D R L Lc A Ad Pr & C Cal P D R L Lc A Ad & C Cal P D R L Lc A Ad \\
Gemma 3 12B & C Cal P D R L Lc A Ad An E T Pr & C Cal P D R L Lc A Ad An E & C Cal P D R L Lc A Ad An \\
Gemma 3 27B & C Cal P D R L Lc A Ad Pr & C Cal P D R L Lc A Ad & C Cal P D R L Lc A Ad \\
MedGemma 4B & C Cal P D R L Lc A Ad Pr & C Cal P D R L Lc A Ad E & C Cal P D R L Lc A Ad \\
MedGemma 27B & C Cal P D R L Lc A Ad & C Cal P D R L Lc A Ad & C Cal P D R L Lc A Ad \\
Lingshu 7B & C Cal P D R L Lc A Ad & C Cal P D R L Lc A Ad & C Cal P D R L Lc A Ad \\
Lingshu 32B & C Cal P D R L Lc A Ad An E T Pr & C Cal P D R L Lc A Ad An E T & C Cal P D R L Lc A Ad An \\
LLaVA-1.5-7B & C Cal P D R L Lc A Ad & C Cal P D R L Lc A Ad & C Cal P D R L Lc A Ad \\
LLaVA-1.5-13B & C Cal P D R L Lc A Ad & C Cal P D R L Lc A Ad & C Cal P D R L Lc A Ad \\
LLaVA-Med 1.5 7B & C Cal P D R L Lc A Ad & C Cal P D R L Lc A Ad & C Cal P D R L Lc A Ad \\
CheXagent-2-3B & C Cal P D R L Lc & C Cal P D R L Lc & C Cal P D R L Lc \\
CheXagent-8B & C Cal P D R L Lc & C Cal P D R L Lc & C Cal P D R L Lc \\
LLaVA-Rad-7B & C Cal P D R L Lc & C Cal P D R L Lc & C Cal P D R L Lc \\
HuatuoGPT-Vision-7B & C Cal P D R L Lc & C Cal P D R L Lc & C Cal P D R L Lc \\
Llama 3.2 Vision 11B & C Cal P D R L Lc A [IB] & C Cal P D R L Lc A [IB] & C Cal P D R L Lc A [IB] \\
\midrule
\multicolumn{4}{p{0.95\textwidth}}{\textit{Legend:} C = write matrix, Cal = calibration pass, P = five additional templates, D = dose sweep, R = two probe refits, L = connector locus, Lc = connector-locus calibration, A = alternative direction constructions, Ad = displacement lifts, An = answer direction, E = extended competitor family, T = token-weighted writes, Pr = fp32 and batch-size-one rescoring.} \\
\bottomrule
\end{tabular}
\end{table}

\subsection{Ownership of every concept}
\label{app:cf-tables}

Tables~\ref{tab:cf-own-nih}--\ref{tab:cf-own-coco} report the ownership contrast of every concept cell, one table per dataset. Its cells are shaded by magnitude at the thresholds $0.02$, $0.10$, and $0.25$, terracotta for positive and slate for negative ownership, and its checkpoints follow the order of Table~\ref{tab:cf-main}. Figure~\ref{fig:cf-w-structure} summarises the structure of the underlying write matrices. Figure~\ref{fig:cf-examples} shows grades for individual images.

\begin{table}[htbp]
\centering
\scriptsize
\setlength{\tabcolsep}{4pt}
\caption{\textbf{Ownership of every concept, NIH ChestX-ray14.} Each cell is the ownership contrast $O_q$ at $\alpha=+0.25$: bold is owned, a dagger marks a stronger competitor, a double dagger an unresolved verdict, and the last row counts the concept's owned cells. Terracotta shades positive and blue negative ownership, darker for larger magnitude.}
\label{tab:cf-own-nih}
\begin{tabular}{lrrrrrr}
\toprule
Checkpoint & \multicolumn{1}{c}{Effusion} & \multicolumn{1}{c}{Atelectasis} & \multicolumn{1}{c}{Pneumothorax} & \multicolumn{1}{c}{Cardiomegaly} & \multicolumn{1}{c}{Mass} & \multicolumn{1}{c}{Nodule}\\
\midrule
Lingshu 32B & \cellcolor{cfSlate1}-0.01 & \cellcolor{cfSlate2}-0.08 & \cellcolor{cfTerra3}\textbf{+0.14} & \cellcolor{cfTerra2}\textbf{+0.05} & \cellcolor{cfSlate2}-0.08 & \cellcolor{cfSlate3}-0.18 \\
CheXagent-2-3B & \cellcolor{cfSlate1}-0.00$^{\ddagger}$ & \cellcolor{cfTerra2}\textbf{+0.04} & \cellcolor{cfSlate1}-0.00$^{\dagger}$ & \cellcolor{cfSlate1}-0.01 & \cellcolor{cfTerra1}\textbf{+0.01} & \cellcolor{cfSlate1}-0.01 \\
Qwen2.5-VL-72B & \cellcolor{cfSlate4}-0.32 & \cellcolor{cfTerra4}\textbf{+0.54} & \cellcolor{cfSlate4}-0.50 & \cellcolor{cfTerra4}\textbf{+0.26} & \cellcolor{cfTerra2}\textbf{+0.07} & \cellcolor{cfSlate3}-0.23 \\
Lingshu 7B & \cellcolor{cfTerra1}+0.02 & \cellcolor{cfTerra1}+0.01 & \cellcolor{cfTerra1}+0.01 & \cellcolor{cfTerra3}\textbf{+0.17} & \cellcolor{cfSlate2}-0.09 & \cellcolor{cfTerra1}+0.02 \\
Qwen3-VL-4B & \cellcolor{cfSlate2}-0.03$^{\dagger}$ & \cellcolor{cfSlate2}-0.07 & \cellcolor{cfSlate4}-0.29 & \cellcolor{cfSlate2}-0.08 & \cellcolor{cfSlate3}-0.13 & \cellcolor{cfSlate2}-0.09 \\
Qwen3-VL-32B & \cellcolor{cfSlate3}-0.14 & \cellcolor{cfSlate2}-0.04 & \cellcolor{cfTerra1}+0.00$^{\ddagger}$ & \cellcolor{cfSlate2}-0.04 & \cellcolor{cfSlate3}-0.11 & \cellcolor{cfTerra3}\textbf{+0.17} \\
InternVL3.5-38B & \cellcolor{cfTerra1}+0.00$^{\ddagger}$ & \cellcolor{cfSlate1}-0.02 & \cellcolor{cfSlate1}-0.00 & \cellcolor{cfSlate2}-0.03 & \cellcolor{cfTerra1}+0.02 & \cellcolor{cfTerra1}+0.00 \\
Gemma 3 12B & \cellcolor{cfSlate4}-0.83 & \cellcolor{cfSlate4}-0.57 & \cellcolor{cfSlate3}-0.17 & \cellcolor{cfSlate3}-0.19 & \cellcolor{cfSlate4}-0.65 & \cellcolor{cfTerra3}\textbf{+0.17} \\
MedGemma 4B & \cellcolor{cfTerra1}+0.01 & \cellcolor{cfSlate4}-0.56 & \cellcolor{cfTerra1}+0.01 & \cellcolor{cfSlate2}-0.07 & \cellcolor{cfSlate4}-0.27 & \cellcolor{cfTerra2}+0.09 \\
InternVL3.5-14B & \cellcolor{cfTerra1}\textbf{+0.01} & \cellcolor{cfSlate1}-0.01 & \cellcolor{cfSlate1}-0.00 & \cellcolor{cfSlate1}-0.00 & \cellcolor{cfTerra1}+0.00 & \cellcolor{cfSlate1}-0.01 \\
Qwen2.5-VL-32B & \cellcolor{cfSlate3}-0.22 & \cellcolor{cfSlate2}-0.06 & \cellcolor{cfSlate2}-0.06 & \cellcolor{cfTerra3}\textbf{+0.11} & \cellcolor{cfSlate2}-0.08 & \cellcolor{cfSlate2}-0.05 \\
InternVL3.5-8B & \cellcolor{cfTerra1}+0.00 & \cellcolor{cfSlate1}-0.02 & \cellcolor{cfSlate1}-0.00 & \cellcolor{cfTerra1}+0.00 & \cellcolor{cfSlate1}-0.01 & \cellcolor{cfSlate1}-0.01 \\
HuatuoGPT-Vision-7B & \cellcolor{cfSlate2}-0.04 & \cellcolor{cfSlate2}-0.03 & \cellcolor{cfSlate2}-0.03 & \cellcolor{cfSlate2}-0.08 & \cellcolor{cfSlate2}-0.04 & \cellcolor{cfTerra1}\textbf{+0.01} \\
LLaVA-Med 1.5 7B & \cellcolor{cfSlate1}-0.00 & \cellcolor{cfSlate1}-0.00 & \cellcolor{cfSlate2}-0.04 & \cellcolor{cfSlate1}-0.01 & \cellcolor{cfSlate1}-0.01 & \cellcolor{cfSlate2}-0.02 \\
CheXagent-8B & \cellcolor{cfTerra1}+0.01 & \cellcolor{cfSlate1}-0.01 & \cellcolor{cfSlate1}-0.00 & \cellcolor{cfSlate1}-0.00 & \cellcolor{cfTerra1}+0.00 & \cellcolor{cfSlate1}-0.01 \\
Qwen2.5-VL-3B & \cellcolor{cfSlate2}-0.09 & \cellcolor{cfSlate2}-0.04 & \cellcolor{cfSlate3}-0.18 & \cellcolor{cfSlate3}-0.11 & \cellcolor{cfSlate2}-0.04 & \cellcolor{cfSlate2}-0.04 \\
Qwen2.5-VL-7B & \cellcolor{cfSlate2}-0.07$^{\dagger}$ & \cellcolor{cfSlate3}-0.23 & \cellcolor{cfSlate3}-0.10 & \cellcolor{cfSlate3}-0.22 & \cellcolor{cfSlate3}-0.17 & \cellcolor{cfSlate2}-0.07 \\
Qwen3-VL-8B & \cellcolor{cfTerra2}+0.07 & \cellcolor{cfSlate1}-0.01 & \cellcolor{cfSlate2}-0.07 & \cellcolor{cfTerra1}+0.00 & \cellcolor{cfSlate3}-0.11 & \cellcolor{cfTerra2}+0.03 \\
Gemma 3 27B & \cellcolor{cfSlate2}-0.04 & \cellcolor{cfSlate2}-0.03 & \cellcolor{cfSlate1}-0.01 & \cellcolor{cfTerra1}+0.00$^{\ddagger}$ & \cellcolor{cfSlate2}-0.08 & \cellcolor{cfSlate1}-0.02 \\
Gemma 3 4B & \cellcolor{cfSlate1}-0.00 & \cellcolor{cfSlate1}-0.00 & \cellcolor{cfSlate1}-0.01 & \cellcolor{cfSlate3}-0.24 & \cellcolor{cfSlate4}-0.28 & \cellcolor{cfSlate3}-0.12 \\
LLaVA-1.5-7B & \cellcolor{cfSlate2}-0.03 & \cellcolor{cfSlate1}-0.01 & \cellcolor{cfSlate2}-0.08 & \cellcolor{cfSlate3}-0.13 & \cellcolor{cfTerra2}+0.02 & \cellcolor{cfSlate1}-0.02 \\
MedGemma 27B & \cellcolor{cfSlate2}-0.06 & \cellcolor{cfSlate4}-0.45 & \cellcolor{cfSlate1}-0.01 & \cellcolor{cfTerra1}+0.00 & \cellcolor{cfSlate3}-0.13 & \cellcolor{cfSlate3}-0.20 \\
LLaVA-1.5-13B & \cellcolor{cfSlate1}-0.01 & \cellcolor{cfSlate2}-0.05 & \cellcolor{cfSlate2}-0.04 & \cellcolor{cfSlate2}-0.03 & \cellcolor{cfTerra1}+0.01 & \cellcolor{cfSlate1}-0.02 \\
Llama 3.2 Vision 11B & \cellcolor{cfSlate2}-0.04 & \cellcolor{cfSlate2}-0.07 & \cellcolor{cfSlate2}-0.05 & \cellcolor{cfTerra2}+0.02 & \cellcolor{cfSlate2}-0.06 & \cellcolor{cfSlate1}-0.01 \\
LLaVA-Rad-7B & \cellcolor{cfSlate1}-0.02 & \cellcolor{cfSlate1}-0.01 & \cellcolor{cfSlate1}-0.00 & \cellcolor{cfTerra1}+0.00 & \cellcolor{cfSlate1}-0.01 & \cellcolor{cfSlate2}-0.03 \\
\midrule
\textbf{Owned} & 1 & 2 & 1 & 4 & 2 & 3 \\
\bottomrule
\end{tabular}
\end{table}

\begin{table}[htbp]
\centering
\scriptsize
\setlength{\tabcolsep}{4pt}
\caption{\textbf{Ownership of every concept, CheXpert Plus.} Each cell is the ownership contrast $O_q$ at $\alpha=+0.25$: bold is owned, a dagger marks a stronger competitor, a double dagger an unresolved verdict, and the last row counts the concept's owned cells. Terracotta shades positive and blue negative ownership, darker for larger magnitude.}
\label{tab:cf-own-chexpert}
\begin{tabular}{lrrrrrr}
\toprule
Checkpoint & \multicolumn{1}{c}{Effusion} & \multicolumn{1}{c}{Atelectasis} & \multicolumn{1}{c}{Pneumothorax} & \multicolumn{1}{c}{Cardiomegaly} & \multicolumn{1}{c}{Consolidation} & \multicolumn{1}{c}{Edema}\\
\midrule
Lingshu 32B & \cellcolor{cfTerra2}\textbf{+0.03} & \cellcolor{cfSlate2}-0.04$^{\dagger}$ & \cellcolor{cfTerra2}\textbf{+0.02} & \cellcolor{cfTerra3}\textbf{+0.11} & \cellcolor{cfTerra2}\textbf{+0.03} & \cellcolor{cfTerra3}\textbf{+0.11} \\
CheXagent-2-3B & \cellcolor{cfTerra2}\textbf{+0.03} & \cellcolor{cfSlate2}-0.02 & \cellcolor{cfTerra1}\textbf{+0.01} & \cellcolor{cfTerra1}\textbf{+0.02} & \cellcolor{cfTerra1}\textbf{+0.01} & \cellcolor{cfSlate1}-0.00$^{\ddagger}$ \\
Qwen2.5-VL-72B & \cellcolor{cfSlate3}-0.12 & \cellcolor{cfTerra2}+0.09 & \cellcolor{cfTerra2}\textbf{+0.03} & \cellcolor{cfSlate3}-0.11 & \cellcolor{cfSlate4}-0.29 & \cellcolor{cfSlate3}-0.13 \\
Lingshu 7B & \cellcolor{cfTerra2}\textbf{+0.10} & \cellcolor{cfSlate2}-0.09 & \cellcolor{cfTerra2}\textbf{+0.09} & \cellcolor{cfTerra1}+0.01$^{\ddagger}$ & \cellcolor{cfSlate2}-0.05 & \cellcolor{cfTerra3}\textbf{+0.21} \\
Qwen3-VL-4B & \cellcolor{cfSlate1}-0.01$^{\dagger}$ & \cellcolor{cfSlate1}-0.02 & \cellcolor{cfSlate3}-0.23 & \cellcolor{cfSlate1}-0.01 & \cellcolor{cfTerra2}\textbf{+0.09} & \cellcolor{cfTerra2}\textbf{+0.09} \\
Qwen3-VL-32B & \cellcolor{cfSlate2}-0.04 & \cellcolor{cfTerra1}+0.00$^{\ddagger}$ & \cellcolor{cfTerra1}+0.00 & \cellcolor{cfSlate2}-0.03 & \cellcolor{cfSlate3}-0.12 & \cellcolor{cfTerra2}\textbf{+0.05} \\
InternVL3.5-38B & \cellcolor{cfTerra1}\textbf{+0.02} & \cellcolor{cfSlate1}-0.01 & \cellcolor{cfSlate1}-0.00 & \cellcolor{cfSlate1}-0.00$^{\ddagger}$ & \cellcolor{cfTerra1}\textbf{+0.01} & \cellcolor{cfSlate2}-0.02 \\
Gemma 3 12B & \cellcolor{cfSlate4}-0.83 & \cellcolor{cfSlate3}-0.13 & \cellcolor{cfSlate4}-0.71 & \cellcolor{cfSlate4}-0.47 & \cellcolor{cfSlate4}-0.26 & \cellcolor{cfTerra3}\textbf{+0.15} \\
MedGemma 4B & \cellcolor{cfSlate3}-0.20 & \cellcolor{cfSlate2}-0.08 & \cellcolor{cfSlate4}-0.60 & \cellcolor{cfTerra2}\textbf{+0.04} & \cellcolor{cfSlate3}-0.20 & \cellcolor{cfTerra3}\textbf{+0.10} \\
InternVL3.5-14B & \cellcolor{cfTerra1}\textbf{+0.01} & \cellcolor{cfSlate1}-0.00 & \cellcolor{cfTerra1}+0.00 & \cellcolor{cfSlate1}-0.01 & \cellcolor{cfTerra1}+0.01 & \cellcolor{cfSlate1}-0.00 \\
Qwen2.5-VL-32B & \cellcolor{cfSlate2}-0.07 & \cellcolor{cfSlate3}-0.19 & \cellcolor{cfSlate3}-0.12 & \cellcolor{cfSlate3}-0.19 & \cellcolor{cfSlate2}-0.05 & \cellcolor{cfSlate2}-0.02 \\
InternVL3.5-8B & \cellcolor{cfSlate1}-0.01 & \cellcolor{cfSlate1}-0.01 & \cellcolor{cfTerra1}\textbf{+0.00} & \cellcolor{cfSlate1}-0.00 & \cellcolor{cfTerra1}+0.01 & \cellcolor{cfSlate1}-0.00 \\
HuatuoGPT-Vision-7B & \cellcolor{cfSlate2}-0.02 & \cellcolor{cfSlate1}-0.02 & \cellcolor{cfSlate1}-0.00 & \cellcolor{cfTerra1}+0.01 & \cellcolor{cfSlate2}-0.03 & \cellcolor{cfSlate1}-0.01 \\
LLaVA-Med 1.5 7B & \cellcolor{cfTerra1}\textbf{+0.00} & \cellcolor{cfSlate1}-0.00 & \cellcolor{cfSlate2}-0.02 & \cellcolor{cfSlate2}-0.03 & \cellcolor{cfSlate1}-0.00$^{\dagger}$ & \cellcolor{cfSlate1}-0.02 \\
CheXagent-8B & \cellcolor{cfSlate2}-0.02 & \cellcolor{cfSlate1}-0.02 & \cellcolor{cfSlate2}-0.02 & \cellcolor{cfSlate2}-0.07 & \cellcolor{cfTerra2}\textbf{+0.04} & \cellcolor{cfSlate2}-0.02 \\
Qwen2.5-VL-3B & \cellcolor{cfSlate3}-0.18 & \cellcolor{cfSlate1}-0.02 & \cellcolor{cfSlate2}-0.04 & \cellcolor{cfSlate2}-0.05 & \cellcolor{cfSlate2}-0.03 & \cellcolor{cfSlate2}-0.08 \\
Qwen2.5-VL-7B & \cellcolor{cfTerra1}+0.01 & \cellcolor{cfSlate2}-0.05 & \cellcolor{cfSlate2}-0.10 & \cellcolor{cfSlate1}-0.00 & \cellcolor{cfSlate3}-0.15 & \cellcolor{cfSlate3}-0.18 \\
Qwen3-VL-8B & \cellcolor{cfSlate2}-0.07$^{\dagger}$ & \cellcolor{cfSlate1}-0.00$^{\dagger}$ & \cellcolor{cfSlate2}-0.10 & \cellcolor{cfSlate1}-0.00 & \cellcolor{cfSlate2}-0.04$^{\dagger}$ & \cellcolor{cfSlate2}-0.06 \\
Gemma 3 27B & \cellcolor{cfSlate2}-0.03 & \cellcolor{cfSlate1}-0.00 & \cellcolor{cfSlate1}-0.01 & \cellcolor{cfSlate1}-0.00 & \cellcolor{cfSlate2}-0.03 & \cellcolor{cfTerra1}+0.00$^{\ddagger}$ \\
Gemma 3 4B & \cellcolor{cfSlate1}-0.00 & \cellcolor{cfSlate1}-0.00 & \cellcolor{cfSlate4}-0.65 & \cellcolor{cfSlate1}-0.02 & \cellcolor{cfSlate4}-0.35 & \cellcolor{cfSlate4}-0.58 \\
LLaVA-1.5-7B & \cellcolor{cfTerra1}+0.00 & \cellcolor{cfSlate2}-0.09 & \cellcolor{cfSlate2}-0.10 & \cellcolor{cfTerra2}+0.03 & \cellcolor{cfSlate2}-0.03 & \cellcolor{cfSlate1}-0.00 \\
MedGemma 27B & \cellcolor{cfSlate1}-0.01 & \cellcolor{cfTerra1}+0.00 & \cellcolor{cfSlate2}-0.06 & \cellcolor{cfTerra2}+0.03 & \cellcolor{cfSlate1}-0.02 & \cellcolor{cfSlate4}-0.28 \\
LLaVA-1.5-13B & \cellcolor{cfSlate1}-0.00 & \cellcolor{cfSlate2}-0.03 & \cellcolor{cfSlate1}-0.01 & \cellcolor{cfTerra1}+0.00 & \cellcolor{cfTerra1}+0.01 & \cellcolor{cfSlate1}-0.01 \\
Llama 3.2 Vision 11B & \cellcolor{cfSlate1}-0.01 & \cellcolor{cfSlate2}-0.05 & \cellcolor{cfSlate2}-0.08 & \cellcolor{cfSlate2}-0.07 & \cellcolor{cfSlate2}-0.05 & \cellcolor{cfSlate1}-0.01 \\
LLaVA-Rad-7B & \cellcolor{cfSlate2}-0.04 & \cellcolor{cfSlate1}-0.01 & \cellcolor{cfTerra1}+0.00 & \cellcolor{cfSlate2}-0.07 & \cellcolor{cfTerra2}+0.02 & \cellcolor{cfSlate1}-0.00 \\
\midrule
\textbf{Owned} & 6 & 0 & 5 & 3 & 5 & 6 \\
\bottomrule
\end{tabular}
\end{table}

\begin{table}[htbp]
\centering
\scriptsize
\setlength{\tabcolsep}{4pt}
\caption{\textbf{Ownership of every concept, COCO.} Each cell is the ownership contrast $O_q$ at $\alpha=+0.25$: bold is owned, a dagger marks a stronger competitor, a double dagger an unresolved verdict, and the last row counts the concept's owned cells. Terracotta shades positive and blue negative ownership, darker for larger magnitude.}
\label{tab:cf-own-coco}
\begin{tabular}{lrrrrrr}
\toprule
Checkpoint & \multicolumn{1}{c}{person} & \multicolumn{1}{c}{dog} & \multicolumn{1}{c}{car} & \multicolumn{1}{c}{chair} & \multicolumn{1}{c}{bottle} & \multicolumn{1}{c}{bicycle}\\
\midrule
Lingshu 32B & \cellcolor{cfTerra4}\textbf{+0.29} & \cellcolor{cfTerra3}\textbf{+0.23} & \cellcolor{cfTerra2}\textbf{+0.08} & \cellcolor{cfTerra4}\textbf{+0.53} & \cellcolor{cfTerra4}\textbf{+0.34} & \cellcolor{cfTerra3}\textbf{+0.22} \\
CheXagent-2-3B & \cellcolor{cfTerra2}\textbf{+0.07} & \cellcolor{cfTerra2}\textbf{+0.07} & \cellcolor{cfTerra2}+0.03 & \cellcolor{cfSlate1}-0.01 & \cellcolor{cfTerra2}+0.04 & \cellcolor{cfTerra1}+0.00 \\
Qwen2.5-VL-72B & \cellcolor{cfTerra4}\textbf{+0.66} & \cellcolor{cfTerra4}\textbf{+0.44} & \cellcolor{cfTerra4}\textbf{+0.95} & \cellcolor{cfTerra4}\textbf{+0.89} & \cellcolor{cfTerra4}\textbf{+0.62} & \cellcolor{cfTerra4}\textbf{+0.40} \\
Lingshu 7B & \cellcolor{cfTerra4}\textbf{+0.39} & \cellcolor{cfTerra4}\textbf{+0.74} & \cellcolor{cfTerra4}\textbf{+0.70} & \cellcolor{cfTerra4}\textbf{+0.76} & \cellcolor{cfTerra4}\textbf{+0.64} & \cellcolor{cfTerra4}\textbf{+0.85} \\
Qwen3-VL-4B & \cellcolor{cfTerra4}\textbf{+0.46} & \cellcolor{cfTerra3}\textbf{+0.13} & \cellcolor{cfTerra4}\textbf{+0.33} & \cellcolor{cfTerra4}\textbf{+0.71} & \cellcolor{cfTerra4}\textbf{+0.34} & \cellcolor{cfTerra2}\textbf{+0.07} \\
Qwen3-VL-32B & \cellcolor{cfTerra3}\textbf{+0.13} & \cellcolor{cfTerra3}\textbf{+0.12} & \cellcolor{cfTerra3}\textbf{+0.11} & \cellcolor{cfTerra3}\textbf{+0.23} & \cellcolor{cfTerra3}\textbf{+0.24} & \cellcolor{cfTerra2}\textbf{+0.07} \\
InternVL3.5-38B & \cellcolor{cfTerra1}\textbf{+0.01} & \cellcolor{cfTerra1}\textbf{+0.00} & \cellcolor{cfTerra1}\textbf{+0.01} & \cellcolor{cfTerra1}\textbf{+0.02} & \cellcolor{cfTerra1}\textbf{+0.01} & \cellcolor{cfTerra1}\textbf{+0.01} \\
Gemma 3 12B & \cellcolor{cfTerra3}\textbf{+0.19} & \cellcolor{cfTerra4}\textbf{+0.76} & \cellcolor{cfTerra4}\textbf{+0.48} & \cellcolor{cfTerra4}\textbf{+0.45} & \cellcolor{cfTerra3}\textbf{+0.18} & \cellcolor{cfTerra4}\textbf{+0.61} \\
MedGemma 4B & \cellcolor{cfTerra4}\textbf{+0.35} & \cellcolor{cfTerra2}\textbf{+0.03} & \cellcolor{cfTerra4}\textbf{+0.43} & \cellcolor{cfTerra3}\textbf{+0.19} & \cellcolor{cfTerra3}\textbf{+0.20} & \cellcolor{cfTerra3}\textbf{+0.14} \\
InternVL3.5-14B & \cellcolor{cfTerra1}\textbf{+0.02} & \cellcolor{cfTerra1}\textbf{+0.01} & \cellcolor{cfTerra1}\textbf{+0.01} & \cellcolor{cfTerra2}\textbf{+0.04} & \cellcolor{cfTerra2}\textbf{+0.03} & \cellcolor{cfTerra1}+0.00$^{\ddagger}$ \\
Qwen2.5-VL-32B & \cellcolor{cfTerra4}\textbf{+0.44} & \cellcolor{cfTerra3}\textbf{+0.12} & \cellcolor{cfTerra4}\textbf{+0.27} & \cellcolor{cfTerra4}\textbf{+0.83} & \cellcolor{cfTerra4}\textbf{+0.41} & \cellcolor{cfTerra3}\textbf{+0.19} \\
InternVL3.5-8B & \cellcolor{cfTerra2}\textbf{+0.03} & \cellcolor{cfTerra1}\textbf{+0.01} & \cellcolor{cfTerra2}\textbf{+0.04} & \cellcolor{cfTerra2}\textbf{+0.04} & \cellcolor{cfTerra2}\textbf{+0.04} & \cellcolor{cfTerra1}\textbf{+0.02} \\
HuatuoGPT-Vision-7B & \cellcolor{cfTerra1}\textbf{+0.02} & \cellcolor{cfTerra1}+0.00$^{\ddagger}$ & \cellcolor{cfTerra1}\textbf{+0.00} & \cellcolor{cfTerra2}\textbf{+0.03} & \cellcolor{cfTerra2}\textbf{+0.03} & \cellcolor{cfTerra1}+0.00$^{\ddagger}$ \\
LLaVA-Med 1.5 7B & \cellcolor{cfTerra1}+0.02 & \cellcolor{cfTerra1}\textbf{+0.01} & \cellcolor{cfTerra1}+0.01 & \cellcolor{cfSlate2}-0.02 & \cellcolor{cfTerra1}+0.00 & \cellcolor{cfSlate1}-0.01 \\
CheXagent-8B & \cellcolor{cfSlate1}-0.00 & \cellcolor{cfTerra1}+0.01 & \cellcolor{cfSlate1}-0.01 & \cellcolor{cfSlate2}-0.03 & \cellcolor{cfSlate3}-0.13 & \cellcolor{cfSlate2}-0.03 \\
Qwen2.5-VL-3B & \cellcolor{cfTerra4}\textbf{+0.46} & \cellcolor{cfTerra4}\textbf{+0.42} & \cellcolor{cfTerra4}\textbf{+0.53} & \cellcolor{cfTerra4}\textbf{+0.68} & \cellcolor{cfTerra4}\textbf{+0.61} & \cellcolor{cfTerra4}\textbf{+0.39} \\
Qwen2.5-VL-7B & \cellcolor{cfTerra4}\textbf{+0.48} & \cellcolor{cfTerra4}\textbf{+0.75} & \cellcolor{cfTerra4}\textbf{+0.61} & \cellcolor{cfTerra4}\textbf{+0.87} & \cellcolor{cfTerra4}\textbf{+0.51} & \cellcolor{cfTerra4}\textbf{+0.81} \\
Qwen3-VL-8B & \cellcolor{cfTerra4}\textbf{+0.42} & \cellcolor{cfTerra4}\textbf{+0.51} & \cellcolor{cfTerra4}\textbf{+0.74} & \cellcolor{cfTerra4}\textbf{+0.79} & \cellcolor{cfTerra4}\textbf{+0.59} & \cellcolor{cfTerra4}\textbf{+0.43} \\
Gemma 3 27B & \cellcolor{cfTerra3}\textbf{+0.24} & \cellcolor{cfTerra4}\textbf{+0.36} & \cellcolor{cfTerra2}\textbf{+0.09} & \cellcolor{cfTerra4}\textbf{+0.28} & \cellcolor{cfTerra3}\textbf{+0.15} & \cellcolor{cfTerra4}\textbf{+0.75} \\
Gemma 3 4B & \cellcolor{cfTerra3}\textbf{+0.16} & \cellcolor{cfTerra4}\textbf{+0.54} & \cellcolor{cfTerra3}\textbf{+0.24} & \cellcolor{cfTerra4}\textbf{+0.34} & \cellcolor{cfTerra4}\textbf{+0.44} & \cellcolor{cfSlate4}-0.27 \\
LLaVA-1.5-7B & \cellcolor{cfTerra2}\textbf{+0.06} & \cellcolor{cfTerra1}\textbf{+0.00} & \cellcolor{cfTerra1}\textbf{+0.01} & \cellcolor{cfTerra2}\textbf{+0.07} & \cellcolor{cfTerra1}+0.01 & \cellcolor{cfTerra1}\textbf{+0.00} \\
MedGemma 27B & \cellcolor{cfTerra4}\textbf{+0.33} & \cellcolor{cfSlate2}-0.04 & \cellcolor{cfTerra2}\textbf{+0.08} & \cellcolor{cfTerra3}\textbf{+0.18} & \cellcolor{cfTerra3}\textbf{+0.21} & \cellcolor{cfSlate2}-0.04 \\
LLaVA-1.5-13B & \cellcolor{cfTerra2}\textbf{+0.05} & \cellcolor{cfTerra1}+0.00$^{\ddagger}$ & \cellcolor{cfTerra1}+0.02 & \cellcolor{cfTerra2}\textbf{+0.08} & \cellcolor{cfTerra1}+0.02 & \cellcolor{cfSlate2}-0.03 \\
Llama 3.2 Vision 11B & \cellcolor{cfSlate2}-0.03 & \cellcolor{cfSlate2}-0.02 & \cellcolor{cfTerra2}+0.02 & \cellcolor{cfSlate1}-0.02 & \cellcolor{cfSlate2}-0.04 & \cellcolor{cfTerra1}\textbf{+0.01} \\
LLaVA-Rad-7B & \cellcolor{cfSlate2}-0.06 & \cellcolor{cfTerra1}+0.01 & \cellcolor{cfSlate2}-0.02 & \cellcolor{cfSlate2}-0.06 & \cellcolor{cfSlate2}-0.04 & \cellcolor{cfSlate1}-0.00 \\
\midrule
\textbf{Owned} & 21 & 19 & 19 & 20 & 18 & 16 \\
\bottomrule
\end{tabular}
\end{table}

\FloatBarrier

\subsection{Controls}
\label{app:cf-controls-sec}

Table~\ref{tab:cf-controls} summarises the reference family and the per-block controls of Section~\ref{sec:controls} across the grid.

\begin{table}[h]
\centering
\scriptsize
\setlength{\tabcolsep}{4pt}
\caption{\textbf{Reference family and controls.} Median and interquartile range per dataset: over all cells in the first four rows, over all blocks with the module in the next six. The last three rows count reference-meeting cells, owned cells, and reference-meeting cells without the simultaneous advantage.}
\label{tab:cf-controls}
\begin{tabular}{lrrrrrr}
\toprule
\multirow{2}{*}{Quantity} & \multicolumn{2}{c}{NIH ChestX-ray14} & \multicolumn{2}{c}{CheXpert Plus} & \multicolumn{2}{c}{COCO (control)}\\
\cmidrule(lr){2-3}\cmidrule(lr){4-5}\cmidrule(lr){6-7}
& \multicolumn{1}{c}{median} & \multicolumn{1}{c}{IQR} & \multicolumn{1}{c}{median} & \multicolumn{1}{c}{IQR} & \multicolumn{1}{c}{median} & \multicolumn{1}{c}{IQR}\\
\midrule
Random-family p95 of $W$ & +0.063 & [+0.021, +0.165] & +0.066 & [+0.027, +0.157] & +0.023 & [+0.008, +0.055] \\
$|$Sham write$|$ & +0.028 & [+0.009, +0.075] & +0.027 & [+0.008, +0.076] & +0.011 & [+0.003, +0.034] \\
Concept write $W_{qq}$ & +0.010 & [-0.006, +0.067] & +0.020 & [-0.001, +0.082] & +0.127 & [+0.020, +0.482] \\
Ownership $O_q$ & -0.017 & [-0.078, +0.000] & -0.014 & [-0.065, +0.001] & +0.102 & [+0.009, +0.433] \\
Signed dose range of $O_q$ & +0.266 & [+0.061, +0.444] & +0.162 & [+0.034, +0.359] & +0.399 & [+0.050, +0.612] \\
Refit SD of $O_q$ (3 seeds) & +0.039 & [+0.015, +0.072] & +0.037 & [+0.017, +0.052] & +0.043 & [+0.006, +0.074] \\
Connector-locus median $O_q$ & -0.013 & [-0.023, -0.005] & -0.009 & [-0.020, -0.004] & +0.003 & [+0.001, +0.007] \\
Label gap (logits) & -0.1 & [-0.9, -0.0] & -0.3 & [-1.3, +0.0] & -5.3 & [-14.8, -0.3] \\
$|$Wording contrast$|$ (IY$-$WY) & +0.033 & [+0.008, +0.078] & +0.021 & [+0.007, +0.035] & +0.028 & [+0.007, +0.065] \\
$|$Mapping contrast$|$ (IA$-$IB) & +0.037 & [+0.008, +0.073] & +0.019 & [+0.003, +0.055] & +0.030 & [+0.007, +0.051] \\
\midrule
Reference met & \multicolumn{2}{c}{20 of 150} & \multicolumn{2}{c}{36 of 150} & \multicolumn{2}{c}{117 of 150} \\
Owned & \multicolumn{2}{c}{13 of 150} & \multicolumn{2}{c}{25 of 150} & \multicolumn{2}{c}{113 of 150} \\
Reference met but not owned & \multicolumn{2}{c}{7 of 150} & \multicolumn{2}{c}{11 of 150} & \multicolumn{2}{c}{4 of 150} \\
\bottomrule
\end{tabular}
\end{table}

\subsection{What each analysis scores its writes against}
\label{app:cf-ledger}

Table~\ref{tab:cf-ledger} is the ledger of references and rules: one row per graded analysis, with the rows it scores, the reference its steering test applies, the rule that decides its verdict, the direction its sham permutes, and the number of cells it grades. It covers \cfLedgerAnalyses{} analyses and \cfLedgerCells{} cells, and it is populated from the packaged per-block statistics, which also check each declared reference and sham before the row is written.

The ledger states every departure from the write matrix's own combination of reference, rule and sham. \cfLedgerFullReference{} of the \cfLedgerAnalyses{} analyses reference the own write against the 95th percentile of the 119-direction random family and the absolute sham. The semantic endpoints reference it against the maximum of a \cfSemLedgerRandom{}-direction family of their own and the endpoint's sham, over \cfSemLedgerCells{} cells. The answer direction under the held-out templates has the random 95th percentile for \cfAnsdirtRefPairs{} of its \cfAnsdirtLedgerPairs{} (cell, template) pairs, and the sham alone for the other \cfAnsdirtShamPairs{}. \cfLedgerMaxT{} of the \cfLedgerAnalyses{} decide the verdict by simultaneous max-$T$ bounds on the fixed contrasts between the own write and each competitor; the logit-scale grade decides by a percentile interval on the ownership contrast, whose maximum it recomputes in every draw, over its \cfLedgerPercentileCells{} cells. \cfLedgerOwnSham{} of the \cfLedgerAnalyses{} compare against a coordinate permutation of the direction they write; \cfLedgerSeedShamWord{}---the probe refits, the alternative direction constructions, the displacement lifts, and the expert-label refits, \cfLedgerSeedShamCells{} cells between them---compare against the permutation of the seed-0 logistic normal, which is not the direction they write. The token-weighted writes compare against the uniformly written permutation of the normal they weight.

\begin{table}[htbp]
\centering
\scriptsize
\setlength{\tabcolsep}{3pt}
\caption{\textbf{What each analysis scores its writes against.} One row per graded analysis: the cohort it scores and its size, the steering test its own write must pass, the rule that decides its verdict, the direction its sham permutes, and the number of graded (question, arm) cells.}
\label{tab:cf-ledger}
\begin{tabular}{>{\raggedright\arraybackslash}p{0.155\textwidth}>{\raggedright\arraybackslash}p{0.088\textwidth}>{\raggedright\arraybackslash}p{0.215\textwidth}>{\raggedright\arraybackslash}p{0.175\textwidth}>{\raggedright\arraybackslash}p{0.125\textwidth}r}
\toprule
Analysis & Rows & Reference & Rule & Sham & \multicolumn{1}{c}{Cells}\\
\midrule
Write matrix, primary template & 600 test & 119 random $p_{95}$, sham & simultaneous max-$T$, 30 contrasts & the direction written & 450 \\
Probe refits, seeds 1 and 2 & 600 test & 119 random $p_{95}$, sham, the write matrix's & simultaneous max-$T$, 30 contrasts & the seed-0 normal & 900 \\
Logit-scale grade & 600 test & 119 random $p_{95}$, sham, on the margin scale & percentile interval on $O^m_q$ & the direction written & 450 \\
Unsaturated-row grade & unsaturated test & 119 random $p_{95}$, sham, on those rows & simultaneous max-$T$, 30 contrasts & the direction written & 450 \\
Alternative direction constructions & 600 test & 119 random $p_{95}$, sham, the write matrix's & simultaneous max-$T$, 30 contrasts & the seed-0 normal & 1{,}440 \\
Displacement lifts & 600 test & 119 random $p_{95}$, sham, the write matrix's & simultaneous max-$T$, 30 contrasts & the seed-0 normal & 684 \\
Extended competitor family & 600 test & 119 random $p_{95}$, sham & simultaneous max-$T$, 8 or 11 contrasts & the direction written & 54 \\
Token-weighted writes & 600 test & 119 random $p_{95}$, sham, the sham unweighted & simultaneous max-$T$, 30 contrasts & the direction written & 96 \\
fp32 and batch-size-one rescoring & 200 test & 119 random $p_{95}$, sham & simultaneous max-$T$, 30 contrasts & the direction written & 108 \\
Answer direction, primary template & 600 test & 119 random $p_{95}$, sham & simultaneous max-$T$, 30 contrasts & the direction written & 54 \\
Answer direction, held-out templates & 600 test & 119 random $p_{95}$, sham in 80 pairs; sham alone in 190 & simultaneous max-$T$, 30 contrasts & the direction written & 270 \\
Non-clinical attributes & 600 test & 119 random $p_{95}$, sham & simultaneous max-$T$, 72 contrasts & the direction written & 114 \\
Findings in the attribute family & 600 test & 119 random $p_{95}$, sham & simultaneous max-$T$, 72 contrasts & the direction written & 228 \\
Radiologist-label regrade & 200 valid & 119 random $p_{95}$, sham & simultaneous max-$T$, 30 contrasts & the direction written & 114 \\
Expert-label refits & 600 test & 119 random $p_{95}$, sham, the write matrix's & simultaneous max-$T$, 30 contrasts & the seed-0 normal & 36 \\
Projection seeds 1 and 2 & 600 test & 119 random $p_{95}$, sham, that projection's & simultaneous max-$T$, 30 contrasts & the direction written & 36 \\
Crossed tower and reader & 200 test & 119 random $p_{95}$, sham, of the tower in use & simultaneous max-$T$, 30 contrasts & the direction written & 132 \\
Identical-tensor replay & 200 test & 119 random $p_{95}$, sham, the source block's & simultaneous max-$T$, 30 contrasts & the direction written & 48 \\
Semantic endpoints & 600 test & 18 random, maximum; sham, the endpoint's own & simultaneous max-$T$, 30 contrasts & the direction written & 144 \\
Fine-grained object family & 600 test & 119 random $p_{95}$, sham, the family's own & simultaneous max-$T$, 30 contrasts & the direction written & 36 \\
\bottomrule
\end{tabular}
\end{table}


\subsection{Robustness of the deficit}
\label{app:cf-robustness}

The projected features are $x=R^{\top}\hvec$, where $R\in\R^{D\times512}$ is the fixed Gaussian projection. Two lifts map vectors from the projected, train-scaled space back to the consumed space. A classifier coefficient $\boldsymbol\beta$ lifts as $\hat\wvec\propto R\,\mathrm{diag}(1/\boldsymbol s)\boldsymbol\beta$. The resulting written direction scores the concept as the probe does (Equation~\ref{eq:lift}). A displacement $\boldsymbol u$ lifts via its minimum-norm preimage as $\hat\wvec\propto R(R^{\top}R)^{-1}\mathrm{diag}(\boldsymbol s)\boldsymbol u$. This lift reproduces the displacement exactly: $R^{\top}\hat\wvec\propto\mathrm{diag}(\boldsymbol s)\boldsymbol u$. The lifts differ because $R^{\top}R$ is not a multiple of the identity. Measurements on every block with the displacement lifts show that $R^{\top}R/D$ is full rank, with eigenvalues $\cfAltdirdEigMin$--$\cfAltdirdEigMax$ and condition number \cfAltdirdCondMin{}--\cfAltdirdCondMax{}. We compute the difference-of-means and pattern vectors in the same projected, train-scaled space as the logistic normal and report both lifts. The alternative-construction module scores the coefficient-lift family. The displacement lift gives the geometrically matched comparison for these two families: a difference of class means and a Haufe pattern are displacements in the projected space rather than classifier coefficients. We form the orthogonalised normal by removing the span of the other five normals from the logistic normal. We refit the residualised normal on features residualised against the other five labels.

\paragraph{Alternative directions, weights, and competitor families.} Token-weighted variants multiply the per-token dose $\alpha\lVert\hvec_{it}\rVert$ by a weight $w_t$ whose mean over the consumed tokens is one. The softmax variant weights tokens by the softmax of their probe scores. It preserves the uniform write's L1 dose. The top-quarter variant ($w_t=4$ on the 25\% highest-scoring tokens) has twice the uniform write's Frobenius norm. The extended competitor family adds a direction for every other dataset label with at least 100 known positives and 100 known negatives (NIH: Consolidation, Edema, Infiltration; CheXpert: Enlarged Cardiomediastinum, Fracture, Lung Lesion, Lung Opacity, Pneumonia, Support Devices). We fit and write each added direction like the protocol normals. Under this family, a cell is owned when its write meets the steering reference and holds the simultaneous advantage over the five protocol competitors and every extra direction. Table~\ref{tab:cf-altdir} reports every construction, lift, token weighting, and competitor family per dataset. Each panel includes the protocol's logistic normal on the same blocks.

\paragraph{Non-clinical attributes.} The attribute directions use three labels from each dataset's metadata: the view field (anteroposterior or portable against posteroanterior), recorded sex, and age at least 60. We fit these directions using the same images, features, projection, and probe settings as the clinical probes. We then lift and write them with the six clinical directions as one nine-direction family. Attribute and clinical questions share the same steering reference in all \cfAttrRefBlocks{} blocks of Table~\ref{tab:cf-attr}: the 119-direction random family written at the same dose and the coordinate-permutation sham. We score each attribute question in \cfAttrqPhrasings{} phrasings on the calibration rows. The phrasing used for the write reaches a median clean-answer AUROC of \cfAttrqWrittenAuroc{}, compared with \cfAttrqSelectedAuroc{} for the phrasing in the same cell with the highest one-sided lower bound. Across the \cfAttrqCells{} attribute cells, the median gap is \cfAttrqGain{} and the largest gap is \cfAttrqGainMax{}. \cfAttrqCapable{} of those cells clear the answer-capable rule as written. We assess the comparison three ways: over every cell, over answer-capable cells alone, and over attribute cells paired with clinical cells in the same block within \cfAttrPairWindow{} clean-answer AUROC. This matching places \cfAttrPairMatched{} of the \cfAttrPairCellsN{} attribute cells into \cfAttrPairAttrN{} pairs and leaves \cfAttrPairUnmatched{} without a partner in their own block. Attribute probes reach a median AUROC of \cfAttrAurocMin{}--\cfAttrAurocMax{} against the attribute label and \cfAttrAnswerAurocMin{}--\cfAttrAnswerAurocMax{} against the model's own clean answer. We summarise each attribute--finding pair by its median over blocks; over the \cfAttrCosPairs{} pairs the largest absolute raw model-space cosine between an attribute direction and a clinical normal is \cfAttrCosMax{}, for \cfAttrCosPairLabel{}. Across the \cfAttrBlocks{} blocks with the module, ownership counts are \cfAttrChestAttrOwned{} of \cfAttrChestAttrN{} cells for attributes and \cfAttrChestClinOwned{} of \cfAttrChestClinN{} for findings in the same family. The corresponding attribute and finding counts are \cfAttrChexAttrOwned{}/\cfAttrChexAttrN{} and \cfAttrChexClinOwned{}/\cfAttrChexClinN{} on CheXpert, and \cfAttrNihAttrOwned{}/\cfAttrNihAttrN{} and \cfAttrNihClinOwned{}/\cfAttrNihClinN{} on NIH.

\paragraph{Answer direction: rescue, loss, and templates.} We lift the answer direction $a_q$ of Section~\ref{sec:ansdir} as we lift the probe normal. We write it at the same dose, against the write matrix's own 119-direction random family and a coordinate-permutation sham of the $a_q$ under test, using the other five $a_d$ as clinical competitors. A concept is \emph{kept} when both $a_q$ and the label direction $\hat w_q$ are owned, \emph{rescued} when only $a_q$ is owned, and \emph{lost} when only $\hat w_q$ is owned (Table~\ref{tab:cf-ansdir}, top). The release reports both ownership contrasts, column selectivity, and the whitened cosine for every cell. We also write the $a_q$ fitted on the primary template IY unchanged under the rewording WY and the answer-mapped templates IA, IB, WA, and WB. Each template uses its own clean baseline, the $a_d$ of the other five concepts, and a coordinate-permutation sham of the $a_q$ under test. The held-out templates score no random family of their own. The random 95th percentile is available where the additional-template module scored one under that template, which is \cfAnsdirtRefPairs{} of the \cfAnsdirtLedgerPairs{} (cell, template) pairs; the remaining \cfAnsdirtShamPairs{} pairs meet the steering reference against their own sham alone. Every pair takes the same $6\times5$ max-$T$ verdict (Table~\ref{tab:cf-ledger}). A block enters when every held-out template is scored on all 600 test rows (Table~\ref{tab:cf-ansdir}, bottom).

\paragraph{Direction and label geometry.} Neither correlations among the six directions nor co-occurrence of the six report labels explains why a competing direction wins (Table~\ref{tab:cf-geometry}). The direction cosine is the raw model-space cosine between two of a block's six write directions, averaged over all blocks' off-diagonal pairs. The label $\phi$ measures the correlation between two of the six test labels. On NIH, the six directions are close to orthogonal: their mean absolute off-diagonal cosine is $\cfGeoAbsCosNih$, against $\cfGeoCosRandom$ between pairs of random unit directions, and their signed mean is $\cfGeoCosNih$. The labels are nearly uncorrelated (mean $\phi=\cfGeoPhiNih$). Yet NIH shows the same off-diagonal dominance as CheXpert, where directions and labels are correlated (mean absolute off-diagonal cosine $\cfGeoAbsCosChex$, signed mean $\cfGeoCosChex$; Effusion--Consolidation $\cfGeoCosEffCons$ averaged over blocks, $\phi=\cfGeoPhiEffCons$). The ownership contrast is positive, $O_q>0$, in \cfGeoOwnWinsNihPct\% of the \cfNihCells{} NIH cells and \cfGeoOwnWinsChexPct\% of the \cfChexCells{} CheXpert cells. Each cell has five competing directions, so its strongest competitor matches any fixed competitor in one cell in five by chance. The strongest competitor matches the most similar direction at this rate (\cfGeoChanceCosMin{}--\cfGeoChanceCosMax\% of chest cells per dataset). It also corresponds to the label with the largest $\phi$ or the most co-occurring label at this rate (\cfGeoChanceLabelMin{}--\cfGeoChanceLabelMax\%). Across the \cfGeoOffdiagCells{} off-diagonal chest cells, the competitor's advantage $W_{q,d}-W_{q,q}$ is uncorrelated with the direction cosine (Spearman over cells $\cfGeoRhoAdvCos$, $p=\cfGeoRhoAdvCosP$) and with the label $\phi$ ($\cfGeoRhoAdvPhi$). At least two competitors exceed the own write in \cfGeoTwoAbove{} of \cfGeoTwoAboveN{} chest cells (\cfGeoTwoAbovePct\%). Removing the strongest competitor therefore leaves $O_q$ negative in those cells. The last three rows of Table~\ref{tab:cf-geometry} count the sign of $O_q$ after removal, not the owned grade. On COCO, we restrict each cell's competitors to a co-occurring family of objects (person with bicycle, $\phi=\cfGeoFamPhiPersonBicycle$; person with dog, $\phi=\cfGeoFamPhiPersonDog$; chair with bottle, $\phi=\cfGeoFamPhiChairBottle$; car with bicycle, $\phi=\cfGeoFamPhiCarBicycle$; person with chair and bottle, mean $\phi=\cfGeoFamPhiPersonChairBottle$). This leaves $O_q>0$ in \cfGeoCocoFamMinPct{}--\cfGeoCocoFamMaxPct\% of each family's cells, compared with \cfGeoCocoFamFullMinPct{}--\cfGeoCocoFamFullMaxPct\% of the same cells with all five competitors. Of the five directions in the seed cell, Nodule is the \emph{least} similar to Effusion (raw model-space cosine $\cfSeedCosNodule$ in this block). It still moves the Effusion answer more than Effusion's own direction ($\cfSeedWNodule$ against $\cfSeedWEff$). The most similar direction, Atelectasis (raw model-space cosine $\cfSeedCosAtel$), moves it by $\cfSeedWAtel$.

\paragraph{Scale, saturation, refits, and numerics.} The grades hold on the logit scale. Recomputing the grade from the margin $\ell_{\mathrm{yes}}-\ell_{\mathrm{no}}$ reproduces the probability-scale owned grade in \cfScaleAgree{} of \cfScaleAgreeN{} cells. Every owned cell retains a positive ownership contrast on that scale, $O^m_q>0$. The full owned grade persists in \cfScaleOwnedKeepFull{} of the \cfScaleOwned{} owned cells, with the steering reference recomputed on the margin scale and a positive percentile interval for $O^m_q$ (Table~\ref{tab:cf-scale}). Saturation counters a ceiling explanation: COCO blocks contain far more rows at $P(\mathrm{yes})>0.99$ or $<0.01$ than chest blocks (median block share $\cfScaleSatCoco$ against $\cfScaleSatNih$ and $\cfScaleSatChex$). Restricting each cell to its unsaturated rows preserves $O_q>0$ in all \cfScaleUnsatOwned{} owned chest cells and in \cfScaleUnsatOtherPos{} of the \cfScaleUnsatOther{} others. Full regrading uses a steering reference recomputed on these rows and the simultaneous max-$T$ advantage over the same six-direction family. It preserves the owned grade in \cfScaleUnsatGradeKept{} of \cfScaleUnsatGradeN{} owned chest cells and raises \cfScaleUnsatGradeGained{} of the \cfScaleUnsatGradeOtherN{} others to owned. The same pattern holds in the clean-No, clean-Yes, label-present, and label-absent strata. Refitting the probes on resampled training patients, and grading each refitted direction against the seed-0 random 95th percentile and the seed-0 sham of the same block and question, changes $O_q$ by a median of $\cfScaleRefitMedian$ (Table~\ref{tab:cf-ledger}); the last column of Table~\ref{tab:cf-refit} regrades seed 0 with the same draws and counts agreement with the paper's grade. Rescoring the write grid on the first 200 test rows of each block, in fp32 and at batch size one, moves the $6\times126$ written cells of a block by at most \cfPrecisionMaxDW{}. Table~\ref{tab:cf-scale} gives that largest change and the largest in a clinical contrast $C_{q,d}$, the difference between the own write and a competing write on the same question. Regraded under either setting, verdicts change in \cfPrecisionVerdictChanges{} and steering references in \cfPrecisionRefChanges{} of the \cfPrecisionGradeCells{} grades, and the two point verdicts change in \cfPrecisionPointVerdictChanges{} (Table~\ref{tab:cf-scale}). The bf16 batch effects declared at the gates (candidate-logit differences of $\cfScaleBatchMin$--$\cfScaleBatchMax$) affect no owned cell. The margin sign agrees in \cfScaleBatchAgree{} of \cfScaleBatchN{} blocks; the \cfScaleBatchDisagree{} disagreeing blocks contain no owned cells. Re-scoring the clean baseline in a second module gives bitwise agreement in 100\% of rows in \cfScaleBaselineIdentical{} of \cfScaleBaselineN{} blocks and in at least \cfScaleBaselineRestPct\% of rows in the rest. Differences satisfying $|\Delta P(\mathrm{yes})|>0.1$ occur in at most \cfScaleBaselineDpPct\% of rows. Writes at the connector locus are an order of magnitude smaller than those at the primary locus (median $|W_{q,q}|$ $\cfConnWqqNih$ against $\cfPrimWqqNih$ on NIH and $\cfConnWqqCoco$ against $\cfPrimWqqCoco$ on COCO). The readable and answerable grades use the 400 calibration rows for support. \cfSupportInelChest{} of the \cfChestCellsN{} chest cells fail that cohort's requirement of ten known positives and ten known negatives: CheXpert Atelectasis in each of the \cfSupportInelChexBlocks{} CheXpert blocks, with \cfCalChexInelPos{} known positives and \cfCalChexInelNeg{} known negatives among the 400 calibration rows. Every NIH concept meets the requirement on that cohort; Pneumothorax has the smallest count, at \cfCalNihMinPos{} positives. On COCO, bicycle is ineligible in each of the \cfSupportInelCocoBlocks{} COCO blocks, with \cfCalCocoInelPos{} positives. Support counts for these two grades come solely from calibration. The test cohort carrying the writes has its own counts (\cfTestChexInelPos{} positives and \cfTestChexInelNeg{} negatives for CheXpert Atelectasis, \cfTestCocoInelPos{} positives for COCO bicycle). Among eligible cells, \cfSupportReadableChest{} of \cfSupportReadableNChest{} are readable (\cfSupportReadablePctChest\%) and \cfSupportAnswerableChest{} of \cfSupportAnswerableNChest{} are answerable (\cfSupportAnswerablePctChest\%).

\paragraph{The answer direction and the label source.} We fit the answer direction of Section~\ref{sec:ansdir} to one quantity: the model's own clean answer margin $\ell_{\mathrm{yes}}-\ell_{\mathrm{no}}$ on the block's primary template, read from clean forward passes with the write disabled. We use the block's own projected, train-scaled features from the first 3{,}000 training-cohort rows in the cohort's fixed manifest order. A seeded shuffle determines this order on NIH, and a hash of patient or image identity determines it on CheXpert and COCO. Neither labels nor answers enter the ordering. For each question, we fit the margin from these features by ridge regression. We select the penalty by five-fold cross-validated $R^2$ over $\{0.1,1,10,100\}$, then refit on all 3{,}000 rows with the chosen penalty. Across the \cfAnsAlphaCells{} fitted questions, the selected penalty is $1$ for \cfAnsAlphaOne{}, $10$ for \cfAnsAlphaTen{}, and $100$ for \cfAnsAlphaHundred{}. The folds split rows: one per patient on CheXpert, one per image on COCO, and rows from \cfAnsTrainUnitsNih{} patients on NIH. We lift and normalise the fitted coefficient vector exactly as we do a probe normal. The radiologist-label cohort provides two distinct tests. First, we re-grade the unchanged directions fitted on report-derived labels on \cfValidRows{} films labelled by radiologist consensus. These patients appear in no campaign cohort (Table~\ref{tab:cf-valid}). The absolute difference $|O_q(\text{valid})-O_q(\text{test})|$ has a 90th percentile of \cfValidNinetyDO{} over all cells. On the \cfValidSupported{} cells with at least ten known positives and ten known negatives on both cohorts, readability agrees in \cfValidReadableAgree{} of \cfValidReadableN{} and answer capability in \cfValidAnswerAgree{} of \cfValidAnswerN{}. Second, we refit the six directions on those radiologist labels using the block's own projection, its train-only scaler, and the protocol's probe settings. We cross-fit over patients in \cfValidfitFolds{} folds, scoring every film with a direction not fitted on it. We write the refitted directions on the 600 test rows at the primary dose and template, against the write matrix's own 119-direction random family and its seed-0 sham, and take the same $6\times5$ max-$T$ verdict; the sham here is the permutation of the report-label direction rather than of the refitted one (Table~\ref{tab:cf-ledger}). Across the \cfValidfitBlocks{} CheXpert blocks with this refit, the expert-label and report-label directions have a median raw model-space cosine of \cfValidfitCosRaw{} and a median covariance-whitened cosine of \cfValidfitCosWhitened{}. The expert-label refit achieves a median held-out AUROC of \cfValidfitAurocExpert{} on the radiologist labels, against \cfValidfitAurocReportDir{} for the report-label direction on the same films. Those two directions differ in two ways at once: the labels they were fitted to, and the number of rows available to fit them (\cfVfArmExpertRows{} per fold against a median of \cfVfArmReportFullRows{} labelled training rows per concept). Three further arms separate the two, all scored on the same films, with the block's own projection, its stored train-only scaler, the protocol's probe settings and the same \cfVfArmFolds{}-fold cross-fitting over patients (Table~\ref{tab:cf-valid}, lower panel). Report-derived labels of those same \cfVfArmRows{} films reach \cfVfArmReportSameFilmsAuroc{}; report-derived labels of a training subsample drawn to leave each concept the same \cfVfArmReportSubRows{} fitting rows as the refit, redrawn \cfVfArmDraws{} times, reach \cfVfArmReportSubAuroc{}; report-derived labels of every training row reach \cfVfArmReportFullAuroc{}. Holding the number of fitting rows fixed, the label source moves the result by $\cfVfArmDeltaLabel$; holding the label source fixed, the sample size moves it by $\cfVfArmDeltaSize$. The size-matched report refit also sits at a raw model-space cosine of \cfVfArmReportSubCos{} to the shipped direction and a whitened cosine of \cfVfArmReportSubCosWhite{}, against \cfVfArmExpertCos{} and \cfVfArmExpertCosWhite{} for the expert refit. The distance between the refit and the shipped direction is the estimation sample, not the label source. The expert-label directions are owned in \cfValidfitOwnedExpert{} of \cfValidfitCells{} cells, against \cfValidfitOwnedReport{} for the report-label directions of the same blocks.

Table~\ref{tab:cf-geometry} reports direction and label geometry. Table~\ref{tab:cf-scale} reports scale, saturation, refit, connector, batch, and numerics statistics supporting Section~\ref{sec:robust-geometry}. Table~\ref{tab:cf-refit} reports the full grade under refitted probes. Table~\ref{tab:cf-pairs} reports paired shared-tower differences supporting Section~\ref{sec:reader}. Table~\ref{tab:cf-altdir} reports alternative directions, lifts, token weights, and competitor families. Table~\ref{tab:cf-ansdir} reports the answer direction on the primary and held-out templates. Table~\ref{tab:cf-attr} reports non-clinical attribute directions. Table~\ref{tab:cf-validation} reports what the answer direction reads and the column selectivity of $W$. Table~\ref{tab:cf-valid} reports the grade on radiologist labels and separates the label source of a refitted direction from the number of rows it was fitted on. Table~\ref{tab:cf-towerswap} reports the crossed tower and reader. Table~\ref{tab:cf-semend} reports the four semantic endpoints. Table~\ref{tab:cf-fgobj} reports the fine-grained object family with the difficulty-matched comparison (Appendix~\ref{app:cf-crossed}).

\begin{table}[h]
\centering
\scriptsize
\setlength{\tabcolsep}{5pt}
\caption{\textbf{Direction and label geometry against off-diagonal dominance.} One column per dataset, with the two chest sets pooled in the third. A ``--'' marks a row that is defined per dataset only.}
\label{tab:cf-geometry}
\begin{tabular}{lrrrr}
\toprule
& \multicolumn{1}{c}{NIH ChestX-ray14} & \multicolumn{1}{c}{CheXpert Plus} & \multicolumn{1}{c}{both chest sets} & \multicolumn{1}{c}{COCO}\\
\midrule
\multicolumn{5}{l}{\textit{Geometry of the six directions and the six labels}} \\
\quad mean off-diagonal cosine between the directions & 0.019 & 0.215 & -- & 0.043 \\
\quad mean off-diagonal $|\cos|$ between the directions & 0.093 & 0.218 & -- & 0.058 \\
\quad mean $|\cos|$ between random unit directions & 0.024 & 0.024 & -- & 0.024 \\
\quad mean off-diagonal $\phi$ between the labels & 0.026 & 0.381 & -- & 0.022 \\
\midrule
\multicolumn{5}{l}{\textit{Cells whose strongest competitor is (chance: one in five)}} \\
\quad the most similar direction & 31/150 & 33/150 & 64/300 & 31/150 \\
\quad the label with the largest $\phi$ & 35/150 & 36/150 & 71/300 & 29/150 \\
\quad the most co-occurring label & 24/150 & 32/150 & 56/300 & 22/150 \\
\midrule
\multicolumn{5}{l}{\textit{Spearman $\rho$ over off-diagonal cells of $W_{q,d}-W_{q,q}$ with}} \\
\quad the cosine between the two directions & 0.085 & 0.048 & 0.037 & 0.042 \\
\quad the $\phi$ between the two labels & 0.036 & 0.037 & -0.001 & 0.007 \\
\midrule
\multicolumn{5}{l}{\textit{Sign of $O_q$}} \\
\quad cells with $O_q>0$ & 38/150 & 43/150 & 81/300 & 129/150 \\
\quad cells with $O_q>0$ without the strongest competitor & 62/150 & 75/150 & 137/300 & 136/150 \\
\quad cells with at least two competitors above the own write & 88/150 & 75/150 & 163/300 & 14/150 \\
\bottomrule
\end{tabular}
\end{table}
\begin{table}[h]
\centering
\scriptsize
\setlength{\tabcolsep}{4pt}
\caption{\textbf{Ownership under alternative scales, strata, refits, loci, and numerics.} Per dataset, cells with the complete owned grade under each alternative, beside the paper's grade on the probability ($p$) and logit-margin ($m$) scales. The lower panel rescores the write grid in fp32 and at batch size one.}
\label{tab:cf-scale}
\begin{tabular}{lrrrrrrrrrr}
\toprule
\multirow{2}{*}{Dataset} & \multicolumn{3}{c}{owned cells} & \multicolumn{2}{c}{unsaturated, $O_q>0$} & \multicolumn{2}{c}{refits} & \multicolumn{2}{c}{median $|W_{q,q}|$} & \multicolumn{1}{c}{connector}\\
\cmidrule(lr){2-4}\cmidrule(lr){5-6}\cmidrule(lr){7-8}\cmidrule(lr){9-10}
& \multicolumn{1}{c}{$p$} & \multicolumn{1}{c}{$m$} & \multicolumn{1}{c}{both} & \multicolumn{1}{c}{owned} & \multicolumn{1}{c}{other} & \multicolumn{1}{c}{med.\ $|\Delta O|$} & \multicolumn{1}{c}{stable} & \multicolumn{1}{c}{primary} & \multicolumn{1}{c}{connector} & \multicolumn{1}{c}{$>$ p95}\\
\midrule
NIH ChestX-ray14 & 13/150 & 12/150 & 11 & 13/13 & 25/131 & 0.020 & 8/13 & 0.032 & 0.008 & 15/150 \\
CheXpert Plus & 25/150 & 27/150 & 25 & 25/25 & 16/121 & 0.023 & 10/25 & 0.032 & 0.010 & 16/144 \\
COCO & 113/150 & 116/150 & 112 & 105/106 & 16/37 & 0.020 & 107/113 & 0.127 & 0.008 & 106/150 \\
\midrule
\multicolumn{11}{p{0.9\linewidth}}{\textit{Batch effects at the gates (61 blocks): batched-vs-single candidate-logit difference 0.00--1.84 (median 0.28); margin sign agreement in 55 of 61 blocks; owned cells in disagreeing blocks: 0.}} \\
\bottomrule
\end{tabular}

\medskip
\setlength{\tabcolsep}{5pt}
\begin{tabular}{llrrrr}
\toprule
\multicolumn{6}{l}{\textit{Numerics: the same grid rescored in fp32 and at batch size one}} \\
\midrule
 & & \multicolumn{2}{c}{fp32 weights and forward} & \multicolumn{2}{c}{bf16, batch size one} \\
\cmidrule(lr){3-4}\cmidrule(lr){5-6}
Checkpoint & Dataset & $\max|\Delta W|$ & $\max|\Delta C|$ & $\max|\Delta W|$ & $\max|\Delta C|$ \\
\midrule
Qwen2.5-VL-7B & NIH ChestX-ray14 & 0.0087 & 0.0063 & 0.0065 & 0.0049 \\
Qwen2.5-VL-7B & COCO & 0.0060 & 0.0032 & 0.0028 & 0.0020 \\
Qwen3-VL-8B & NIH ChestX-ray14 & 0.0177 & 0.0089 & 0.0204 & 0.0089 \\
InternVL3.5-8B & NIH ChestX-ray14 & 0.0018 & 0.0014 & 0.0018 & 0.0014 \\
Gemma 3 4B & NIH ChestX-ray14 & 0.0544 & 0.0493 & 0.0261 & 0.0377 \\
Gemma 3 12B & NIH ChestX-ray14 & 0.0558 & 0.0359 & 0.0111 & 0.0109 \\
Gemma 3 27B & NIH ChestX-ray14 & 0.0413 & 0.0186 & 0.0149 & 0.0053 \\
MedGemma 4B & NIH ChestX-ray14 & 0.0125 & 0.0069 & 0.0056 & 0.0054 \\
Lingshu 32B & NIH ChestX-ray14 & 0.0032 & 0.0023 & 0.0009 & 0.0007 \\
\bottomrule
\end{tabular}
\end{table}
\begin{table}[h]
\centering
\scriptsize
\setlength{\tabcolsep}{3pt}
\caption{\textbf{The full grade under refitted probes.} Per dataset: cells whose seed-0 grade is owned, a stronger competitor, unresolved, or a fixed-family advantage without the reference, and how many keep that grade under both refits, under at least one, or under none.}
\label{tab:cf-refit}
\begin{tabular}{lrrrrrrrrrrrrrrr}
\toprule
\multirow{2}{*}{Dataset} & \multirow{2}{*}{blocks} & \multirow{2}{*}{cells} & \multicolumn{4}{c}{owned} & \multicolumn{3}{c}{stronger competitor} & \multicolumn{3}{c}{unresolved} & \multicolumn{2}{c}{adv.\ no ref.} & \multicolumn{1}{c}{seed 0}\\
\cmidrule(lr){4-7}\cmidrule(lr){8-10}\cmidrule(lr){11-13}\cmidrule(lr){14-15}
 &  &  & \multicolumn{1}{c}{$n$} & \multicolumn{1}{c}{both} & \multicolumn{1}{c}{$\geq1$} & \multicolumn{1}{c}{none} & \multicolumn{1}{c}{$n$} & \multicolumn{1}{c}{both} & \multicolumn{1}{c}{$\geq1$} & \multicolumn{1}{c}{$n$} & \multicolumn{1}{c}{both} & \multicolumn{1}{c}{$\geq1$} & \multicolumn{1}{c}{$n$} & \multicolumn{1}{c}{both} & \multicolumn{1}{c}{agrees}\\
\midrule
NIH ChestX-ray14 & 25 & 150 & 13 & 4 & 8 & 5 & 101 & 69 & 97 & 21 & 5 & 12 & 15 & 0 & 150/150 \\
CheXpert Plus & 25 & 150 & 25 & 6 & 17 & 8 & 95 & 63 & 89 & 21 & 3 & 10 & 9 & 3 & 150/150 \\
COCO & 25 & 150 & 113 & 99 & 110 & 3 & 17 & 13 & 17 & 10 & 4 & 7 & 10 & 2 & 149/150 \\
\textit{chest} & 50 & 300 & 38 & 10 & 25 & 13 & 196 & 132 & 186 & 42 & 8 & 22 & 24 & 3 & 300/300 \\
\textit{all} & 75 & 450 & 151 & 109 & 135 & 16 & 213 & 145 & 203 & 52 & 12 & 29 & 34 & 5 & 449/450 \\
\bottomrule
\end{tabular}
\end{table}
\begin{table}[h]
\centering
\scriptsize
\setlength{\tabcolsep}{4pt}
\caption{\textbf{Same write, different reader: paired differences.} For every pair of checkpoints that share a vision tower, the difference in ownership $\Delta O_q=O_q(\text{larger})-O_q(\text{smaller})$ on the same 600 test rows, with a paired patient-bootstrap 95\% interval. The last two columns divide $|\Delta O_q|$ by the two blocks' mean refit standard deviation.}
\label{tab:cf-pairs}
\begin{tabular}{llcrrrrr}
\toprule
\multirow{2}{*}{Pair} & \multirow{2}{*}{Dataset} & \multirow{2}{*}{relation} & \multicolumn{1}{c}{$\Delta O_q$ CI} & \multicolumn{2}{c}{$|\Delta O_q|$} & \multicolumn{2}{c}{ratio to refit SD}\\
\cmidrule(lr){5-6}\cmidrule(lr){7-8}
&  &  & \multicolumn{1}{c}{$\not\ni 0$} & \multicolumn{1}{c}{median} & \multicolumn{1}{c}{max} & \multicolumn{1}{c}{median} & \multicolumn{1}{c}{max}\\
\midrule
Gemma 3 4B $\to$ Gemma 3 12B & NIH ChestX-ray14 & bit-identical & 5/6 & 0.33 & 0.83 & 2.7 & 6.9 \\
Gemma 3 4B $\to$ Gemma 3 12B & CheXpert Plus & bit-identical & 6/6 & 0.29 & 0.83 & 2.1 & 5.9 \\
Gemma 3 4B $\to$ Gemma 3 12B & COCO & bit-identical & 5/6 & 0.23 & 0.87 & 1.6 & 6.0 \\
Gemma 3 12B $\to$ Gemma 3 27B & NIH ChestX-ray14 & bf16-equal & 6/6 & 0.37 & 0.79 & 3.7 & 8.0 \\
Gemma 3 12B $\to$ Gemma 3 27B & CheXpert Plus & bf16-equal & 6/6 & 0.35 & 0.81 & 1.4 & 3.3 \\
Gemma 3 12B $\to$ Gemma 3 27B & COCO & bf16-equal & 5/6 & 0.15 & 0.40 & 1.7 & 4.3 \\
Gemma 3 4B $\to$ Gemma 3 27B & NIH ChestX-ray14 & bf16-equal & 5/6 & 0.07 & 0.24 & 1.8 & 6.0 \\
Gemma 3 4B $\to$ Gemma 3 27B & CheXpert Plus & bf16-equal & 6/6 & 0.17 & 0.65 & 4.6 & 17.3 \\
Gemma 3 4B $\to$ Gemma 3 27B & COCO & bf16-equal & 6/6 & 0.16 & 1.02 & 1.4 & 8.8 \\
MedGemma 4B $\to$ MedGemma 27B & NIH ChestX-ray14 & bf16-equal & 5/6 & 0.09 & 0.30 & 1.5 & 5.0 \\
MedGemma 4B $\to$ MedGemma 27B & CheXpert Plus & bf16-equal & 5/6 & 0.19 & 0.54 & 1.8 & 5.0 \\
MedGemma 4B $\to$ MedGemma 27B & COCO & bf16-equal & 3/6 & 0.05 & 0.35 & 0.6 & 4.6 \\
\bottomrule
\end{tabular}
\end{table}
\begin{table}[h]
\centering
\scriptsize
\setlength{\tabcolsep}{4pt}
\caption{\textbf{Ownership under alternative directions, lifts, token weights, and competitor families.} Cells owned under the paper's rule within each family, and the median $O_q$, per dataset. Each panel repeats the reference write on its own blocks.}
\label{tab:cf-altdir}
\begin{tabular}{llrrrrrr}
\toprule
\multirow{2}{*}{Direction} & \multirow{2}{*}{Lift} & \multicolumn{2}{c}{NIH ChestX-ray14} & \multicolumn{2}{c}{CheXpert Plus} & \multicolumn{2}{c}{COCO}\\
\cmidrule(lr){3-4}\cmidrule(lr){5-6}\cmidrule(lr){7-8}
&  & \multicolumn{1}{c}{owned} & \multicolumn{1}{c}{med.\ $O_q$} & \multicolumn{1}{c}{owned} & \multicolumn{1}{c}{med.\ $O_q$} & \multicolumn{1}{c}{owned} & \multicolumn{1}{c}{med.\ $O_q$}\\
\midrule
\multicolumn{8}{l}{\textit{Direction construction (20 blocks per dataset)}} \\
\quad logistic normal & coefficient & 7/120 & -0.03 & 19/120 & -0.01 & 101/120 & +0.18 \\
\quad difference of means & coefficient & 16/120 & -0.02 & 18/120 & -0.01 & 62/120 & +0.02 \\
\quad Haufe pattern & coefficient & 17/120 & -0.02 & 17/120 & -0.01 & 53/120 & +0.01 \\
\quad orthogonalised normal & coefficient & 8/120 & -0.02 & 14/120 & -0.01 & 98/120 & +0.13 \\
\quad residualised normal & coefficient & 8/120 & -0.02 & 13/120 & -0.01 & 102/120 & +0.20 \\
\quad owned under at least one &  & 25/120 & -- & 41/120 & -- & 108/120 & -- \\
\midrule
\multicolumn{8}{l}{\textit{Displacement lift (19 blocks per dataset)}} \\
\quad logistic normal & coefficient & 7/114 & -0.03 & 19/114 & -0.01 & 100/114 & +0.19 \\
\quad difference of means & displacement & 21/114 & -0.03 & 19/114 & -0.00 & 58/114 & +0.02 \\
\quad Haufe pattern & displacement & 17/114 & -0.02 & 17/114 & -0.00 & 49/114 & +0.01 \\
\midrule
\multicolumn{8}{l}{\textit{Token weights of the logistic normal (4 NIH, 2 CheXpert, 2 COCO blocks)}} \\
\quad uniform & coefficient & 3/24 & -0.08 & 5/12 & +0.00 & 12/12 & +0.60 \\
\quad softmax of the probe score & coefficient & 0/24 & -0.01 & 0/12 & -0.00 & 12/12 & +0.16 \\
\quad top quarter of tokens & coefficient & 2/24 & -0.12 & 3/12 & -0.05 & 12/12 & +0.57 \\
\midrule
\multicolumn{8}{l}{\textit{Competitor family of the logistic normal (5 NIH, 4 CheXpert blocks)}} \\
\quad six clinical directions & coefficient & 3/30 & -0.07 & 8/24 & -0.07 & -- & -- \\
\quad with every extra dataset label & coefficient & 3/30 & -0.07 & 5/24 & -0.10 & -- & -- \\
\bottomrule
\end{tabular}
\end{table}
\begin{table}[h]
\centering
\scriptsize
\setlength{\tabcolsep}{3.5pt}
\caption{\textbf{The answer direction.} Top, on the primary template: concepts owned by the label direction $\hat w_q$ and by the answer direction $a_q$, split into kept, rescued, and lost, then medians over the six questions. Bottom: $a_q$ under each held-out template, with the concepts owned, in parentheses those already owned under IY, and the kept (concept, template) pairs.}
\label{tab:cf-ansdir}
\begin{tabular}{llrrrrrrrrrr}
\toprule
\multirow{2}{*}{Checkpoint} & \multirow{2}{*}{Dataset} & \multicolumn{5}{c}{owned concepts} & \multirow{2}{*}{$a_q>$ comp.} & \multicolumn{4}{c}{median}\\
\cmidrule(lr){3-7}\cmidrule(lr){9-12}
&  & \multicolumn{1}{c}{$\hat w_q$} & \multicolumn{1}{c}{$a_q$} & \multicolumn{1}{c}{kept} & \multicolumn{1}{c}{rescued} & \multicolumn{1}{c}{lost} &  & \multicolumn{1}{c}{$W^a_{q,q}$} & \multicolumn{1}{c}{$\cos$} & \multicolumn{1}{c}{$\cos_\Sigma$} & \multicolumn{1}{c}{$R^2$}\\
\midrule
Qwen2.5-VL-7B & NIH ChestX-ray14 & 0/6 & 3/6 & 0 & 3 & 0 & 6/6 & 0.61 & 0.07 & 0.35 & 0.54 \\
Qwen2.5-VL-7B & CheXpert Plus & 0/6 & 6/6 & 0 & 6 & 0 & 6/6 & 0.55 & 0.06 & 0.36 & 0.48 \\
Qwen2.5-VL-7B & COCO & 6/6 & 6/6 & 6 & 0 & 0 & 6/6 & 0.78 & 0.64 & 0.84 & 0.58 \\
Gemma 3 12B & NIH ChestX-ray14 & 1/6 & 2/6 & 1 & 1 & 0 & 5/6 & 0.47 & 0.18 & 0.33 & 0.17 \\
Gemma 3 12B & CheXpert Plus & 1/6 & 0/6 & 0 & 0 & 1 & 1/6 & 0.37 & 0.17 & 0.10 & 0.19 \\
Gemma 3 12B & COCO & 6/6 & 6/6 & 6 & 0 & 0 & 6/6 & 0.61 & 0.59 & 0.79 & 0.37 \\
Lingshu 32B & NIH ChestX-ray14 & 2/6 & 3/6 & 2 & 1 & 0 & 6/6 & 0.24 & 0.30 & 0.60 & 0.77 \\
Lingshu 32B & CheXpert Plus & 5/6 & 5/6 & 5 & 0 & 0 & 6/6 & 0.22 & 0.33 & 0.87 & 0.80 \\
Lingshu 32B & COCO & 6/6 & 6/6 & 6 & 0 & 0 & 6/6 & 0.36 & 0.54 & 0.73 & 0.55 \\
\midrule
NIH ChestX-ray14 & 3 blocks & 3/18 & 8/18 & 3 & 5 & 0 & 17/18 & 0.47 & 0.18 & 0.43 & 0.54 \\
CheXpert Plus & 3 blocks & 6/18 & 11/18 & 5 & 6 & 1 & 13/18 & 0.35 & 0.16 & 0.36 & 0.48 \\
\textit{chest} & 6 blocks & 9/36 & 19/36 & 8 & 11 & 1 & 30/36 & 0.41 & 0.18 & 0.39 & 0.50 \\
COCO & 3 blocks & 18/18 & 18/18 & 18 & 0 & 0 & 18/18 & 0.61 & 0.61 & 0.81 & 0.54 \\
\bottomrule
\end{tabular}

\medskip
\begin{tabular}{llrrrrrr}
\toprule
\multirow{2}{*}{Checkpoint} & \multirow{2}{*}{Dataset} & \multicolumn{5}{c}{held-out template: owned (IY-owned kept)} & \multirow{2}{*}{kept pairs}\\
\cmidrule(lr){3-7}
&  & \multicolumn{1}{c}{WY} & \multicolumn{1}{c}{IA} & \multicolumn{1}{c}{IB} & \multicolumn{1}{c}{WA} & \multicolumn{1}{c}{WB} & \\
\midrule
Qwen2.5-VL-7B & NIH ChestX-ray14 & 4 (3) & 3 (3) & 4 (3) & 4 (3) & 4 (3) & 15/15 \\
Qwen2.5-VL-7B & CheXpert Plus & 4 (4) & 6 (6) & 5 (5) & 5 (5) & 5 (5) & 25/30 \\
Qwen2.5-VL-7B & COCO & 6 (6) & 6 (6) & 6 (6) & 6 (6) & 6 (6) & 30/30 \\
Gemma 3 12B & NIH ChestX-ray14 & 2 (1) & 1 (1) & 0 (0) & 1 (0) & 1 (0) & 2/10 \\
Gemma 3 12B & CheXpert Plus & 0 (0) & 1 (0) & 1 (0) & 1 (0) & 1 (0) & 0/0 \\
Gemma 3 12B & COCO & 5 (5) & 6 (6) & 6 (6) & 6 (6) & 6 (6) & 29/30 \\
Lingshu 32B & NIH ChestX-ray14 & 3 (3) & 3 (3) & 3 (3) & 3 (3) & 3 (3) & 15/15 \\
Lingshu 32B & CheXpert Plus & 5 (5) & 3 (3) & 3 (3) & 4 (4) & 4 (4) & 19/25 \\
Lingshu 32B & COCO & 6 (6) & 6 (6) & 6 (6) & 6 (6) & 6 (6) & 30/30 \\
\midrule
\textit{chest} & 6 blocks & 18 (16) & 17 (16) & 16 (14) & 18 (15) & 18 (15) & 76/95 \\
COCO & 3 blocks & 17 (17) & 18 (18) & 18 (18) & 18 (18) & 18 (18) & 89/90 \\
\bottomrule
\end{tabular}
\end{table}
\begin{table}[htbp]
\centering
\scriptsize
\setlength{\tabcolsep}{4pt}
\caption{\textbf{Non-clinical attributes of the same radiographs.} Top: per checkpoint and dataset, owned attribute and owned clinical cells inside the same nine-direction family, and which attributes are owned. Bottom: per attribute and dataset, medians over blocks of the probe AUROC, selectivity, clean-answer AUROC, and cosine to the clinical normals.}
\label{tab:cf-attr}
\begin{tabular}{llrrccc}
\toprule
Checkpoint & Dataset & \multicolumn{1}{c}{attributes owned} & \multicolumn{1}{c}{findings owned} & \multicolumn{1}{c}{view AP} & \multicolumn{1}{c}{sex F} & \multicolumn{1}{c}{age $\geq60$}\\
\midrule
Qwen2.5-VL-3B & NIH ChestX-ray14 & 0/3 & 0/6 & -- & -- & -- \\
Qwen2.5-VL-3B & CheXpert Plus & 0/3 & 0/6 & -- & -- & -- \\
Qwen2.5-VL-7B & NIH ChestX-ray14 & 0/3 & 0/6 & -- & -- & -- \\
Qwen2.5-VL-7B & CheXpert Plus & 0/3 & 0/6 & -- & -- & -- \\
Qwen2.5-VL-32B & NIH ChestX-ray14 & 0/3 & 0/6 & -- & -- & -- \\
Qwen2.5-VL-32B & CheXpert Plus & 0/3 & 0/6 & -- & -- & -- \\
Qwen3-VL-4B & NIH ChestX-ray14 & 1/3 & 0/6 & -- & yes & -- \\
Qwen3-VL-4B & CheXpert Plus & 1/3 & 2/6 & -- & yes & -- \\
Qwen3-VL-8B & NIH ChestX-ray14 & 0/3 & 0/6 & -- & -- & -- \\
Qwen3-VL-8B & CheXpert Plus & 1/3 & 0/6 & -- & -- & yes \\
Qwen3-VL-32B & NIH ChestX-ray14 & 0/3 & 1/6 & -- & -- & -- \\
Qwen3-VL-32B & CheXpert Plus & 1/3 & 1/6 & -- & -- & yes \\
InternVL3.5-8B & NIH ChestX-ray14 & 0/3 & 0/6 & -- & -- & -- \\
InternVL3.5-8B & CheXpert Plus & 0/3 & 1/6 & -- & -- & -- \\
InternVL3.5-14B & NIH ChestX-ray14 & 0/3 & 1/6 & -- & -- & -- \\
InternVL3.5-14B & CheXpert Plus & 0/3 & 1/6 & -- & -- & -- \\
InternVL3.5-38B & NIH ChestX-ray14 & 0/3 & 0/6 & -- & -- & -- \\
InternVL3.5-38B & CheXpert Plus & 0/3 & 2/6 & -- & -- & -- \\
Gemma 3 4B & NIH ChestX-ray14 & 0/3 & 0/6 & -- & -- & -- \\
Gemma 3 4B & CheXpert Plus & 0/3 & 0/6 & -- & -- & -- \\
Gemma 3 12B & NIH ChestX-ray14 & 0/3 & 1/6 & -- & -- & -- \\
Gemma 3 12B & CheXpert Plus & 0/3 & 1/6 & -- & -- & -- \\
Gemma 3 27B & NIH ChestX-ray14 & 0/3 & 0/6 & -- & -- & -- \\
Gemma 3 27B & CheXpert Plus & 0/3 & 0/6 & -- & -- & -- \\
MedGemma 4B & NIH ChestX-ray14 & 0/3 & 0/6 & -- & -- & -- \\
MedGemma 4B & CheXpert Plus & 1/3 & 1/6 & -- & -- & yes \\
MedGemma 27B & NIH ChestX-ray14 & 1/3 & 0/6 & -- & -- & yes \\
MedGemma 27B & CheXpert Plus & 0/3 & 0/6 & -- & -- & -- \\
Lingshu 7B & NIH ChestX-ray14 & 2/3 & 1/6 & -- & yes & yes \\
Lingshu 7B & CheXpert Plus & 0/3 & 3/6 & -- & -- & -- \\
Lingshu 32B & NIH ChestX-ray14 & 2/3 & 2/6 & -- & yes & yes \\
Lingshu 32B & CheXpert Plus & 1/3 & 4/6 & -- & -- & yes \\
LLaVA-1.5-7B & NIH ChestX-ray14 & 0/3 & 0/6 & -- & -- & -- \\
LLaVA-1.5-7B & CheXpert Plus & 0/3 & 0/6 & -- & -- & -- \\
LLaVA-1.5-13B & NIH ChestX-ray14 & 1/3 & 0/6 & yes & -- & -- \\
LLaVA-1.5-13B & CheXpert Plus & 0/3 & 0/6 & -- & -- & -- \\
LLaVA-Med 1.5 7B & NIH ChestX-ray14 & 0/3 & 0/6 & -- & -- & -- \\
LLaVA-Med 1.5 7B & CheXpert Plus & 0/3 & 1/6 & -- & -- & -- \\
\midrule
NIH ChestX-ray14 & 19 blocks & 7/57 & 6/114 & 1/19 & 3/19 & 3/19 \\
CheXpert Plus & 19 blocks & 5/57 & 17/114 & 0/19 & 1/19 & 4/19 \\
\textit{chest} & 38 blocks & 12/114 & 23/228 & 1/38 & 4/38 & 7/38 \\
\midrule
\multicolumn{7}{l}{\textit{Attribute probes and directions: medians over the blocks of each dataset}} \\
Attribute & Dataset & probe AUROC & selectivity & answer AUROC & med.\ $|\cos|$ & owned \\
\midrule
AP (portable) projection & NIH ChestX-ray14 & 0.999 & 0.32 & 0.76 & 0.03 & 1/19 \\
AP (portable) projection & CheXpert Plus & 0.998 & 0.31 & 0.73 & 0.05 & 0/19 \\
AP (portable) projection & \textit{chest} & 0.999 & 0.31 & 0.74 & 0.04 & 1/38 \\
female sex & NIH ChestX-ray14 & 0.984 & 0.30 & 0.73 & 0.05 & 3/19 \\
female sex & CheXpert Plus & 0.957 & 0.27 & 0.65 & 0.06 & 1/19 \\
female sex & \textit{chest} & 0.971 & 0.28 & 0.71 & 0.05 & 4/38 \\
age $\geq60$ & NIH ChestX-ray14 & 0.896 & 0.21 & 0.62 & 0.05 & 3/19 \\
age $\geq60$ & CheXpert Plus & 0.898 & 0.21 & 0.64 & 0.06 & 4/19 \\
age $\geq60$ & \textit{chest} & 0.896 & 0.21 & 0.63 & 0.06 & 7/38 \\
\midrule
\multicolumn{7}{l}{\textit{The same comparison read three ways, pooled over the 38 chest blocks}} \\
Comparison & unit & attribute cells & owned & finding cells & owned & difference \\
\midrule
every cell & cell & 114 & 12 & 228 & 23 & 0.004 \\
answer-capable cells only & cell & 100 & 11 & 141 & 22 & -0.046 \\
answerability-matched pairs & pair & 188 & 14 & 188 & 17 & -0.016 \\
\bottomrule
\end{tabular}
\end{table}
\begin{table}[t]
\centering
\small
\setlength{\tabcolsep}{4pt}
\caption{\textbf{Attributes and findings by clean-answer AUROC.} The nine-direction family's cells binned by the clean answer's AUROC against the question's own label; each entry is owned cells of cells in the bin.}
\label{tab:cf-attr-strat}
\begin{tabular}{lrrrrrr}
\toprule
Cells & \multicolumn{1}{c}{$<$ 0.50} & \multicolumn{1}{c}{0.50--0.70} & \multicolumn{1}{c}{0.70--0.80} & \multicolumn{1}{c}{0.80--0.90} & \multicolumn{1}{c}{$\geq$ 0.90} & \multicolumn{1}{c}{all cells}\\
\midrule
attributes, NIH ChestX-ray14 & 1/4 & 1/26 & 3/14 & 2/6 & 0/7 & 7/57 \\
attributes, CheXpert Plus & 0/5 & 3/29 & 0/12 & 2/5 & 0/6 & 5/57 \\
\textit{attributes, chest} & 1/9 & 4/55 & 3/26 & 4/11 & 0/13 & 12/114 \\
findings, NIH ChestX-ray14 & 0/6 & 2/63 & 2/32 & 2/13 & 0/0 & 6/114 \\
findings, CheXpert Plus & 0/13 & 4/47 & 2/9 & 5/21 & 6/24 & 17/114 \\
\textit{findings, chest} & 0/19 & 6/110 & 4/41 & 7/34 & 6/24 & 23/228 \\
\bottomrule
\end{tabular}
\end{table}
\begin{table}[h]
\centering
\scriptsize
\setlength{\tabcolsep}{4pt}
\caption{\textbf{The answer direction and column selectivity.} Top, per block with the answer-direction module: medians over the six questions of the AUROCs, the Pearson correlation, and the raw and whitened cosines between $a_q$ and $\hat w_q$. Bottom, per dataset: the column-selectivity index $S_d=W_{d,d}/\sum_q|W_{q,d}|$, and ownership recomputed on the rows whose label for $q$ is known.}
\label{tab:cf-validation}
\begin{tabular}{llrrrrrrr}
\toprule
\multirow{2}{*}{Checkpoint} & \multirow{2}{*}{Dataset} & \multicolumn{2}{c}{AUROC, label} & \multicolumn{2}{c}{AUROC, clean answer} & \multirow{2}{*}{Pearson} & \multicolumn{2}{c}{cosine}\\
\cmidrule(lr){3-4}\cmidrule(lr){5-6}\cmidrule(lr){8-9}
&  & \multicolumn{1}{c}{$a_q$} & \multicolumn{1}{c}{$\hat w_q$} & \multicolumn{1}{c}{$a_q$} & \multicolumn{1}{c}{$\hat w_q$} &  & \multicolumn{1}{c}{raw} & \multicolumn{1}{c}{$\Sigma$}\\
\midrule
Qwen2.5-VL-7B & NIH ChestX-ray14 & 0.60 & 0.73 & 0.93 & 0.66 & 0.36 & 0.07 & 0.35 \\
Qwen2.5-VL-7B & CheXpert Plus & 0.67 & 0.86 & 0.85 & 0.66 & 0.34 & 0.06 & 0.36 \\
Qwen2.5-VL-7B & COCO & 0.97 & 0.97 & 0.98 & 0.98 & 0.84 & 0.64 & 0.84 \\
Lingshu 32B & NIH ChestX-ray14 & 0.75 & 0.76 & 0.95 & 0.79 & 0.62 & 0.30 & 0.60 \\
Lingshu 32B & CheXpert Plus & 0.85 & 0.86 & 0.95 & 0.89 & 0.88 & 0.33 & 0.87 \\
Lingshu 32B & COCO & 0.94 & 0.97 & 0.93 & 0.92 & 0.75 & 0.54 & 0.73 \\
Gemma 3 12B & NIH ChestX-ray14 & 0.59 & 0.65 & 0.72 & 0.61 & 0.31 & 0.18 & 0.33 \\
Gemma 3 12B & CheXpert Plus & 0.61 & 0.78 & 0.73 & 0.56 & 0.13 & 0.17 & 0.10 \\
Gemma 3 12B & COCO & 0.93 & 0.96 & 0.80 & 0.76 & 0.81 & 0.59 & 0.79 \\
\midrule
\multicolumn{9}{p{0.92\linewidth}}{\textit{Alignment against ownership over the 54 answer-direction cells: Spearman correlation between $\cos_\Sigma(a_q,\hat w_q)$ and the label direction's $O_q$ $\rho=0.75$ ($p=8\times10^{-11}$); with the label direction's owned indicator $\rho=0.74$ ($p=1\times10^{-10}$); median $\cos_\Sigma$ 0.80 in owned against 0.34 in not-owned cells; within the chest sets alone $\rho=0.59$ ($p=2\times10^{-4}$, 36 cells).}} \\
\midrule
\multicolumn{9}{l}{\textit{Column selectivity and known-label ownership, per dataset}} \\
Dataset & blocks & cells & med.\ $S_d$ & $S_d\geq0.5$ & known rows & sign kept & owned kept & \\
\midrule
NIH ChestX-ray14 & 25 & 150 & 0.12 & 2/150 & 600 & 150/150 & 13/13 & \\
CheXpert Plus & 25 & 150 & 0.12 & 2/150 & 165 & 143/150 & 23/25 & \\
COCO & 25 & 150 & 0.57 & 92/150 & 600 & 150/150 & 113/113 & \\
\bottomrule
\end{tabular}
\end{table}
\begin{table}[h]
\centering
\scriptsize
\setlength{\tabcolsep}{4pt}
\caption{\textbf{The grade on radiologist labels, and what refitting on them changes.} Top, per CheXpert block with the valid module: the complete grade recomputed on 200 radiologist-labelled films, against the same grade on the 600 labeler-labelled test rows. Bottom: the six directions refitted on the radiologist labels.}
\label{tab:cf-valid}
\begin{tabular}{lrrrrr}
\toprule
Checkpoint & \multicolumn{1}{c}{owned (test)} & \multicolumn{1}{c}{owned (valid)} & \multicolumn{1}{c}{verdict agrees} & \multicolumn{1}{c}{owned grade agrees} & \multicolumn{1}{c}{median $|\Delta O_q|$}\\
\midrule
Qwen2.5-VL-3B & 0 & 0 & 5/6 & 6/6 & 0.005 \\
Qwen2.5-VL-7B & 0 & 0 & 6/6 & 6/6 & 0.005 \\
Qwen2.5-VL-32B & 0 & 0 & 6/6 & 6/6 & 0.004 \\
Qwen3-VL-4B & 2 & 2 & 4/6 & 6/6 & 0.009 \\
Qwen3-VL-8B & 0 & 0 & 5/6 & 6/6 & 0.003 \\
Qwen3-VL-32B & 1 & 1 & 6/6 & 6/6 & 0.004 \\
InternVL3.5-8B & 1 & 0 & 5/6 & 5/6 & 0.001 \\
InternVL3.5-14B & 1 & 1 & 3/6 & 6/6 & 0.002 \\
InternVL3.5-38B & 2 & 2 & 6/6 & 6/6 & 0.001 \\
Gemma 3 4B & 0 & 0 & 4/6 & 6/6 & 0.007 \\
Gemma 3 12B & 1 & 1 & 6/6 & 6/6 & 0.016 \\
Gemma 3 27B & 0 & 0 & 3/6 & 6/6 & 0.002 \\
MedGemma 4B & 2 & 1 & 5/6 & 5/6 & 0.025 \\
MedGemma 27B & 0 & 0 & 5/6 & 6/6 & 0.010 \\
Lingshu 7B & 3 & 3 & 6/6 & 6/6 & 0.007 \\
Lingshu 32B & 5 & 5 & 6/6 & 6/6 & 0.005 \\
LLaVA-1.5-7B & 0 & 0 & 5/6 & 6/6 & 0.003 \\
LLaVA-1.5-13B & 0 & 0 & 5/6 & 6/6 & 0.000 \\
LLaVA-Med 1.5 7B & 1 & 1 & 6/6 & 6/6 & 0.000 \\
\midrule
\textit{all} (19 blocks) & 19 & 17 & 97/114 & 112/114 & 0.003 \\
\bottomrule
\end{tabular}

\medskip
\setlength{\tabcolsep}{4pt}
\begin{tabular}{llrrrrr}
\toprule
\multicolumn{7}{p{0.95\linewidth}}{\textit{The label source of a refitted direction against the number of rows it was fitted on. The expert-label refit is estimated on 200 radiologist-labelled films and the shipped direction on about 20{,}000 report-labelled rows, so comparing only those two confounds the two. Every arm keeps the block's own projection, its stored train-only scaler, the protocol's probe settings and the same 5-fold cross-fitting over patients, and is scored on the same films by a direction not fitted on them. The two training subsample arms redraw the fitting rows 20 times and report the median; the second draws enough rows to leave the concept with as many labelled rows as the expert refit has. Cosines are to the shipped direction. Medians over the 36 cells of 6 CheXpert blocks.}} \\
\midrule
 & & & \multicolumn{2}{c}{cross-fitted AUROC against} & \multicolumn{2}{c}{cosine to shipped} \\
\cmidrule(lr){4-5}\cmidrule(lr){6-7}
Labels & Fitted on & rows & radiologist & report & raw & whitened \\
\midrule
radiologist & 200 valid films & 160 & 0.69 & 0.71 & 0.12 & 0.56 \\
report-derived & the same 200 valid films & 44 & 0.67 & 0.71 & 0.12 & 0.59 \\
report-derived & 200 training rows & 42 & 0.66 & 0.70 & 0.12 & 0.58 \\
report-derived & 200 labelled training rows & 160 & 0.70 & 0.74 & 0.22 & 0.69 \\
report-derived & every training row & 5332 & 0.79 & 0.86 & 1.00 & 1.00 \\
\midrule
\multicolumn{7}{l}{\textit{What each contrast isolates, in AUROC against the radiologist labels}} \\
\midrule
\multicolumn{6}{l}{label source, both arms fitting the same number of rows} & -0.01 \\
\multicolumn{6}{l}{sample size, both arms on report-derived labels} & 0.09 \\
\multicolumn{6}{l}{the expert refit against the shipped direction, which mixes the two} & -0.11 \\
\bottomrule
\end{tabular}
\end{table}

\subsection{Crossing the factors, replaying the tensors, and changing the endpoint}
\label{app:cf-crossed}

\paragraph{Crossed tower and reader.} The swap replaces the host checkpoint's vision tower with the donor checkpoint's tower, leaving the host's connector and language model untouched. For every swap, the module records a receipt listing the replaced tensors of the tower root, their parameter count, a digest of the donor tower, and tensor-by-tensor comparisons of the loaded tower against the donor and of everything outside the tower against the host's own weights. Of the host tower's \cfSwapTowerTensors{} tensors, \cfSwapTensors{} carry weights; the module replaces each and verifies bitwise equality with the donor's (\cfSwapParamsM{} million parameters). The module rebuilds the remaining tensor, a position buffer. None of the \cfSwapOutside{} tensors outside the tower changes. Write directions are fitted in the representation each tower produces, so a swapped arm writes the donor's directions. The two factors therefore carry different loads: at a fixed tower the swap changes the connector and the language model and nothing else, while replacing the tower replaces the representation and the six directions fitted on it together. The tower factor is the joint effect of both; the reader factor is the effect of neither. Each of the \cfSwapCombos{} combinations of tower and reader is graded on the same \cfSwapRows{} test rows of the block, at the primary dose and template, against the 119-direction random family and the sham of the tower in use, with the $6\times5$ max-$T$ verdict (Table~\ref{tab:cf-ledger}). The two native combinations give the two checkpoints' own grades on that cohort. Across the \cfSwapNativeCells{} native cells, the ownership contrast changes by at most \cfSwapNativeMaxDO{} relative to the grade on the block's whole test cohort. On the smaller cohort, \cfSwapNativeUnresolved{} of the \cfSwapNativeOwnedFull{} cells owned in that grade become unresolved. The crossover contrasts hold one factor fixed and change the other, and report three quantities of the same six questions: the own write $W_{q,q}$, the ownership contrast $O_q$, and the margin of the own write over the strongest competitor in the whole reference family, $M_q=W_{q,q}-\max(\max_{d\neq q}W_{q,d},\,\text{random }p_{95},\,|\text{sham}|)$, each as the mean over the six questions of the absolute paired change, with a percentile interval over the same 2{,}000 draws of the rows and, in brackets, the number of questions of six whose change is simultaneously nonzero. The last two rows of each block average each factor over its two levels. The three come from one pass over the four arms' packaged outcomes, so a row's columns are the same arms read three ways. The four arms of a dataset are the same four arms whichever of the pair's blocks indexes them, and the two host blocks give the same crossover: the middle panel counts each dataset's four combinations once, and an arm graded in both host blocks agrees in all \cfSwapArmsTwice{} cases where it appears twice. A dash marks an arm that is not graded in that host block. Table~\ref{tab:cf-towerswap} reports every combination and contrast.

\paragraph{What the crossing owns.} Gemma~3 4B and MedGemma 4B give \cfSwapCombos{} combinations of tower and reader on each dataset, and written directions travel with the tower that produced them. On NIH ChestX-ray14 and CheXpert Plus, no combination owns a finding: \cfSwapChestOwned{} of \cfSwapChestCells{} cells across the \cfSwapChestCombos{} combinations. On COCO, the same crossing owns \cfSwapCocoOwned{} of \cfSwapCocoCells{} cells. Each native combination owns \cfSwapCocoNativeMin{} of the six objects, the Gemma tower behind the MedGemma reader owns \cfSwapCocoGemmaTowerCross{}, and the MedGemma tower behind the Gemma~3 reader owns \cfSwapCocoMedTowerCross{}. The crossover holds one factor fixed and changes the other on three quantities of the same six questions. On NIH, the mean absolute reader effect is \cfSwapReaderNihMin{} $[\cfSwapReaderNihLo,\cfSwapReaderNihHi]$ on the own write and \cfSwapReaderNihOMin{} $[\cfSwapReaderNihOLo,\cfSwapReaderNihOHi]$ on the ownership contrast, against a tower effect of \cfSwapTowerNihMin{} $[\cfSwapTowerNihLo,\cfSwapTowerNihHi]$ and \cfSwapTowerNihOMin{} $[\cfSwapTowerNihOLo,\cfSwapTowerNihOHi]$. On CheXpert the reader effect is \cfSwapReaderChexMin{} on the own write and \cfSwapReaderChexOMin{} on the ownership contrast, against \cfSwapTowerChexMin{} and \cfSwapTowerChexOMin{} for the tower. On the family margin the two chest sets give \cfSwapReaderNihMMin{} and \cfSwapReaderChexMMin{} for the reader against \cfSwapTowerNihMMin{} and \cfSwapTowerChexMMin{} for the tower. On COCO the tower effect is the larger of the two on all three quantities: \cfSwapTowerCocoMin{}, \cfSwapTowerCocoOMin{}, and \cfSwapTowerCocoMMin{} against \cfSwapReaderCocoMin{}, \cfSwapReaderCocoOMin{}, and \cfSwapReaderCocoMMin{}.

\paragraph{How well the crossed arms answer.} The hybrid arms answer worse than the native arms. On the labelled rows the crossing grades, with no write applied, the median clean-answer AUROC over a block's six concepts falls from \cfSwapAnsNativeGemCoco{} in the native Gemma~3 4B arm on COCO to \cfSwapAnsHybridGemCoco{} in the MedGemma tower behind that reader, from \cfSwapAnsNativeMedNih{} in the native MedGemma 4B arm on NIH ChestX-ray14 to \cfSwapAnsHybridMedNih{} in the Gemma~3 tower behind that reader, and from \cfSwapAnsNativeMedChex{} to \cfSwapAnsHybridMedChex{} on CheXpert Plus. The remaining three arms give \cfSwapAnsNativeGemNih{} against \cfSwapAnsHybridGemNih{} and \cfSwapAnsNativeGemChex{} against \cfSwapAnsHybridGemChex{} in the Gemma~3 4B blocks, and \cfSwapAnsNativeMedCoco{} against \cfSwapAnsHybridMedCoco{} in the MedGemma 4B block on COCO. All \cfSwapAnsWorse{} of the \cfSwapAnsArms{} hybrid arms fall below the native arm of their own block, by \cfSwapAnsDropMin{} to \cfSwapAnsDropMax{}.

\begin{table}[htbp]
\centering
\scriptsize
\setlength{\tabcolsep}{4pt}
\caption{\textbf{Crossing the tower and the reader.} The vision tower of one checkpoint runs behind the other's connector and language model; G is Gemma~3 4B, M is MedGemma 4B, and each column of the top two panels is tower/reader. The bottom panel holds one factor, changes the other, and gives the mean absolute paired change over the six questions.}
\label{tab:cf-towerswap}
\begin{tabular}{llcccccc}
\toprule
\multirow{2}{*}{Host checkpoint} & \multirow{2}{*}{Dataset} & \multicolumn{2}{c}{native} & \multicolumn{2}{c}{swapped} & \multirow{2}{*}{concepts} & \multirow{2}{*}{crossover}\\
\cmidrule(lr){3-4}\cmidrule(lr){5-6}
&  & \multicolumn{1}{c}{G/G} & \multicolumn{1}{c}{M/M} & \multicolumn{1}{c}{M/G} & \multicolumn{1}{c}{G/M} &  & \\
\midrule
Gemma 3 4B & NIH ChestX-ray14 & 0 & 0 & 0 & 0 & 6 & complete \\
Gemma 3 4B & CheXpert Plus & 0 & 0 & 0 & 0 & 6 & complete \\
Gemma 3 4B & COCO & 5 & 5 & 0 & 4 & 6 & complete \\
MedGemma 4B & NIH ChestX-ray14 & 0 & 0 & 0 & 0 & 6 & complete \\
MedGemma 4B & CheXpert Plus & 0 & 0 & -- & 0 & 6 & -- \\
MedGemma 4B & COCO & 5 & 5 & -- & 4 & 6 & -- \\
\midrule
\multicolumn{8}{l}{\textit{The four combinations of each dataset, shared arms counted once}} \\
Combinations & Dataset & G/G & M/M & M/G & G/M & concepts & owned of cells \\
\midrule
all four & NIH ChestX-ray14 & 0 & 0 & 0 & 0 & 6 & 0/24 \\
all four & CheXpert Plus & 0 & 0 & 0 & 0 & 6 & 0/24 \\
all four & COCO & 5 & 5 & 0 & 4 & 6 & 14/24 \\
\bottomrule
\end{tabular}
\vspace{4pt}
\setlength{\tabcolsep}{2.5pt}
\begin{tabular}{llccc}
\toprule
\multicolumn{5}{l}{\textit{Crossover: the effect of replacing one factor with the other held fixed}} \\
\midrule
Block & Changed, held & own write $W_{q,q}$ & ownership $O_q$ & family margin $M_q$ \\
\midrule
Gemma 3 4B, NIH & reader, own tower & 0.133 [0.106, 0.169] (1) & 0.158 [0.136, 0.187] (4) & 0.425 [0.393, 0.458] (5) \\
 & reader, partner tower & 0.240 [0.221, 0.259] (5) & 0.210 [0.192, 0.233] (4) & 0.286 [0.268, 0.309] (6) \\
 & tower, own reader & 0.164 [0.140, 0.189] (4) & 0.061 [0.046, 0.080] (1) & 0.137 [0.120, 0.164] (4) \\
 & tower, partner reader & 0.160 [0.125, 0.193] (6) & 0.265 [0.239, 0.299] (5) & 0.496 [0.462, 0.525] (6) \\
 & \quad mean reader effect & 0.186 [0.170, 0.207] & 0.184 [0.169, 0.203] & 0.356 [0.339, 0.376] \\
 & \quad mean tower effect & 0.162 [0.140, 0.182] & 0.163 [0.149, 0.185] & 0.316 [0.299, 0.338] \\
\addlinespace
Gemma 3 4B, CheXpert & reader, own tower & 0.382 [0.353, 0.412] (5) & 0.303 [0.273, 0.332] (6) & 0.490 [0.465, 0.515] (5) \\
 & reader, partner tower & 0.125 [0.102, 0.151] (4) & 0.214 [0.194, 0.240] (5) & 0.240 [0.214, 0.263] (6) \\
 & tower, own reader & 0.269 [0.245, 0.293] (5) & 0.223 [0.205, 0.241] (6) & 0.372 [0.352, 0.390] (5) \\
 & tower, partner reader & 0.183 [0.153, 0.217] (4) & 0.297 [0.267, 0.332] (5) & 0.517 [0.485, 0.553] (6) \\
 & \quad mean reader effect & 0.254 [0.235, 0.274] & 0.259 [0.240, 0.280] & 0.365 [0.347, 0.383] \\
 & \quad mean tower effect & 0.226 [0.208, 0.246] & 0.260 [0.243, 0.279] & 0.444 [0.425, 0.466] \\
\addlinespace
Gemma 3 4B, COCO & reader, own tower & 0.072 [0.049, 0.104] (2) & 0.229 [0.187, 0.271] (5) & 0.143 [0.114, 0.185] (3) \\
 & reader, partner tower & 0.202 [0.170, 0.239] (3) & 0.425 [0.390, 0.459] (6) & 0.432 [0.398, 0.464] (6) \\
 & tower, own reader & 0.295 [0.261, 0.337] (5) & 0.497 [0.453, 0.540] (6) & 0.408 [0.368, 0.450] (6) \\
 & tower, partner reader & 0.200 [0.175, 0.227] (4) & 0.335 [0.307, 0.364] (5) & 0.298 [0.273, 0.326] (5) \\
 & \quad mean reader effect & 0.137 [0.119, 0.161] & 0.327 [0.298, 0.354] & 0.288 [0.265, 0.315] \\
 & \quad mean tower effect & 0.248 [0.228, 0.271] & 0.416 [0.392, 0.441] & 0.353 [0.331, 0.379] \\
\addlinespace
MedGemma 4B, NIH & reader, own tower & 0.240 [0.221, 0.259] (5) & 0.210 [0.192, 0.233] (4) & 0.286 [0.268, 0.309] (6) \\
 & reader, partner tower & 0.133 [0.106, 0.169] (1) & 0.158 [0.136, 0.187] (4) & 0.425 [0.393, 0.458] (5) \\
 & tower, own reader & 0.160 [0.125, 0.193] (6) & 0.265 [0.239, 0.299] (5) & 0.496 [0.462, 0.525] (6) \\
 & tower, partner reader & 0.164 [0.140, 0.189] (4) & 0.061 [0.046, 0.080] (1) & 0.137 [0.120, 0.164] (4) \\
 & \quad mean reader effect & 0.186 [0.170, 0.207] & 0.184 [0.169, 0.203] & 0.356 [0.339, 0.376] \\
 & \quad mean tower effect & 0.162 [0.140, 0.182] & 0.163 [0.149, 0.185] & 0.316 [0.299, 0.338] \\
\bottomrule
\end{tabular}
\end{table}

\paragraph{Identical-tensor replay.} Checkpoints within one family share a vision tower, but two readers' inputs agree only up to bf16 rounding. The replay module removes this rounding difference. It stores the consumed block's token tensors once from the group's source checkpoint and feeds them to every reader in the group, replacing each reader's own tower output so the two readers consume the same bits. Before replacement, it measures drift as the largest absolute difference between the reader's own tower output and the stored tensor on the same rows. A reader replaying the tensors of its own block controls for the machinery: its replayed grade is its own grade computed through the replacement path. The module also recomputes each pair's between-reader difference twice: once with each reader running its own tower and once with both reading the stored tensors. The two computations differ only in their input. Across the \cfReplaySelf{} replays in which a reader reads its own block's tensors, no verdict and no ownership grade changes; \cfReplaySelfExact{} reproduce the block's own grade exactly and the third differs by at most $\cfReplayMaxDWSelf$ in a write magnitude. In \cfReplayDriftZero{} of the \cfReplayBlocks{} replays the reader's own tower output already matches the stored tensor bit for bit. The exception is LLaVA-Med, whose fine-tuned tower departs from it by up to \cfReplayDriftMax{}, against a mean absolute tensor entry of \cfReplayTensorMin{}--\cfReplayTensorMax{}. Across the \cfReplayBlocks{} replays, the change is distinguishable from zero in \cfReplayMoved{} of \cfReplayCells{} concept cells, the largest change in a write magnitude is $\cfReplayMaxDW$, and \cfReplayVerdictChanges{} verdicts and \cfReplayOwnedChanges{} ownership grades change. Across the \cfReplayPairs{} reader pairs, the mean absolute difference between the readers' own-question write magnitudes is \cfReplayPairOwnMin{}--\cfReplayPairOwnMax{} when each reader runs its own tower, and feeding both readers the stored tensors changes that difference by at most \cfReplayChangeMax{}, or \cfReplayRelMax\% of its size.

\paragraph{Semantic endpoints.} The four endpoints change what is read after the write, not what is written. The negated question asks whether there is no evidence of the finding. A direction that carries the concept raises the affirmative answer and lowers the negated one, so the two raw effects must have opposite signs. The forced choice asks which is present: the finding or its strongest competitor, taken from the block's own primary-template write matrix. We ask it in both presentation orders, so a position preference appears as a gap between them. The report continuation requests the findings section and reads the log-probability of the finding word. Each endpoint carries its own clean baseline, a coordinate-permutation sham of each direction under test, and its own family of \cfSemLedgerRandom{} random directions written at the same dose. An endpoint cell meets the steering reference when its own signed effect is positive, above every one of those \cfSemLedgerRandom{} random writes, and above the absolute sham, and it is owned when it also holds the campaign's simultaneous $6\times5$ max-$T$ advantage. The random bar here is the maximum of a family of \cfSemLedgerRandom{} rather than the 95th percentile of the campaign's 119; everything else is the campaign's rule (Table~\ref{tab:cf-ledger}). The incremental-validity test regresses the endpoint effect on write magnitude, probe selectivity, and clean-answer AUROC, with endpoint dummies, then adds the ownership contrast and its owned indicator. The penalty and the standardisation are fitted inside each training fold, so no held-out cell enters its own prediction. The predictors are per-concept block-level quantities, so the effective sample is the six concepts rather than the cells. The patient bootstrap preserves the dependence between cells that share rows. We cross-validate both models twice, leaving out one concept at a time inside a block and one block at a time across the \cfSemCvBlocks{} blocks. Held-out $R^2$ changes by $\cfSemCvLocoDelta$ under the concept split and $\cfSemCvLoboDelta$ under the block split. A null that permutes ownership across the concepts of a block over \cfSemCvPerms{} draws puts both at $p=\cfSemCvLocoP$ and $p=\cfSemCvLoboP$, and no single split reaches $p<0.05$. Table~\ref{tab:cf-semend} reports every endpoint, test, and regression.

\begin{table}[htbp]
\centering
\scriptsize
\setlength{\tabcolsep}{3pt}
\caption{\textbf{Four endpoints no direction was fitted against.} Both direction families of Section~\ref{sec:ansdir} written at the primary dose and read at the negated question, the two presentation orders of a forced choice, and a report continuation. The bottom panels add the ownership contrast to a regression of the endpoint effect and cross-validate it by concept and by block.}
\label{tab:cf-semend}
\begin{tabular}{lllcccc}
\toprule
\multirow{2}{*}{Checkpoint} & \multirow{2}{*}{Dataset} & \multirow{2}{*}{Endpoint} & \multicolumn{2}{c}{label direction} & \multicolumn{2}{c}{answer direction}\\
\cmidrule(lr){4-5}\cmidrule(lr){6-7}
&  &  & \multicolumn{1}{c}{owned} & \multicolumn{1}{c}{reference} & \multicolumn{1}{c}{owned} & \multicolumn{1}{c}{reference}\\
\midrule
Lingshu 32B & CheXpert Plus & negated question & 2/6 & 2/6 & 3/6 & 4/6 \\
Lingshu 32B & CheXpert Plus & forced choice, finding first & 3/6 & 3/6 & 4/6 & 4/6 \\
Lingshu 32B & CheXpert Plus & forced choice, competitor first & 4/6 & 4/6 & 5/6 & 5/6 \\
Lingshu 32B & CheXpert Plus & report continuation & 0/6 & 0/6 & 0/6 & 0/6 \\
Qwen2.5-VL-7B & NIH ChestX-ray14 & negated question & 0/6 & 1/6 & 3/6 & 6/6 \\
Qwen2.5-VL-7B & NIH ChestX-ray14 & forced choice, finding first & 0/6 & 0/6 & 3/6 & 3/6 \\
Qwen2.5-VL-7B & NIH ChestX-ray14 & forced choice, competitor first & 0/6 & 0/6 & 2/6 & 2/6 \\
Qwen2.5-VL-7B & NIH ChestX-ray14 & report continuation & 0/6 & 0/6 & 0/6 & 1/6 \\
Qwen2.5-VL-7B & CheXpert Plus & negated question & 1/6 & 1/6 & 5/6 & 6/6 \\
Qwen2.5-VL-7B & CheXpert Plus & forced choice, finding first & 0/6 & 0/6 & 3/6 & 4/6 \\
Qwen2.5-VL-7B & CheXpert Plus & forced choice, competitor first & 0/6 & 0/6 & 2/6 & 2/6 \\
Qwen2.5-VL-7B & CheXpert Plus & report continuation & 0/6 & 0/6 & 0/6 & 0/6 \\
\midrule
\textit{all} & \textit{3 blocks} & \textit{4 endpoints} & 10/72 & 11/72 & 30/72 & 37/72 \\
\midrule
\multicolumn{7}{l}{\textit{The three endpoint-specific tests}} \\
Checkpoint & Dataset & opposite on the negation & order gap & forced choice owned & report owned & core owned \\
\midrule
Lingshu 32B & CheXpert Plus & 4/6 & 0.026 & 3/6 & 0/6 & 5/6 \\
Qwen2.5-VL-7B & NIH ChestX-ray14 & 6/6 & 0.089 & 0/6 & 0/6 & 0/6 \\
Qwen2.5-VL-7B & CheXpert Plus & 6/6 & 0.111 & 0/6 & 0/6 & 0/6 \\
\midrule
\multicolumn{7}{l}{\textit{What ownership adds, held out: leaving out one concept at a time}} \\
Checkpoint & Dataset & Direction & held-out $R^2$ & with ownership & $\Delta R^2$ & permutation $p$ \\
\midrule
Lingshu 32B & CheXpert Plus & label direction & -0.034 & 0.147 & 0.181 & 0.39 \\
Lingshu 32B & CheXpert Plus & answer direction & 0.069 & -0.223 & -0.291 & 0.70 \\
Qwen2.5-VL-7B & CheXpert Plus & label direction & -0.139 & -0.149 & -0.010 & 0.55 \\
Qwen2.5-VL-7B & CheXpert Plus & answer direction & -3.168 & -0.112 & 3.057 & 0.18 \\
Qwen2.5-VL-7B & NIH ChestX-ray14 & label direction & -0.195 & -0.147 & 0.048 & 0.31 \\
Qwen2.5-VL-7B & NIH ChestX-ray14 & answer direction & -1.310 & -0.176 & 1.134 & 0.50 \\
\textit{pooled} & \textit{6 tests} & \textit{both} & -- & -- & 0.687 & 0.26 \\
\midrule
\multicolumn{7}{l}{\textit{The same, leaving out one block at a time}} \\
\midrule
\textit{all 3} & \textit{pooled} & label direction & -0.578 & -0.500 & 0.078 & 0.30 \\
\textit{all 3} & \textit{pooled} & answer direction & -0.461 & -0.398 & 0.063 & 0.26 \\
\textit{pooled} & \textit{2 tests} & \textit{both} & -- & -- & 0.070 & 0.26 \\
\midrule
\multicolumn{7}{l}{\textit{In-sample reference only, since adding a predictor cannot lower a training $R^2$}} \\
Checkpoint & Dataset & Direction & training $R^2$ & with ownership & $\Delta R^2$ & \\
\midrule
Lingshu 32B & CheXpert Plus & label direction & 0.769 & 0.788 & 0.019 & \\
Lingshu 32B & CheXpert Plus & answer direction & 0.368 & 0.423 & 0.055 & \\
Qwen2.5-VL-7B & NIH ChestX-ray14 & label direction & 0.241 & 0.246 & 0.004 & \\
Qwen2.5-VL-7B & NIH ChestX-ray14 & answer direction & 0.267 & 0.392 & 0.125 & \\
Qwen2.5-VL-7B & CheXpert Plus & label direction & 0.175 & 0.250 & 0.074 & \\
Qwen2.5-VL-7B & CheXpert Plus & answer direction & 0.473 & 0.509 & 0.036 & \\
\bottomrule
\end{tabular}
\end{table}

\paragraph{Fine-grained objects and difficulty matching.} The rule that picks the six additional COCO categories, and the two matching windows, are stated in full in Appendix~\ref{app:repro-protocol}. The rule reads the label counts and the mean instance areas of the fixed cohorts and nothing else; the registered protocol file carries it under the date of amendment A15, and the module's first outcome shard is later than that date, so the categories and the windows were fixed before any of them was scored and before any ownership outcome of theirs existed. We fit, lift, write, and grade these categories exactly like the six easy objects, in the same blocks and on the same rows, as their own six-direction family, against a random family and a sham of their own (Table~\ref{tab:cf-ledger}). We measure their clean-answer AUROC on the same calibration rows and by the same rule as every other cell, giving the matching variable the same meaning on both sides. The matched comparison pairs each chest cell with every natural-image cell whose matching variables lie within \cfFgWindow{} of it. It reports the ownership rate on each side of the surviving pairs under two estimators. The \emph{pooled} difference is the ownership rate among the matched chest cells minus the rate among the natural-image cells that served as a partner at least once. The \emph{matched} difference is the mean over matched chest cells of that cell's own ownership minus the ownership rate of its own partner set, so a chest cell with many partners and one with few count equally. Neither estimator can be obtained from the other. The \emph{chest cells} and \emph{partners} columns give the effective sample of both: matched chest cells of all chest cells, and distinct natural-image cells used as a partner of all of them. Both intervals are percentile intervals over 5{,}000 draws of a cluster bootstrap that resamples models, keeping each model's chest and natural-image cells together. Chest cells without a partner leave the comparison. The stratified comparison answers the same question without dropping a cell. It bins every cell of both pools by its clean-answer AUROC at the edges \cfStratEdges{}, fixed in the released analysis code, and reports the ownership rate inside each bin with the same cluster bootstrap over models on the difference in the top bin. The last bin is closed on the right and every cell of each pool falls in exactly one bin. The \cfStratChestTopN{} chest cells of that bin are CheXpert Plus cells; no NIH ChestX-ray14 cell reaches a clean-answer AUROC of \cfStratTopEdge{}. The same stratification of the nine-direction attribute family appears in Table~\ref{tab:cf-attr}. Table~\ref{tab:cf-fgobj} reports the per-block counts, the three pools, each matching, and every bin.

\begin{table}[htbp]
\centering
\scriptsize
\setlength{\tabcolsep}{4pt}
\setlength{\abovecaptionskip}{2pt}
\setlength{\belowcaptionskip}{0pt}
\renewcommand{\arraystretch}{0.95}
\caption{\textbf{Small and fine-grained objects, and the difficulty-matched comparison.} Top and middle: the six fine-grained categories graded like the six easy objects, against the easy objects of every COCO block and of the 6 blocks the fine-grained cells come from. Bottom: chest cells matched to natural-image cells on difficulty, then the same cells stratified on clean-answer AUROC.}
\label{tab:cf-fgobj}
\begin{tabular}{lrrrrrr}
\toprule
\multirow{2}{*}{Checkpoint} & \multicolumn{2}{c}{owned of six} & \multicolumn{2}{c}{median answer AUROC} & \multicolumn{2}{c}{median selectivity}\\
\cmidrule(lr){2-3}\cmidrule(lr){4-5}\cmidrule(lr){6-7}
& \multicolumn{1}{c}{fine} & \multicolumn{1}{c}{easy} & \multicolumn{1}{c}{fine} & \multicolumn{1}{c}{easy} & \multicolumn{1}{c}{fine} & \multicolumn{1}{c}{easy}\\
\midrule
Qwen2.5-VL-3B & 6 & 6 & 0.947 & 0.971 & 0.024 & 0.037 \\
Qwen3-VL-4B & 5 & 6 & 0.984 & 0.990 & 0.036 & 0.033 \\
InternVL3.5-8B & 6 & 6 & 0.981 & 0.978 & 0.104 & 0.117 \\
Gemma 3 4B & 4 & 5 & 0.948 & 0.924 & 0.076 & 0.070 \\
Lingshu 7B & 6 & 6 & 0.950 & 0.968 & 0.025 & 0.023 \\
LLaVA-1.5-7B & 3 & 5 & 0.956 & 0.979 & 0.266 & 0.285 \\
\midrule
\multicolumn{7}{l}{\textit{Paired within these blocks:} fine-grained minus easy is +0, -1, +0, -1, -2, +0 cells of six; fine higher in 0,} \\
\multicolumn{7}{l}{easy higher in 3, tied in 3; exact two-sided sign test $p=0.25$} \\
\midrule
\multicolumn{7}{l}{\textit{The four pools of cells}} \\
Cells & blocks & owned & rate & median answer AUROC & median selectivity & readable \\
\midrule
chest findings & 50 & 38/300 & 0.127 & 0.695 & 0.137 & 178/300 \\
easy objects, every COCO block & 25 & 113/150 & 0.753 & 0.977 & 0.064 & 90/150 \\
easy objects, these blocks & 6 & 34/36 & 0.944 & 0.980 & 0.065 & 24/36 \\
fine-grained objects, these blocks & 6 & 30/36 & 0.833 & 0.963 & 0.058 & 23/36 \\
\midrule
\multicolumn{7}{l}{\textit{Chest cells matched to natural-image cells on difficulty, under both estimators}} \\
Matching & chest cells & partners & chest rate & natural rate & difference & 95\% interval \\
\midrule
probe selectivity alone & 300/300 & 186/186 & 0.127 & 0.769 &  &  \\
\quad pooled &  &  &  &  & -0.642 & [-0.763, -0.512] \\
\quad matched &  &  &  &  & -0.539 & [-0.678, -0.359] \\
clean-answer AUROC alone & 300/300 & 186/186 & 0.127 & 0.769 &  &  \\
\quad pooled &  &  &  &  & -0.642 & [-0.749, -0.505] \\
\quad matched &  &  &  &  & -0.066 & [-0.304, -0.004] \\
both variables & 212/300 & 101/186 & 0.137 & 0.624 &  &  \\
\quad pooled &  &  &  &  & -0.487 & [-0.653, -0.201] \\
\quad matched &  &  &  &  & -0.024 & [-0.198, 0.076] \\
both variables, easy partners only & 211/300 & 79/150 & 0.137 & 0.570 &  &  \\
\quad pooled &  &  &  &  & -0.432 & [-0.625, -0.150] \\
\quad matched &  &  &  &  & -0.008 & [-0.198, 0.085] \\
\midrule
\multicolumn{7}{l}{\textit{Chest cells and natural-image cells stratified on clean-answer AUROC, owned of cells}} \\
Cells & $<$ 0.50 & 0.50--0.70 & 0.70--0.80 & 0.80--0.90 & $\geq$ 0.90 & all cells \\
\midrule
NIH ChestX-ray14 & 0/7 & 6/80 & 3/43 & 4/19 & 0/1 & 13/150 \\
CheXpert Plus & 0/13 & 6/52 & 2/20 & 6/29 & 11/36 & 25/150 \\
\textit{chest findings} & 0/20 & 12/132 & 5/63 & 10/48 & 11/37 & 38/300 \\
easy objects, every COCO block & 0/6 & 0/6 & 0/7 & 1/4 & 112/127 & 113/150 \\
fine-grained objects, these blocks & 0/0 & 0/0 & 0/0 & 1/1 & 29/35 & 30/36 \\
\textit{natural-image objects} & 0/6 & 0/6 & 0/7 & 2/5 & 141/162 & 143/186 \\
\multicolumn{7}{l}{\quad chest minus natural-image in the $\geq$ 0.90 bin: -0.573, 95\% interval [-0.763, -0.388]} \\
\bottomrule
\end{tabular}
\end{table}

\subsection{Write structure across the grid}
\label{app:cf-figures}

Figure~\ref{fig:cf-w-structure} summarises the $6\times6$ write matrices across all completed checkpoints as the cell-wise median matrix for each dataset. Figure~\ref{fig:write-flow} reads the same matrices as flow: band width is proportional to $P_{q,d}$, the mean positive answer change across checkpoints, coloured bands reach the writer's own question, and grey bands reach other questions. The write-flow share is the total width of own-question bands divided by the total width of all bands, $\sum_d P_{d,d}/\sum_{q,d}P_{q,d}$: \cfOwnShareNih\% on NIH and \cfOwnShareChex\% on CheXpert, compared with \cfOwnShareCoco\% on COCO. The column-selectivity index $S_d$ in Table~\ref{tab:cf-validation} measures a different quantity. It retains the sign of a direction's own effect. Its denominator includes negative changes. It gives every cell equal weight. We report its median over cells (\cfValColSelNih{}, \cfValColSelChex{}, and \cfValColSelCoco{}). The largest sinks on the chest sets are the Cardiomegaly question on NIH and the Edema and Effusion questions on CheXpert. These questions receive writes from every direction.

\begin{figure}[p]
    \centering
    \resizebox{\textwidth}{!}{\includegraphics{figures/figA1_write_structure.pdf}}
    \caption{\textbf{Write structure across the grid.} Cell-wise median of the write matrix $W_{q,d}$ over completed checkpoints on NIH ChestX-ray14, CheXpert Plus, and COCO. Figure~\ref{fig:write-flow} reads the same matrices as flow.}
    \label{fig:cf-w-structure}
\end{figure}

\subsection{Per-image examples}
\label{app:cf-examples}

Figure~\ref{fig:cf-examples} shows grades for individual test images, with $P(\mathrm{yes})$ in parentheses on every answer line. The first four chest examples are deficit cells of Qwen2.5-VL-7B, the next two are owned cells (Lingshu 32B Edema, Qwen2.5-VL-72B Atelectasis), and the natural-image examples are owned cells of Qwen2.5-VL-7B, Lingshu 32B, and MedGemma 4B. For a chest example, we pool test rows where the finding is labelled present, the probe reads it (probability above $0.5$), and the clean answer is No. For a natural-image example, we pool rows where the probe reads the label correctly and the clean answer is No. We split each pool by whether the concept's own write or the strongest competing write raises the answer more. Among rows where either write raises the answer by at least $0.2$, we show the row with the gap closest to the median. Each example's last line reports how many pool rows share its ordering, giving the example its own base rate.

\begin{figure}[p]
    \centering
    \resizebox{\textwidth}{!}{\includegraphics{figures/figA2_examples.pdf}}
    \caption{\textbf{Reading, answering, and writing on single images.} Per image: the question, the report label with the probe's reading, the clean answer, and the answers under the concept's own write and the strongest competing write, at $\alpha=+0.25$. A check marks an image the own direction wins, a cross one a competitor wins.}
    \label{fig:cf-examples}
\end{figure}
\section{Seed study: registered protocol}
\label{app:seed}

The two-model study that motivated the evaluation registered the procedures below; Appendix~\ref{app:seed-results} reports its results.

\paragraph{Data and split.}
We use NIH ChestX-ray14 \citep{wang2017chestxray}, whose findings are mined from radiology reports. A deterministic hash of patient identity assigns rows to disjoint train, validation, and test partitions. The registered manifest contains 18,212 training images and 5,343 test images from 1,659 test patients. The three cells pair LLaVA-1.5-7B with Effusion and Edema, and Qwen2.5-VL-7B with Effusion. All reported uncertainty clusters by patient. Prompts appear in Appendix~\ref{app:protocol}.

\paragraph{Models and consumed loci.}
For LLaVA-1.5-7B \citep{liu2024improved} and Qwen2.5-VL-7B \citep{bai2025qwen25vl}, we use the final visual block consumed by the connector or merger. Writing $\alpha$ for the signed intervention dose, synthetic hook tests establish downstream-path membership, a bitwise $\alpha=0$ no-op, and changed downstream representations and logits at nonzero dose. Probe and intervention therefore share one consumed representation.

\paragraph{Controlled decodability.}
For image $i$, let $\hvec_{it}\in\R^D$ denote the $D$-dimensional representation of visual token $t$ at the consumed final block; mean pooling over tokens yields $\hvec_i$. A fixed seed-0 Gaussian matrix $R\in\R^{D\times512}$ projects every model to a common readout dimension; featurewise scale $\boldsymbol{s}$ is fitted on training rows only. We then fit logistic regression with $C=1$, maximum 2,000 iterations, and seed 0. Each of 20 control seeds maps the recurring input type (AP view indicator, sex indicator, and age decade, giving 39 strata) to a random label. Let $y$ be the real test labels and $r$ the fitted probe's continuous scores. For control seed $k\in\{0,\ldots,19\}$, $y_k^{\mathrm{ctrl}}$ denotes the random labels and $r_k^{\mathrm{ctrl}}$ the corresponding control probe scores. Selectivity compares their AUROCs:
\begin{equation}
S = \operatorname{AUROC}(y,r)-\frac{1}{20}\sum_{k=0}^{19}\operatorname{AUROC}(y_k^{\mathrm{ctrl}},r_k^{\mathrm{ctrl}}).
\label{eq:seed-selectivity}
\end{equation}
We report the real AUROC, control distribution, and $S$ with 2,000 patient-cluster bootstrap resamples. A positive interval for $S$ means that the target is more readable than same-capacity nuisance-type label maps. It does not yet imply downstream use. Paired contrasts reuse identical ordered test rows and bootstrap draws.

\paragraph{Probe-normal intervention.}
Let $\boldsymbol\beta_c$ be the projected-space logistic coefficient for the positive class of concept $c$. We lift it to the consumed $D$-dimensional representation and normalize it exactly as
\begin{equation}
\hat\wvec_c = \frac{R\left(\boldsymbol\beta_c\oslash\max(\boldsymbol{s},10^{-8})\right)}{\left\lVert R\left(\boldsymbol\beta_c\oslash\max(\boldsymbol{s},10^{-8})\right)\right\rVert_2}.
\label{eq:seed-lift}
\end{equation}
The positive-class coefficient fixes the direction's clinical pole. At every visual token $t$, the registered intervention is
\begin{equation}
\hvec'_{it}=\hvec_{it}+\alpha\lVert\hvec_{it}\rVert_2\hat\wvec_c.
\label{eq:seed-intervention}
\end{equation}
Here the sign of $\alpha$ selects the positive or negative clinical pole, while $|\alpha|$ is the perturbation norm as a fraction of each token's activation norm. At the first answer position, $P(\mathrm{yes})=\sigma(\ell_{\mathrm{yes}}-\ell_{\mathrm{no}})$, where each log probability is the maximum over case and leading-space variants that tokenize to one token; the endpoint is its paired mean change from the common unsteered baseline. Original cells use concept doses $\{-1,-0.5,-0.25,-0.1,0,0.1,0.25,0.5,1\}$ on 200 held-out test images. Controls are evaluated at the six nonzero doses excluding $\pm0.1$: 20 isotropic random unit directions, one sham formed by a fixed coordinate permutation of $\hat\wvec_c$ and renormalization, and fixed alternative clinical probe normals. Every control receives the same signed dose on the same rows. Dose monotonicity over $[-0.5,0.5]$ is descriptive.

\paragraph{Original-cell decision rule.}
Let $\Delta_c(\alpha)$ be the target direction's paired mean probability change, and let $\mathcal A=\{-1,-0.5,-0.25,0.25,0.5,1\}$ be the six controlled doses. The primary effect is $\max_{\alpha\in\mathcal A}\operatorname{sign}(\alpha)\Delta_c(\alpha)$, attained at dose $\alpha^\star$. A cell clears the generic intervention gate when this selected effect exceeds the 95th percentile of the 20 random directions' mean probability changes at that same $\alpha^\star$, each multiplied by $\operatorname{sign}(\alpha^\star)$, and the maximum absolute sham effect over $\mathcal A$. The random directions are not separately maximized over doses. This asymmetric selection defines a registered reference gate, not a cross-dose-corrected 5\% test. The conjunction is fixed across the three cells. The independent supplement below instead evaluates a dose locked before observing its patients.

\paragraph{Direction-specificity supplement.}
The initial experiment prespecified the random and sham comparisons; its clinical-direction analysis was descriptive. The Nodule response motivated a separately registered supplement at the locked Qwen dose, $\alpha=+0.25$. We select one row by label-blind hash for each of 400 ChestX-ray14 test patients absent from the original 200-image intervention cohort. The semantic family contains all five non-target normals from the pre-existing Qwen direction bundle (Atelectasis, Pneumothorax, Cardiomegaly, Mass, and Nodule), rather than directions chosen from the supplement outcome. Each uses the same training rows, projection, scaling, classifier, lifting rule, and normalization as Effusion. Together with a common baseline, 20 random directions, and the sham, the grid produces 28 configurations and 11,200 per-image outcomes; exact row selection appears in Appendix~\ref{app:seed-results}.

Semantic ownership is a contrast on a direction--question response surface, not a property of a direction in isolation. Let $W_{q,d}$ denote the paired mean change in question $q$'s yes probability after intervening with direction $d$ at the locked dose. A large $W_{q,d}$ establishes that $d$ can write $q$'s answer; it does not establish that $d$ owns that answer. Within a fixed clinical family, the ownership contrast is
\begin{equation}
O_q=W_{q,q}-\max_{d\neq q}W_{q,d}.
\label{eq:seed-margin}
\end{equation}
The registered supplement estimates the Effusion row of this surface, so the reported ownership contrast $O_{\mathrm{Effusion}}$ is Effusion minus its maximum alternative. The maximum is recomputed inside each of 5,000 patient-bootstrap draws. Thus $O_{\mathrm{Effusion}}>0$ means Effusion beats every member of the fixed clinical family, while $O_{\mathrm{Effusion}}<0$ means at least one alternative is stronger. Adding competitors can only maintain or decrease $O_q$: a positive advantage is conditional on the candidate family, whereas a negative counterexample does not require an exhaustive family. This operational ownership criterion measures relative advantage, not a necessary condition for a concept to participate in computation. Direction specificity requires a positive Effusion effect above random p95 and absolute sham, together with a one-sided 95\% lower bound for $O_{\mathrm{Effusion}}$ above zero.

\paragraph{Causal-ownership matrix.}
To test whether another clinical normal better matches the answer pathway, we prospectively cross six Qwen directions (Effusion, Atelectasis, Pneumothorax, Cardiomegaly, Mass, and Nodule) with their six matching questions on 400 new patients. Every question uses the same $\alpha=+0.25$, 119 random directions, and a sham. For question $q$, ownership requires the diagonal effect $W_{q,q}$ to exceed every other clinical direction, every random direction, and the absolute sham effect. A centered studentized max-$T$ bootstrap gives simultaneous 95\% lower bounds over all 30 clinical contrasts. The shared-alias family contains two statistics per direction: mean excess over the matching diagonals on its five off-diagonal questions and signed mean effect on those questions. These give 12 max-$T$ bounds; the effect must also exceed the random alias maximum and maximum absolute sham effect, detailed in Appendix~\ref{app:seed-results}.

For a family of $m$ statistics, let $\widehat C_j$ be statistic $j$'s estimate and $C^*_{bj}$ its value in joint patient-bootstrap draw $b$, with $B=5{,}000$ draws. Write $s_j=\operatorname{SD}_{b}(C^*_{bj})$ for the bootstrap standard deviation with denominator $B-1$, held fixed across draws, and $Q_{.95}$ for the empirical 95th percentile with linear interpolation. For $s_j>10^{-12}$, the centered standardized draw $Z_{bj}$, common critical value $c_{.95}$, and simultaneous lower bound $L_j$ are
\begin{equation}
Z_{bj}=\frac{C^*_{bj}-\widehat C_j}{s_j},\qquad
c_{.95}=Q_{.95}\!\left(\left\{\max_{1\leq j\leq m}Z_{bj}\right\}_{b=1}^{B}\right),\qquad
L_j=\widehat C_j-c_{.95}s_j.
\label{eq:max-t}
\end{equation}
Joint resampling preserves paired conditions; no standard error is re-estimated within a draw. Appendix~\ref{app:simultaneous-bounds} specifies the families, degenerate components, and the separate percentile intervals for target-minus-maximum margins.

\paragraph{Answer-encoding crossover.}
We next register a paired sensitivity study on 200 patients from the causal-ownership matrix cohort; the other 200 form a patient-disjoint capability-calibration set. This is not the earlier Effusion specificity cohort. Six questions cross the original yes/no prompt with A-present and B-present prompts, which exchange the finding-to-letter mapping. Each uses both $\alpha=\pm0.25$, six clinical directions, 20 random directions, and a sham. A question enters the primary endpoint only if both clean coded mappings have AUROC one-sided 95\% lower bounds above 0.5 and at least ten patients per class on calibration. At $\alpha=+0.25$, we compare squared mean effects in components that retain or reverse raw A-minus-B sign under the mapping exchange, using an unbiased estimator and 2,000 joint patient-bootstrap draws. Appendix~\ref{app:seed-results} defines the components and full control rule. Descriptive clinical competition uses the contrast in Equation~\ref{eq:seed-margin}, replacing probability change with finding-present logit-margin change within each prompt.

\paragraph{Independent Mass confirmation.}
We prospectively test Mass on 200 patients absent from all preceding Qwen intervention cohorts, retaining the crossover's 200 calibration patients. Two question wordings (\textit{show}, \textit{is}) cross explicit yes/no, A-present, and B-present instructions; the original yes/no prompt is a seventh anchor. At $\alpha=+0.25$, each cell receives six clinical directions, 119 random directions, and a sham. The four coded cells share 20 clinical comparisons with simultaneous 95\% lower bounds. Confirmation requires positive bounds and Mass effects above every random effect and absolute sham, either across all four coded cells or both original-wording cells. Appendix~\ref{app:seed-results} specifies calibration and paired descriptive contrasts.

\section{Seed study: results}
\label{app:seed-results}

\begin{table}[t]
\centering
\scriptsize
\setlength{\tabcolsep}{4pt}
\input{tables/cf_colors}
\caption{\textbf{Seed replication: Qwen2.5-VL-7B on NIH ChestX-ray14, 600 new patients.} One row per concept, with the concept write $W_{qq}$ against the random 95th percentile (p95), the absolute sham, and the 120 compared effects it is ranked among. Blue shades a negative ownership contrast, darker for larger magnitude.}
\label{tab:cf-seed}
\begin{tabular}{lrrrrrrrll}
\toprule
Concept & \multicolumn{1}{c}{$S$} & \multicolumn{1}{c}{AUROC} & \multicolumn{1}{c}{$W_{qq}$} & \multicolumn{1}{c}{p95} & \multicolumn{1}{c}{$|$sham$|$} & \multicolumn{1}{c}{rank} & \multicolumn{1}{c}{$O_q$ [95\% CI]} & Competitor & Verdict\\
\midrule
Effusion & 0.102 & 0.593 & +0.190 & 0.137 & 0.017 & 2/120 & \cellcolor{cfSlate2}-0.071 [-0.080, -0.062] & Nodule & competitor \\
Atelectasis & 0.062 & 0.566 & +0.101 & 0.297 & 0.096 & 55/120 & \cellcolor{cfSlate3}-0.226 [-0.233, -0.219] & Nodule & competitor \\
Pneumothorax & 0.152 & 0.508 & +0.089 & 0.320 & 0.113 & 52/120 & \cellcolor{cfSlate3}-0.102 [-0.108, -0.096] & Effusion & competitor \\
Cardiomegaly & 0.205 & 0.742 & +0.085 & 0.311 & 0.158 & 33/120 & \cellcolor{cfSlate3}-0.219 [-0.227, -0.211] & Atelectasis & competitor \\
Mass & 0.021 & 0.678 & +0.359 & 0.372 & 0.092 & 10/120 & \cellcolor{cfSlate3}-0.169 [-0.178, -0.160] & Effusion & competitor \\
Nodule & -0.041 & 0.642 & +0.118 & 0.158 & 0.055 & 15/120 & \cellcolor{cfSlate2}-0.072 [-0.081, -0.063] & Effusion & competitor \\
\bottomrule
\end{tabular}
\end{table}

\subsection{A non-generic write effect reproduces}
\label{sec:qwen-replication}

The first finding starts with a successful intervention: Qwen Effusion reproduces beyond random and sham controls before clinical alternatives change its attribution.

Controlled decodability is evaluated at the final visual block each connector consumes: the penultimate CLIP encoder layer for LLaVA and the last of Qwen's 32 visual blocks. Intervals are 95\% patient-cluster bootstrap intervals over 2,000 draws; the control AUROC is the mean over the 20 recurring-type random-label maps, with its 5th--95th percentile spread in parentheses.

Qwen Effusion first satisfies the reader test. Its probe reaches AUROC \cfSeedDecQwenEffAuroc{} $[\cfSeedDecQwenEffAurocLo,\cfSeedDecQwenEffAurocHi]$ against a control AUROC of \cfSeedDecQwenEffCtrl{} (\cfSeedDecQwenEffCtrlLo--\cfSeedDecQwenEffCtrlHi), a controlled selectivity of \cfSeedDecQwenEffSel{} (95\% CI $[\cfSeedDecQwenEffSelLo,\cfSeedDecQwenEffSelHi]$). LLaVA-1.5-7B reads both of its concepts at its consumed block: Effusion at AUROC \cfSeedDecLlavaEffAuroc{} $[\cfSeedDecLlavaEffAurocLo,\cfSeedDecLlavaEffAurocHi]$ with selectivity \cfSeedDecLlavaEffSel{} $[\cfSeedDecLlavaEffSelLo,\cfSeedDecLlavaEffSelHi]$, and Edema at AUROC \cfSeedDecLlavaEdemaAuroc{} $[\cfSeedDecLlavaEdemaAurocLo,\cfSeedDecLlavaEdemaAurocHi]$ with selectivity \cfSeedDecLlavaEdemaSel{} $[\cfSeedDecLlavaEdemaSelLo,\cfSeedDecLlavaEdemaSelHi]$, against a control AUROC of \cfSeedDecLlavaEffCtrl{} (\cfSeedDecLlavaEffCtrlLo--\cfSeedDecLlavaEffCtrlHi). Original-cohort intervention outcomes maximize the concept-consistent effect over six fixed nonzero doses. Qwen makes a strong write-direction case in the original 200-image cohort: the maximum concept-consistent change is 0.2343 at $\alpha=+0.25$, above the random-direction 95th percentile of 0.0963 and the maximum absolute sham effect of 0.0852.

The response at the locked dose $\alpha=+0.25$ reproduces on the patient-disjoint 400-patient cohort. Effusion changes mean $P(\mathrm{yes})$ by 0.1945, again above the random-direction 95th percentile of 0.0878 and absolute sham effect of 0.0157. None of these patients appears in the original intervention cohort, and the dose was not reselected, supporting a large and reproducible non-generic effect within ChestX-ray14.


\subsection{Semantic competition changes attribution, not efficacy}
\label{sec:specificity-results}

The fixed clinical family provides a harder identification test than non-semantic controls. Nodule changes the Effusion answer by 0.2519 and realizes the point-estimate maximum among the five alternatives. The Effusion ownership contrast $O_{\mathrm{Effusion}}$ is $-0.0573$. The registered fixed-family endpoint recomputes the maximum inside each of 5,000 patient-bootstrap draws, giving a two-sided 95\% CI of $[-0.0694,-0.0450]$ and a one-sided 95\% lower bound of $-0.0678$. The entire interval is below zero. Descriptively, all five alternative clinical normals move the Effusion endpoint in the same positive direction; Effusion exceeds four by 0.1172--0.1778, while Nodule alone has a larger point estimate. This shared-sign pattern motivates an answer-sensitive-alias hypothesis but cannot distinguish it from question-specific ownership.

Qwen Effusion therefore remains a reproducible direction that changes the answer beyond random and sham perturbations, while its specific semantic attribution fails. The familywise endpoint establishes that at least one fixed clinical alternative is stronger; Nodule is the point-estimate maximizer, not a separately familywise-identified semantic owner.

\subsection{Wording changes clinical competition within the answer interface}
\label{sec:encoding-results}

The second finding concerns the conditions of clinical competition. The Qwen ownership pattern persists in logit coordinates, so we test whether changing the answer interface changes a direction's relative advantage.

In the registered crossover, Atelectasis, Pneumothorax, Mass, and Nodule pass independent capability calibration. On these four questions, mapping-even minus mapping-odd response energy is $-2.1786$ (95\% CI $[-2.4061,-1.9810]$): effects are larger in the component that reverses raw A-minus-B sign when the finding-to-letter mapping is exchanged. Yet all 20 random directions and the sham are also mapping-odd dominant. Aggregate mapping response therefore does not distinguish the clinical normals from generic perturbations.

The crossover identifies Mass as a prompt-conditioned candidate: its clinical margin changes from $-0.7181$ under the original yes/no prompt to $0.2825$ and $0.2900$ under A-present and B-present, respectively (Table~\ref{tab:seed-mass}, top). A-present assigns A to presence and B to absence; B-present reverses those assignments. Both coded effects exceed the discovery study's 20 random directions and sham. The top rows of Table~\ref{tab:seed-mass} carry pointwise 95\% intervals. Because wording and instructions also change, we test them separately on independent patients.

\begin{table}[H]
\centering
\scriptsize
\setlength{\tabcolsep}{5pt}
\caption{\textbf{Mass clinical competition across prompts.} Clinical margin per prompt: the Mass effect minus the largest of the five other clinical effects, in finding-present logit units. Top, the discovery crossover; bottom, the registered confirmation on new patients.}
\label{tab:seed-mass}
\begin{tabular}{lrrrrr}
\toprule
Prompt & Clinical margin & 95\% CI & Mass effect & Random max & $|$Sham$|$ \\
\midrule
Standard yes/no & $-0.7181$ & $[-0.7956,-0.6412]$ & 2.8625 & 3.0412 & 0.5869 \\
A-present & $0.2825$ & $[0.2269,0.3406]$ & 1.6125 & 1.1650 & 0.7550 \\
B-present & $0.2900$ & $[0.2300,0.3506]$ & 1.7469 & 1.2581 & 1.1475 \\
\midrule
Condition & Clinical margin & Minimum LCB & Mass effect & Random max & \\
\midrule
\textit{show}/A-present & 0.2525 & 0.1675 & 1.7019 & 2.2100 & \\
\textit{show}/B-present & 0.2338 & 0.1462 & 1.8188 & 2.2169 & \\
\textit{is}/A-present & $-0.1087$ & $-0.1751$ & 1.9225 & 2.4750 & \\
\textit{is}/B-present & $-0.2006$ & $-0.2686$ & 2.0019 & 2.4944 & \\
\bottomrule
\end{tabular}
\end{table}

The bottom of Table~\ref{tab:seed-mass} reports the registered confirmation on 200 new patients. Clinical margins subtract the largest of five alternative clinical effects in finding-present logit units. Minimum LCB is the smallest simultaneous lower bound among a cell's five contrasts, from the joint 20-comparison family. All seven prompt cells pass separate capability calibration.

The clinical advantage reproduces under both \textit{show} coded prompts, but the Mass effects remain below the maxima of 119 random directions, with reference ranks of 3/120. These 119 directions form a fixed reference family: each rank places Mass third among the 120 compared effects. This finite descriptive ranking is not a $p$-value, and failure to exceed this family's maximum does not establish nonspecificity across arbitrary direction distributions or reference-family sizes. Under \textit{is}, both clinical margins are negative. Neither the cross-wording nor original-wording specificity conjunction is met.

Explicit \textit{show} yes/no also has a positive clinical margin, 0.2369; its paired differences from A-present and B-present span zero. Within each response instruction, changing \textit{show} to \textit{is} reduces the clinical margin. These paired wording contrasts are descriptive. They support wording-conditioned clinical competition, rather than an advantage attributable to answer letters alone. Mass steering increases Brier error in every cell, while all AUROC-change intervals span zero.

\FloatBarrier
\section{Prompts, Scores, and Notation}
\label{app:protocol}

\subsection{A common question and answer endpoint}

At the first answer position, let $\ell_{\mathrm{yes}}$ and
$\ell_{\mathrm{no}}$ denote the maximum log probabilities among the valid
single-token case and leading-space variants of each answer.  The registered
score is $P(\mathrm{yes})=\sigma(\ell_{\mathrm{yes}}-\ell_{\mathrm{no}})$,
where $\sigma$ is the logistic sigmoid. The original Edema cell uses ``pulmonary
edema,'' and the closure experiment uses ``consolidation'' in the same
template. Taking the best valid token variant for each answer makes the
endpoint insensitive to which listed capitalization or leading-space form
receives the highest score. This endpoint measures the relative support for
yes versus no at the first answer position; it does not measure the accuracy
of a generated clinical report.


The shared interface makes the question--direction grid interpretable:
changing a clinical direction does not change the question being scored.
Differences across rows reflect different clinical questions (questions are
rows of $W$, directions are columns), so the ownership test compares directions
within each row before drawing a conclusion about the named concept.

\subsection{Reading the intervention notation}

Table~\ref{tab:notation} collects the notation introduced in the article.
The distinction between question $q$ and direction $d$ is central: a large
$W_{q,d}$ shows that a direction moves an answer, while $O_q$ asks whether
the direction bearing that answer's clinical name moves it more than its
competitors. A positive ownership margin alone is not the full confirmation
rule; the uncertainty and non-semantic control comparisons must also pass.
Keeping these roles separate allows the following analyses to distinguish
readability, intervention strength, and semantic attribution.

\begin{table}[H]
\centering
\caption{Notation used in the manuscript.}
\label{tab:notation}
\begin{tabular}{l>{\raggedright\arraybackslash}p{0.8\linewidth}}
\toprule
Symbol & Definition \\
\midrule
$q,d,c$ & Clinical question, steering direction, and probe concept, respectively. \\
$\hvec_{it}$ & Consumed visual representation for image $i$ and visual token $t$. \\
$R,\boldsymbol{s}$ & Fixed Gaussian projection and train-fitted featurewise scale. \\
$\hat\wvec_c$ & Unit probe-normal direction for concept $c$ after lifting to the consumed representation. \\
$\alpha$ & Signed intervention dose; $|\alpha|$ is the perturbation norm relative to each token norm. \\
$W_{q,d}$ & Paired mean change in question $q$'s $P(\mathrm{yes})$ under direction $d$. \\
$O_q$ & Ownership margin $W_{q,q}-\max_{d\neq q}W_{q,d}$; positive values favor the named direction. \\
$S$ & Real-label AUROC minus the mean AUROC of the 20 nuisance-type control tasks. \\
$\gamma_i$ & Positive-minus-negative paired representation displacement along the Consolidation normal. \\
\bottomrule
\end{tabular}
\end{table}

The notation separates three levels of evidence: $S$ describes label
availability, $W_{q,d}$ describes an intervention effect, and $O_q$ compares
the named direction with clinical alternatives. The remaining appendix
follows this progression while keeping the evaluation cohorts explicit.

\FloatBarrier

\section{Joint bootstrap and simultaneous lower bounds}
\label{app:simultaneous-bounds}

Equation~\ref{eq:max-t} uses the standard deviation across bootstrap
estimates as a fixed scale for each statistic. With $\overline C^*_j$
denoting their mean, that scale is
\begin{equation}
s_j=\left[\frac{1}{B-1}\sum_{b=1}^{B}
       (C^*_{bj}-\overline C^*_j)^2\right]^{1/2}.
\label{eq:bootstrap-scale}
\end{equation}
The standardized deviations are centered on the observed $\widehat C_j$,
not on $\overline C^*_j$. For a component with $s_j\leq10^{-12}$,
all its $Z_{bj}$ values are set to zero. If every component satisfies
that threshold, the critical value is zero. Otherwise, each draw's
maximum includes these zero components, and its empirical 95th percentile
uses linear interpolation. All 5,000 draws are retained. The final
subtraction $L_j=\widehat C_j-c_{.95}s_j$ still uses the actual $s_j$,
including for thresholded components.

For the ownership matrix, we resample 400 patients with replacement using seed 20260904. Each draw uses the same patient indices for every question and intervention, preserving all paired baseline differences. The clinical family comprises the 30 fixed contrasts $W_{q,q}-W_{q,d}$ for the six questions and their five alternatives. Each question's simultaneous ownership lower bound is the minimum of its five contrast bounds. Simultaneous upper bounds use the same critical value with the sign of every contrast reversed. A question has a stronger competitor if any upper bound is below zero. A question is unresolved if it has neither a positive minimum lower bound nor a negative upper bound. The shared-alias calculation uses a separate 12-statistic family: for each of six directions, its mean off-diagonal effect and its mean excess over the five corresponding matched-direction effects. The Mass confirmation applies the same construction to 200 independently selected patients, with seed 20260910 and a separate family of 20 Mass-minus-alternative contrasts across four coded prompt conditions.

Direct target-minus-maximum intervals summarise a different bootstrap quantity. For each draw, we recompute the largest competing mean and subtract it from the target mean. The displayed two-sided percentile intervals use the 2.5th and 97.5th percentiles of these margins. The earlier Effusion specificity supplement uses their 5th percentile as its one-sided lower bound. These intervals are distinct from the simultaneous bounds on the fixed linear contrasts in Equation~\ref{eq:max-t}.

The campaign applies the same construction to every block, using 600 test rows and 2{,}000 shared bootstrap draws of rows (seed 2026090601). The 30 contrasts $W_{q,q}-W_{q,d}$ receive simultaneous 95\% bounds from the centred studentised max-$T$ bootstrap. $O_q$ receives a direct percentile interval with the maximum recomputed in every draw. The logit-scale check in Table~\ref{tab:cf-scale} applies the same rules to margin-scale deltas.
\section{Reproducibility specification}
\label{app:repro}

This appendix specifies each family's consumed block, the consumed-token mask, the projection and probe, the control constructions, the cohort roles, and the protocol record at the registered protocol's level of detail. The released code implements each component.

\subsection{Consumed loci per family}
\label{app:repro-loci}

The primary locus is the output of the last visual block that feeds the consumer. The connector locus is the consumer's output that enters the language model. We identify the block from each checkpoint's own configuration when loading the checkpoint. Each block's preflight record stores the identified block, its hidden dimension, its consumer, and any branch that bypasses it.

\begin{itemize}[leftmargin=*,itemsep=1pt,topsep=2pt]
\item Qwen2.5-VL, Lingshu, and Qwen3-VL: the merger reads the last block of the vision transformer, which has 32 blocks in Qwen2.5-VL and Lingshu, 24 in Qwen3-VL 4B, and 27 in Qwen3-VL 8B and 32B. A Qwen2.5-VL or Lingshu image contributes 576 patch tokens, which the merger joins in $2\times2$ groups into 144; a Qwen3-VL image contributes at most 1{,}024. The Qwen3-VL DeepStack mergers, which read earlier blocks, are not written.
\item Gemma~3 and MedGemma: the projector reads the last of the 27 SigLIP encoder layers through the tower's final layer normalisation, and consumes all 4{,}096 patch tokens of the $896\times896$ image.
\item InternVL3.5: the projector reads the last InternViT encoder layer, of 24 at 8B and 14B and of 45 at 38B; the tower's final normalisation is an identity in the released configuration. An image contributes 1{,}024 patch tokens from a single $448\times448$ tile.
\item LLaVA-1.5 and LLaVA-Med: the projector reads the penultimate of the 24 CLIP encoder layers, and the final layer is computed and discarded. An image contributes 576 patch tokens.
\item Llama~3.2 Vision: the projector input concatenates the last of the eight global-encoder layers with five intermediate layers of the local encoder, which bypass the primary locus. An image contributes 1{,}601 tokens, the CLS token and 1{,}600 patches, from its one real $560\times560$ tile.
\end{itemize}
Table~\ref{tab:cf-models} lists the consumed-token count of every checkpoint.

\subsection{The consumed-token mask}
\label{app:repro-mask}

The write in Equation~\ref{eq:intervention} changes only the consumed tokens at the positions the consumer reads at the locus. We derive these positions for each batch from that batch's processor outputs:

\begin{itemize}[leftmargin=*,itemsep=1pt,topsep=2pt]
\item Qwen2.5-VL, Lingshu, and Qwen3-VL: the tokens of all images in a batch form one sequence in which every image occupies one contiguous span, whose bounds follow from the image's patch grid. At the connector locus the span is a quarter as long, one merged token per $2\times2$ patch group.
\item Gemma~3 and MedGemma: all 4{,}096 patch tokens; the tower produces no CLS token. The number of written tokens must equal the number of image placeholder tokens in the language-model input.
\item InternVL3.5: the patch tokens of the single $448\times448$ tile. The CLS token, which the consumer drops, is excluded. The number of written tokens must equal the number of image placeholder tokens.
\item LLaVA-1.5 and LLaVA-Med: the 576 patch tokens. The CLS token, which the checkpoint's feature selection drops, is excluded.
\item Llama~3.2 Vision: the 1{,}601 tokens, CLS and 1{,}600 patches, of the one real $560\times560$ tile. The seven padding tokens of each tile and the three padding tiles are excluded, in agreement with the model's own aspect-ratio and cross-attention masks.
\end{itemize}

Scoring stops with an error and no fallback when an image span does not exactly cover the locus, a mask length differs from the token count at the locus, an image has no consumed token, or one forward pass reaches the locus more than once.

Preflight check (C) verifies the mask. It captures the locus tensor on a clean pass and a pass steered at $\alpha=+0.25$ along the first concept direction. It reports the maximum absolute difference outside the mask and the maximum absolute changes in the consumer's output and answer logits. Across all \cfBlocks{} blocks, the change outside the consumed tokens is $0$ at both loci. The consumer's output and answer logits both change. The same records show a bitwise-identical repeated clean pass (check A) and a bitwise-identical $\alpha=0$ pass (check B) for every block.

\subsection{Projection, lift, scaler, and probe}
\label{app:repro-probe}

We construct the projection $R\in\R^{D\times512}$ by drawing standard normal variates once from a PCG64 generator seeded with $0$, dividing by $\sqrt{512}$, and storing the result in float32. Its entries are independent $\mathcal{N}(0,1/512)$ variates. The seed is $0$ for every checkpoint, dataset, and locus. $D$ is the family's hidden dimension. We neither normalise nor orthogonalise the columns of $R$. We normalise the lifted direction of Equation~\ref{eq:lift} once to unit $\ell_2$ norm in model space. The Gram matrix $R^{\top}R/D$ has full rank 512 in all \cfAltdirdBlocks{} blocks that carry the displacement module, with eigenvalues $\cfAltdirdEigMin$--$\cfAltdirdEigMax$ and condition number \cfAltdirdCondMin{}--\cfAltdirdCondMax{}.

Projecting the token-mean read vector $\hvec_i$ gives features $R^{\top}\hvec_i$. A featurewise standardiser fitted only on training rows supplies the mean $\boldsymbol\mu$ and scale $\boldsymbol s$. Every row is standardised as $(R^{\top}\hvec-\boldsymbol\mu)\oslash\max(\boldsymbol s,10^{-8})$. Each concept probe is an $\ell_2$-regularised logistic regression with inverse regularisation strength $C=1$. We fit each probe by L-BFGS for at most 2{,}000 iterations with unweighted classes and seed $0$, using training rows with a known label for that concept. The positive-class coefficient fixes the direction's clinical pole. For each of the two refit seeds, we resample training units with replacement, retaining all rows of each sampled unit at its multiplicity. We then refit the scaler and six probes using the same projection. The random family, shams, and control probes remain at seed 0.

\subsection{Control constructions}
\label{app:repro-controls}

\paragraph{Random directions.} We draw one $119\times D$ standard-normal matrix using a PCG64 generator seeded with $0$ and normalise each row to unit $\ell_2$ norm. Every question in the block uses the same 119 directions.

\paragraph{Coordinate-permutation sham.} Immediately after the random draw, the same generator produces one permutation of $\{0,\dots,D{-}1\}$ per concept, in protocol concept order. We form concept $c$'s sham by rearranging the coordinates of $\hat\wvec_c$ with concept $c$'s own permutation. This preserves the direction's coordinate multiset and unit norm. We write each question with its own sham.

\paragraph{Matched random-label probes.} Each row has a type: view $\times$ sex $\times$ age decade on the chest sets, and aspect-ratio $\times$ area buckets on COCO. For control seed $k\in\{0,\dots,19\}$, we traverse the types in sorted order and use a PCG64 generator seeded with $k$ to draw one label per type uniformly from $\{0,1\}$. If all types receive the same label, we flip the first. Every row of a type receives that type's label. We give each control probe the concept probe's capacity and nuisance structure by fitting it on the same training rows and known-label subset with the same hyperparameters. A block fits control probes when its training rows span at least eight types. Selectivity is Equation~\ref{eq:selectivity}. Its bootstrap over calibration units retains a draw when the real AUROC is finite and at least ten of the twenty control AUROCs are finite. We grade a cell for readability when at least 1{,}900 of the 2{,}000 draws are retained.

\subsection{Cohort roles}
\label{app:repro-cohorts}

Table~\ref{tab:cf-roles} lists each role's size, its independent units, and its use. One role holds more than one row per unit: the NIH training role, which takes every study of each training patient. Every other role of every dataset holds one row per unit, and no unit appears in two roles of a dataset. We use the calibration role only to determine readability and answer capability. We neither fit probes nor score writes on it. Every write uses the test role.

\begin{table}[h]
\centering
\scriptsize
\setlength{\tabcolsep}{3.5pt}
\caption{\textbf{Cohort roles.} Rows per role with independent units in parentheses and the role's use; units are patients on the chest sets and images on COCO. The last two rows give the largest rows per unit and the largest overlap between two roles.}
\label{tab:cf-roles}
\begin{tabular}{>{\raggedright\arraybackslash}p{0.115\textwidth}lll>{\raggedright\arraybackslash}p{0.335\textwidth}}
\toprule
Role & NIH ChestX-ray14 & CheXpert Plus & COCO & Use \\
\midrule
train & \cfCohortNihTrainRows{} (\cfCohortNihTrainUnits{}) & \cfCohortChexTrainRows{} (\cfCohortChexTrainUnits{}) & \cfCohortCocoTrainRows{} (\cfCohortCocoTrainUnits{}) & Scaler, six concept probes, twenty control probes, the two nuisance probes of the chest sets, and the answer-direction ridge. No row of this role is scored. \\
preflight & \cfCohortNihPreflightRows{} (\cfCohortNihPreflightUnits{}) & \cfCohortChexPreflightRows{} (\cfCohortChexPreflightUnits{}) & \cfCohortCocoPreflightRows{} (\cfCohortCocoPreflightUnits{}) & The seven acceptance gates, at both loci. \\
calibration & \cfCohortNihCalibrationRows{} (\cfCohortNihCalibrationUnits{}) & \cfCohortChexCalibrationRows{} (\cfCohortChexCalibrationUnits{}) & \cfCohortCocoCalibrationRows{} (\cfCohortCocoCalibrationUnits{}) & Clean pass at $\alpha=0$: probe readability and answer capability. \\
test & \cfCohortNihTestRows{} (\cfCohortNihTestUnits{}) & \cfCohortChexTestRows{} (\cfCohortChexTestUnits{}) & \cfCohortCocoTestRows{} (\cfCohortCocoTestUnits{}) & Every write matrix, dose, refit, template, locus, and addendum module. \\
valid & -- & \cfCohortChexValidRows{} (\cfCohortChexValidUnits{}) & -- & The regrade on radiologist consensus labels (Table~\ref{tab:cf-valid}). \\
\midrule
rows per unit & \cfCohortNihTrainMax{} (\cfCohortNihTrainMulti{} units) & \cfCohortChexTrainMax{} (\cfCohortChexTrainMulti{}) & \cfCohortCocoTrainMax{} (\cfCohortCocoTrainMulti{}) & Largest number of rows one unit contributes, with the units contributing more than one; the NIH training role is the campaign's only role above one. \\
shared units & \cfCohortNihMaxOverlap{} & \cfCohortChexMaxOverlap{} & \cfCohortCocoMaxOverlap{} & Largest number of units two roles of the dataset share, over all \cfCohortPairsAll{} role pairs of the three datasets. \\
\bottomrule
\end{tabular}
\end{table}
\FloatBarrier

\subsection{Protocol record and amendments}
\label{app:repro-protocol}

Each block writes a run record containing the protocol identifier \texttt{cf-transfer-v1}, a run identifier naming the checkpoint and dataset, the evaluation code version that produced the block, the pinned checkpoint and processor revisions, the library versions, the numeric precisions, and the resolved processing settings, including the rendered chat template. All \cfBlocks{} blocks carry the protocol identifier \texttt{cf-transfer-v1}. The protocol file registers the code version it was frozen against.

Table~\ref{tab:cf-amendments} lists every addition to the protocol file after the first wave, what it fixes, and the date the file records for it. Every analysis the paper reports that is not one of the planned modules appears here, with the date the file gives it. The \cfCoveragePlannedWord{} planned modules have no amendment date. Each amendment entry names its own rows, references, and decision rule before it names a result, and the robustness tables in Appendix~\ref{app:cf-robustness} identify the blocks that carry each one. A1 and A2 also carry the identifiers of the execution log.

\begin{table}[h]
\centering
\scriptsize
\setlength{\tabcolsep}{5pt}
\caption{\textbf{Protocol amendments.} Every addition to the registered protocol file after the first wave, with the date the file records for it. Each date precedes the first outcome shard of the module it registers.}
\label{tab:cf-amendments}
\begin{tabular}{ll>{\raggedright\arraybackslash}p{0.62\textwidth}}
\toprule
Amendment & Date & What it adds and what it fixes \\
\midrule
A1 & 2026-09-14 & Per-block primary template: the single-template modules score the first eligible template in the order (IY, IB), so a checkpoint whose yes/no template fails the image-free mapping gate is graded on its B-positive template instead of losing its grade. Llama~3.2 Vision 11B re-ran on all three datasets. \\
A2 & 2026-09-14 & CheXpert gains the dose sweep, the two probe refits, the connector locus, the connector-locus calibration, and the five additional templates, with Effusion and Edema as its two template concepts, so the second chest set carries the same modules as the first. \\
A3 & 2026-09-14 & Four alternative direction constructions per concept: difference of class means, Haufe pattern, orthogonalised normal, and residualised normal. Fixes the reading that the deficit belongs to the logistic normal. \\
A4 & 2026-09-14 & Extra clinical directions for every additional dataset label with at least 100 known positives and negatives, added to the competitor family. Fixes the reading that the six protocol concepts are too narrow a competitor set. \\
A5 & 2026-09-14 & Softmax-weighted and top-quarter token weights for the six logistic directions. Fixes the reading that a uniform write dilutes the concept over tokens that do not carry it. \\
A6 & 2026-09-14 & The write grid rescored with fp32 weights and forward pass, and in bf16 at batch size one. Fixes the reading that the grades are a numerical artefact. \\
A7 & 2026-09-14 & The answer direction: a ridge regression of the clean answer margin on the same projected, train-scaled features. Fixes the reading that nothing in the consumed representation moves the answer. \\
A8 & 2026-09-15 & The difference-of-means and pattern vectors lifted as displacements through the minimum-norm preimage. Fixes the reading that the coefficient lift is what fails. \\
A9 & 2026-09-15 & Three non-clinical attribute directions written with the six clinical directions as one nine-direction family. Fixes the reading that the deficit follows from something about radiographs rather than from the concepts. \\
A10 & 2026-09-15 & The answer directions fitted on the primary template written under the five held-out templates. Fixes the reading that the answer direction is a property of one prompt. \\
A11 & 2026-09-15 & The write grid regraded on the CheXpert validation split with radiologist consensus labels. Fixes the reading that report-derived labels are the source of the deficit. \\
A12 & 2026-09-16 & The crossed tower and reader: one checkpoint's vision tower loaded behind the other's connector and language model, and all four combinations of tower and reader graded on the same rows. Fixes the reading that ownership is a property of the representation alone or of the reader alone. \\
A13 & 2026-09-16 & The identical-tensor replay: the consumed block's stored token tensors fed to every reader of a shared-tower group. Fixes the reading that two readers of one tower differ because their inputs differ by bf16 rounding. \\
A14 & 2026-09-16 & The four semantic endpoints: the negated question, the counterbalanced two-finding forced choice, and the finding word's log-probability in a report continuation, each with its own clean baseline, sham, and random family, together with the incremental-validity test of the ownership contrast. Fixes the reading that ownership is validated only against the model's answers under the six templates. \\
A15 & 2026-09-16 & The fine-grained object family: six further COCO categories chosen by the rule stated below, fitted, written, and graded as their own six-direction family in the same blocks and rows. Fixes the reading that the chest-to-COCO contrast is a contrast in task difficulty. \\
A16 & 2026-09-16 & The clean pass of the fine-grained categories on the calibration rows, so their clean-answer AUROC is measured on the same rows and by the same rule as every other cell and the difficulty matching compares like with like. \\
A17 & 2026-09-16 & The matched attribute reference: the protocol's own 119-direction random family written against the three attribute questions. Fixes the attribute-versus-finding comparison, which until then referenced attribute cells against their sham alone and clinical cells against the random 95th percentile as well. \\
A18 & 2026-09-16 & Three phrasings of each attribute question scored clean on the calibration rows, with the selection rule and its reason recorded per cell. Fixes the reading that an attribute cell is unowned because the model cannot answer the phrasing that was written. \\
A19 & 2026-09-16 & The expert-label refits: the six concept probes refitted on the radiologist-labelled films, cross-fitted over patients, and written on the test rows. Fixes the reading that the label source, not the concept, decides ownership. \\
A20 & 2026-09-16 & Projection seeds 1 and 2: the probes, the random family, and the sham redrawn in two further random projections and regraded against that projection's own references. Fixes the reading that the grades depend on the seed-0 projection. \\
\bottomrule
\end{tabular}
\end{table}

\FloatBarrier

\paragraph{The fine-grained category selection rule.} The rule that picks the six fine-grained COCO categories of amendment A15 is stated in the protocol file in full, and it reads on the labels alone. A category is eligible when its training positives are at least as many as the weakest of the six existing concepts, when it has at least ten positives among the 400 calibration rows and at least ten among the 600 test rows, and when its mean relative instance area is below the median of the six existing concepts. Eligible categories are ranked by calibration positives in descending order, ties by mean relative instance area in ascending order, and remaining ties by category identifier; the first six are taken. The rule uses the label counts and the instance areas of the fixed cohorts and nothing else: no probe, no answer, and no write enters it. The protocol file carries the rule under the date of amendment A15, and the module's first outcome shard is later than that date, so the six categories and the two matching windows were fixed before any of them was scored.
\end{document}
